\documentclass[a4paper,11pt]{article}
\usepackage[a4paper,margin=2.4cm]{geometry}
\usepackage{amsmath,amssymb,amsthm}
\usepackage{fontspec}
\usepackage{unicode-math}
\usepackage{graphicx}
\usepackage{booktabs}
\usepackage{hyperref}
\usepackage{xcolor}
\definecolor{neonmagenta}{HTML}{FF10F0}
\usepackage{enumitem}
\usepackage{seqsplit}
\usepackage{fvextra}
\usepackage{etoolbox}
% Allow long filenames to break across lines
\newcommand{\filename}[1]{\texttt{\seqsplit{#1}}}
% fvextra Verbatim for Lean listings (Unicode-safe, keeps characters in order)
\newcounter{leanlisting}
\newenvironment{leanlisting}[2]{%
  \refstepcounter{leanlisting}\label{#2}%
  \begin{center}\textit{\small Listing \theleanlisting: #1}\end{center}%
  \Verbatim[breaklines=true,breakanywhere=true,fontsize=\footnotesize,frame=single,framesep=4pt,rulecolor=\color{gray!50},fillcolor=\color{gray!5}]%
}{%
  \endVerbatim%
}
% fvextra Verbatim for code listings (Unicode-safe)
\newcounter{codelisting}
\newenvironment{codelisting}[2]{%
  \refstepcounter{codelisting}\label{#2}%
  \begin{center}\textit{\small Listing \thecodelisting: #1}\end{center}%
  \Verbatim[breaklines=true,breakanywhere=true,fontsize=\footnotesize,frame=single,framesep=4pt,rulecolor=\color{gray!50},fillcolor=\color{gray!5}]%
}{%
  \endVerbatim%
}
% Fonts: Fira Sans (text), Fira Mono (code), Fira Math (math) are primary.
% Where a glyph is missing from the Fira fonts, it is substituted from
% Noto Sans (text) or Noto Sans Math (math and listing symbols):
%  - U+2208 (in) is missing from Fira Mono; in text/listings it is taken
%    from Noto Sans Math. Fira Math has it, so math-mode \in is Fira.
%  - U+2090 (subscript a) is missing from Fira Mono; in text/listings it
%    is taken from Noto Sans. In math it stays a roman subscript.
%  - Math glyphs missing from Fira Math are substituted by range from
%    Noto Sans Math (range assignments added below the \setmathfont line
%    for exactly the codepoints the Fira Math coverage lacks).
\setmainfont{Fira Sans}[BoldFont={Fira Sans Bold},ItalicFont={Fira Sans Italic},BoldItalicFont={Fira Sans Bold Italic}]
\setmonofont{Fira Mono}[BoldFont={Fira Mono Bold}]
\setmathfont{Fira Math}[Scale=1.15]
% NOTO_RANGE_SUBS
\newfontfamily\notosans{Noto Sans}
\newfontfamily\notosansmath{Noto Sans Math}
\catcode`\∈=\active
\protected\def∈{\ifmmode\in\else{\notosansmath\char"2208}\fi}
\catcode`\ₐ=\active
\protected\defₐ{\ifmmode{}_{\mathrm{a}}\else{\notosans\char"2090}\fi}
\hypersetup{colorlinks=true,linkcolor=blue!60!black,citecolor=blue!60!black,urlcolor=blue!60!black}
\setlength{\parskip}{0.6em}
\setlength{\parindent}{0pt}
\setlength{\emergencystretch}{2em}

% Section, equation, figure, table and listing numbering restarts in each
% part, matching the standalone papers; hyperref anchors stay unique.
\makeatletter
\@addtoreset{section}{part}
\@addtoreset{equation}{part}
\@addtoreset{figure}{part}
\@addtoreset{table}{part}
\@addtoreset{leanlisting}{part}
\@addtoreset{codelisting}{part}
\makeatother
\renewcommand{\thesection}{\arabic{section}}
\renewcommand{\theHsection}{\thepart.\arabic{section}}
\renewcommand{\theHequation}{\thepart.\arabic{equation}}
\renewcommand{\theHfigure}{\thepart.\arabic{figure}}
\renewcommand{\theHtable}{\thepart.\arabic{table}}
\renewcommand{\theHleanlisting}{\thepart.\arabic{leanlisting}}
\renewcommand{\theHcodelisting}{\thepart.\arabic{codelisting}}

\theoremstyle{plain}
\newtheorem{theorem}{Theorem}[section]
\newtheorem{proposition}[theorem]{Proposition}
\newtheorem{lemma}[theorem]{Lemma}
\newtheorem{corollary}[theorem]{Corollary}

\theoremstyle{definition}
\newtheorem{definition}[theorem]{Definition}
\newtheorem{example}[theorem]{Example}

\theoremstyle{remark}
\newtheorem{remark}[theorem]{Remark}
\renewcommand{\theHtheorem}{\thepart.\thesection.\arabic{theorem}}

\begin{document}

\title{\textbf{\Large Elimination of the GGHV Normal Forms at Degree~Pair $(72,108)$, and Closure of the Two-Dicritical Line:}\\[0.6em]
\large A Certified Computational Audit in Three Parts\\[0.4em]
\normalsize Branch-(a,b) Elimination \textperiodcentered{} Branch-(c) Elimination \textperiodcentered{} Dicriticals at Infinity of Keller Maps}
\author{\textbf{Brandon D. McCrary}\\[0.25em]
\textit{Independent Researcher}\\[0.25em]
\href{https://orcid.org/0009-0008-2804-7165}{ORCID: 0009-0008-2804-7165}\\[0.6em]
{\small \parbox{0.92\textwidth}{\centering \textbf{Computation, Lean and LaTeX Source:}\\
{\color{neonmagenta}\hypersetup{urlcolor=neonmagenta}\href{https://github.com/brandonmccraryresearch-cloud/2D-Jacobian-Conjecture-Workspace}{\nolinkurl{https://github.com/brandonmccraryresearch-cloud/2D-Jacobian-Conjecture-Workspace}}}}}}
\date{September 28 -- October 1, 2026; combined edition, October 22, 2026; finalized November 4, 2026 (Part~II validation update)}
\maketitle

\begin{abstract}
\textbf{Part~I.}\enspace We present a computational elimination of the branch~(a,b) normal form arising from Proposition~4.3 of Guccione--Guccione--Horruitiner--Valqui~\cite{GGHV} for a hypothetical $(72,108)$ counterexample to the planar Jacobian conjecture. Every base object is source-gated directly against GGHV v1 before any proof is drafted.

Starting from the verified Newton polygons
\[
\begin{aligned}
N(P)&=\operatorname{conv}\{(0,0),(1,0),(8,14),(8,16)\},\\
N(Q)&=\operatorname{conv}\{(0,0),(2,1),(12,21),(12,24)\},
\end{aligned}
\]
with $[P,Q]=x^2$ post-$\varphi$, we reduce the $72$-unknown bilinear system to five one-variable layers $E_1,\dots,E_5$ totaling $92$ equations, verified as an exact symbolic polynomial identity.

The top layer $E_5$ is the Davenport--Stothers extremal equation $\alpha\beta+u(2\alpha\beta'-3\alpha'\beta)=1$~\cite{Stothers,Davenport}. The Belyi classification gives the geometric upper bounds $3$, $10$, $35$ chart points for $m=3,5,7$; for $m=7$ the bound is sharp, with the degree-$35$ eliminating polynomial $\mathcal{W}(a_6)$ irreducible over $\mathbb{Q}$, a quintic in $S=a_6^7$ with nonzero discriminant and exactly one real root. The $35$ solutions are distinct and form a single Galois orbit: one real solution and $17$ conjugate pairs, grouped in five $S$-orbits. The classification is proved twice: the Belyi bound combined with the $35$ exhibited points gives a complete classification (relying on Riemann's existence theorem and the computed msolve parametrization), while \S\ref{ab:sec:chartproof} gives an independent algebraic proof, valid over every field of characteristic~$0$, which is the route formalized in Lean. In reversed coordinates the chart equations are the coefficient recursion of $(1+y_1v+\dots+y_7v^7)^{3/2}$. Six weighted-homogeneous relations, of weights $4,5,5,6,6,7$ and with explicit certificates, cut out the five torus orbits. We formalize this proof in Lean.

For $m=3$ and $m=5$ the bounds $3$ and $10$ are also sharp. These two classifications are kernel-checked in Lean as well, over every field of characteristic~$0$. In the normalized chart, the system holds if and only if two things hold: the eliminant vanishes ($3a_2^3-32$ for $m=3$, $9a_4^{10}+37200a_4^5+95051008$ for $m=5$), and the remaining unknowns take explicit values. In each case the eliminant has no repeated root. So over every algebraically closed field of characteristic~$0$ there are exactly $3$ solutions for $m=3$ and $10$ for $m=5$, and this count is kernel-checked too. Both eliminants are irreducible over~$\mathbb{Q}$, and this is kernel-checked as well. Theorem~\ref{ab:thm:main} does not need these two cases.

The descent proceeds through $E_4$ (an $18\times19$ matrix of rank $17$, kernel dimension $2$), $E_3$ (a $19\times20$ matrix of rank $18$, kernel dimension $2$, unconditionally solvable for every direction $t$), and $E_2$ (injective $19\times 12$, rank $12$), leaving seven explicit compatibility conditions. We expand these as $b_i(t)+s_1M_i(t)+s_2L_i(t)=0$ with $b_i$ cubic and $M_i,L_i$ linear, eliminate $(s_1,s_2)$ to obtain $35$ quintic $3\times3$ minors. The entire descent is computed exactly over the quintic field $K_5$ in characteristic $0$: the $35$ minors span all binary quintics, forcing $t=0$. The case $t=0$ forces $B_0$ constant, contradicting the required vertex $(12,24)$. The Lean formalization quantifies over all roots of $R$ and all torus parameters, covering all $35$ top-layer solutions.

Hence no polynomials $P,Q$ satisfying $\mathrm{NewtonNF2}$ have $[P,Q]=\lambda x^2$ with $\lambda\neq0$ (Theorem~\ref{ab:thm:main}). In Lean this NewtonNF2 form is a single kernel-checked theorem, \texttt{main\_theorem}, on the standard axioms only, independent of GGHV Proposition~4.3. Its top-layer input is \texttt{chartClassification\_holds}, a formalized algebraic proof of the $m=7$ top-layer classification of Proposition~\ref{ab:prop:top-class}, in the normalized chart (\S\ref{ab:sec:chartproof}). The further rigidity $a_{8,16}=0$ (Corollary~\ref{ab:cor:a816}) is kernel-checked in Lean as well, together with the vanishing of every other coefficient of $P$ and $Q$ off the top edges except the constant terms (\texttt{lower\_edge\_rigidity}); it uses the layer $E_1$. A second proof of $a_{8,16}=0$, outside Lean and independent of the descent, is an explicit certificate over $K_5$: the identity $a_{8,16}^2=\sum_{k=1}^{75}H_ke_k$, where the $e_k$ are the $75$ equations of the layers $E_4,\dots,E_1$. Two independent programs check it exactly. The application to the Jacobian conjecture is conditional on the correctness and exhaustiveness of~\cite[Proposition~4.3]{GGHV} (Corollary~\ref{ab:cor:gghv}).

\medskip

\textbf{Part~II.}\enspace We present a computational elimination of the branch-(c) normal form arising from case~(1) of Proposition~4.3 of Guccione--Guccione--Horruitiner--Valqui~\cite{GGHV} for a hypothetical $(72,108)$ counterexample to the planar Jacobian conjecture. Every base object is source-gated directly against GGHV v1 before any proof is drafted. This paper is the branch-(c) companion to the branch-(a,b) elimination~\cite{BranchABv19}, whose format it follows.

Starting from the verified Newton polygons
\[
\begin{aligned}
N(P)&=\operatorname{conv}\{(0,0),(1,0),(8,14),(8,16),(0,8)\},\\
N(Q)&=\operatorname{conv}\{(0,0),(2,1),(12,21),(12,24),(0,12)\},
\end{aligned}
\]
with $[P,Q]=\lambda x^2$ post-$\varphi$ and $\lambda\ne0$, we reduce the bracket condition to layer identities $E_n$ in the grading $w(x)=2$, $w(y)=-1$, extract $132$ exact scalar equations, and eliminate the layers $E_4,\dots,E_{-2}$ to twelve chart conditions $\Omega,\Psi,\Phi_1,\Phi_2,\Theta_1,\Theta_2,\Theta_3$ plus five pure rows, over the quintic field $K_5=\mathbb{Q}[w]/(w^5-w^4+3w^3+3w^2+26)$. The top layer is the branch-(a,b) problem, whose exact $K_5$ classification is inherited. The descent is proved twice: symbolically, as $600$ kernel-checked reflective polynomial identities in Lean~4 (no \texttt{sorry}, standard axioms only), and numerically, by exact finite-field linear algebra.

The chart is empty. On the stratum $b_{11,20}=0$ this is proved in the Lean kernel from an exact $K_5$ certificate (degree~$2$, $75$ terms). On the whole chart it follows from a rank lemma: full row rank of the fixed-degree Macaulay matrix lifts from $\mathbb{F}_p$ to $K_5$ through a nonzero minor (retained as independent corroboration; the Lean certificate of Theorem~\ref{bc:thm:unconditional} is now the load-bearing route). At weight $W=24$ the $3199\times 6054$ Macaulay matrix has full row rank $3199$ at two primes ($p=1000003$, $p=32003$), verified by two independent implementations (FLINT and NumPy), with negative controls. An earlier reduction lemma used in a preliminary version of this work is false as stated and is retracted here; the rank lemma is its correct replacement, and the distinction---rank lifts, mere consistency does not---is proved by counterexample.

In Lean, with all $790$ modules built and \texttt{\#print axioms} reporting only \texttt{propext}, \texttt{Classical.choice}, \texttt{Quot.sound},
\[
\begin{aligned}
\texttt{ChartEmptyC\_T1ne0}\ L\ &\Rightarrow\
\lnot\exists\,P\,Q\,\lambda,\ \lambda\ne0\land\texttt{NewtonNFc}\ P\,Q \\
&\qquad\land[P,Q]=\lambda x^2,
\end{aligned}
\]
and, since November 4, 2026, \texttt{ChartEmptyC\_T1ne0} is proved inside Lean with no hypotheses (\texttt{challenge\_target}; Theorem~\ref{bc:thm:unconditional} of Part~II), so branch-(c) elimination is unconditional in Lean. The rank computation remains as independent corroboration. The application to the Jacobian conjecture is conditional on the correctness and exhaustiveness of~\cite[Proposition~4.3]{GGHV} (Corollary~\ref{bc:cor:gghv}).

\medskip

\textbf{Part~III.}\enspace Let $F = (P,Q) \colon \mathbb{A}^2 \to \mathbb{A}^2$ be a Keller map, that is,
$[P,Q] := P_x Q_y - P_y Q_x \in K^\times$ for $K$ algebraically closed of characteristic $0$,
and let $\lambda P + \mu Q$ be its pencil. The October 1 draft of this paper claimed that no
Keller map admits two distinct dicritical divisors, unconditionally. That claim is false: every
Keller map of degree at least $2$ has at least two exceptional dicriticals in total
(Corollary~\ref{td:cor:count-fails}), and the draft's $E_1$ degree obstruction
theorem is false as stated, with an explicit Keller
counterexample in \S\ref{td:sec:e1}.

The corrected target is \textbf{Theorem~R} (\S\ref{td:sec:statement}), which is \textbf{open}:
for every Keller map, the pencil $\lambda P + \mu Q$ has at most one dicritical divisor
\textbf{at infinity}. Theorem~R is equivalent to the planar Jacobian conjecture
(Proposition~\ref{td:prop:equiv}). This paper proves, or records as sketch or fixed-degree
computation, only the following steps toward it, each labelled with its grade:
\begin{itemize}[nosep]
\item the count of dicriticals (Proposition~\ref{td:prop:count}, grade B) and the position a
counterexample's dicriticals would occupy (\S\ref{td:sec:where});
\item the line at infinity is not dicritical for a non-affine Keller map
(Lemma~\ref{td:lem:linf}, grade B), and no first blow-up at infinity is dicritical
(Lemma~\ref{td:lem:e1}, algebra grade A, chart step grade B), so every dicritical at infinity
of a non-affine Keller map has depth at least $2$;
\item the $m = 1$ bridge equations collapse under two stated premises (P1) and (P2)
(Theorem~\ref{td:thm:bridge}, equations grade A, geometric reading grade B), where (P2) is
\textbf{not} derived and remains open;
\item the constant-leading-form argument for $m \geq 2$ on $E_k$, $k \geq 2$, as a
\textbf{proof sketch for free centers only} (Proposition~\ref{td:prop:constant-form}, grade C);
\item the adjacent $m = 1$ case as a computation at $d = 3$ only (grade C).
\end{itemize}
What remains open is collected in \S\ref{td:sec:open} (obligations O1--O6) and the tests any
proof must pass (T1--T4). \S\ref{td:sec:correspondence} maps each archived script to the exact
statement it certifies, under the corrected statements.
\end{abstract}

\section*{Validation note --- November 4, 2026 (Part~II update)}
This note records the deepest-research validation (fact-check mode: verification of a provided text against its record, not a systematic review) run on the November 4, 2026 draft before any repo incorporation. It is not part of the mathematical argument.
\begin{itemize}[nosep]
\item \textbf{What changed.} Part~II no longer rests on a grade-B premise. On November 4, 2026, \texttt{challenge\_target} was proved in Lean: \texttt{ChartEmptyC\_T1ne0} holds for every field of characteristic~$0$ with no hypotheses, with \texttt{\#print axioms} exactly \texttt{[propext, Classical.choice, Quot.sound]}; \texttt{branch\_c\_eliminated} follows from it by \texttt{main\_theorem\_c\_of\_chartEmpty\_T1ne0} with the same axiom triple (Theorem~\ref{bc:thm:unconditional}). The rank-lemma route of \S\ref{bc:sec:wholechart} is retained as independent corroboration at grade~B, not as the load-bearing step.
\item \textbf{Scope guard.} This eliminates one case --- case~(1) of GGHV Proposition~4.3, branch~(c) --- unconditionally in Lean. It does \emph{not} by itself exclude $(72,108)$ unconditionally and does \emph{not} prove the planar Jacobian conjecture: those still require the correctness and exhaustiveness of GGHV Proposition~4.3 and the reduction chain, which were ledgered, not re-derived, in this validation.
\item \textbf{Case numbering.} Verified throughout: Proposition~4.3 case~(1) is the polygon set \emph{with} vertices $(0,8)$ resp.\ $(0,12)$ $=$ \texttt{NewtonNFc} $=$ branch~(c); case~(2) is the set \emph{without} them $=$ \texttt{NewtonNF2} $=$ branches~(a),(b). No occurrence of the older shorthand ``(a),(b) $=$ case~(1)'' appears in this paper. Two Part~II line citations were corrected in this finalization (GGHV v1 lines~495/599--600 $\to$ 494/617--618) after re-extraction of the v1 text, which shows case~(1) at line~494 and the post-$\varphi$ case~c list at lines~617--618.
\item \textbf{Provenance caveat.} \texttt{challenge\_target} and \texttt{branch\_c\_eliminated} were proved November 4, 2026 as scratch files (\texttt{bridge\_attempt\_2026\_11\_04.lean}, \texttt{composition\_2026-11-04.lean}, run on lean-build-2 with logs \texttt{bridge\_challenge\_target\_2026-11-04.log} and \texttt{branch\_c\_eliminated\_2026-11-04.log}); they are not yet tracked repo modules. Incorporation as tracked modules is a separate, later step and is listed as open item~(5) in \S\ref{bc:sec:epistemic}.
\item \textbf{Title reading.} ``Closure of the Two-Dicritical Line'' in the title means closure (withdrawal and correction) of the October 1 draft's false two-dicritical claim, not a proof of Theorem~R: Theorem~R remains \textbf{open} (Part~III, \S\ref{td:sec:statement}).
\end{itemize}



\vspace{1.2cm}
\hrule
\vspace{1.4cm}

\tableofcontents
\newpage

\part{\texorpdfstring{Elimination of the GGHV Branch-(a,b) Normal Form at Degree Pair (72,108)}{Elimination of the GGHV Branch-(a,b) Normal Form at Degree Pair (72,108)}}

\section{Introduction}

\subsection{Context and prior work}

The planar Jacobian conjecture asserts that a polynomial map $(P,Q):\mathbb{C}^2\to\mathbb{C}^2$ with $[P,Q]:=P_xQ_y-P_yQ_x\in\mathbb{C}^\times$ is invertible. In~\cite{GGHV}, Guccione, Guccione, Horruitiner, and Valqui raise the lower bound on $\max\{\deg P,\deg Q\}$ for a hypothetical counterexample from $100$ (Moh~\cite{Moh}) to $108$, leaving open only the degree pair $(72,108)$ and its symmetric $(108,72)$.

Santib\'a\~nez-Leal has published machine evidence related to the $(72,108)$ candidate~\cite{Santibanez,SantibanezPaper}, explicitly not claiming elimination. The candidate itself is documented in~\cite{JC125note}.

For the degree pair $(72,108)$, \cite{GGHV} (an unrefereed preprint, arXiv:2204.14178v1) identifies two cases. The corner $(9,27)$ is discarded in \cite[\S5]{GGHV}; the open case is the corner $(8,28)$ with $(m,n)=(3,2)$. For this case, \cite[Proposition~4.3]{GGHV} shows that any counterexample gives rise, after an explicit (non-invertible) morphism $\varphi$, to one of two Newton polygon configurations. Case~(1), with additional vertices $(0,8)$ resp.\ $(0,12)$, is branch~(c). Case~(2), the subject of this paper, is branches~(a) and~(b):
\begin{equation}\label{ab:eq:polygons}
\begin{aligned}
N(P)&=\operatorname{conv}\{(0,0),(1,0),(8,14),(8,16)\},\\
N(Q)&=\operatorname{conv}\{(0,0),(2,1),(12,21),(12,24)\},
\end{aligned}
\end{equation}
with $[P,Q]=x^2$ in post-$\varphi$ coordinates. See Figure~\ref{ab:fig:polygons} for the polygon shapes.

\begin{figure}[htbp]
\centering
\includegraphics[width=0.95\linewidth]{figures_ab/fig1_newton_polygons.png}
\caption{The branch~(a,b) Newton polygons from~\cite[Proposition~4.3(2)]{GGHV}, source-gated at v1 lines~495 and 599--600. $N(P)$ has $25$ lattice points; $N(Q)$ has $47$.}
\label{ab:fig:polygons}
\end{figure}

\subsection{Main result}

\begin{theorem}[Elimination of the GGHV branch-(a,b) normal form]\label{ab:thm:main}
Let $K$ be a field of characteristic~$0$. Let $P,Q\in K[x,y]$ have
$\operatorname{supp}P\subset N_P$ and $\operatorname{supp}Q\subset N_Q$,
with $a_{1,0}\cdot a_{8,14}\cdot b_{2,1}\cdot b_{12,21}\neq0$ and
$[P,Q]=\lambda x^2$ for some $\lambda\neq0$. Then $b_{12,24}=0$.
In particular, no $P,Q$ satisfy $\mathrm{NewtonNF2}(P,Q)$
(listing~\ref{ab:lst:def-newtonnf2}; the polygons are in~\S\ref{ab:sec:source})
with $[P,Q]=\lambda x^2$.
\end{theorem}

\begin{corollary}[Lower-edge rigidity]\label{ab:cor:a816}
Under the hypotheses of Theorem~\ref{ab:thm:main}, $a_{8,16}=0$ as well. More precisely, every coefficient
$a_{i,j}$ with $j\geq 2i-1$ and every coefficient $b_{i,j}$ with $j\geq 2i-2$ vanishes, except $a_{0,0}$ and $b_{0,0}$:
$P=P_2+a_{0,0}$ and $Q=Q_3+b_{0,0}$.
In Lean this is \texttt{lower\_edge\_rigidity} (listing~\ref{ab:lst:a816}), kernel-checked on the standard axioms;
$\lambda\neq0$ is not needed. A second proof of $a_{8,16}=0$, outside Lean, is an explicit certificate over $K_5$
(\filename{scripts/a816\_certificate/}), checked exactly by {\sc Singular} and by python-flint;
see~\S\ref{ab:sec:lower-edge}. The corollary is not needed for the elimination: $b_{12,24}=0$
already contradicts the required Newton vertex $(12,24)$.
\end{corollary}

\begin{remark}[Lean-to-mathematics correspondence]\label{ab:rem:nf2-bridge}
The Lean predicate \texttt{NewtonNF2} (listing~\ref{ab:lst:def-newtonnf2}) is
\[
\texttt{NewtonNF2}(P,Q)\iff
\begin{cases}
\operatorname{supp}P\subseteq N_P,\ \operatorname{supp}Q\subseteq N_Q
\text{ (via \texttt{inNP}, \texttt{inNQ}),}\\
a_{0,0},a_{1,0},a_{8,14},a_{8,16}\neq0,\\
b_{0,0},b_{2,1},b_{12,21},b_{12,24}\neq0,
\end{cases}
\]
i.e.\ exactly the support conditions of Theorem~\ref{ab:thm:main} together with the nonvanishing of the eight vertex coefficients. The half-plane predicates are tied to the vertex sets by \texttt{decide} checks run on duplicate predicates \texttt{latticeNP}/\texttt{latticeNQ} (which restate the same inequalities): the lattice-point counts ($25$ and $47$) and the vertex-membership check confirm the eight vertices lie on the polygons. This bridge is what lets the kernel-checked \texttt{main\_theorem} speak about the same $P,Q$ as Theorem~\ref{ab:thm:main}.
\end{remark}

\begin{corollary}[Conditional on GGHV]\label{ab:cor:gghv}
Assuming \cite[Prop.~4.3]{GGHV}, no counterexample to the plane Jacobian conjecture arises in case~$(8,28)$, sub-cases~a) and~b).
\end{corollary}

The proof is computational but exact. We reduce the $72$-unknown system to triangular layers (Figure~\ref{ab:fig:layers}), determine the top layer via the Belyi-map/Riemann-existence argument (Figure~\ref{ab:fig:top}), and independently by an algebraic argument formalized in Lean (\S\ref{ab:sec:chartproof}). We then obstruct the descent via the Lean-formalized layer certificates (\S\ref{ab:sec:descent}). Key steps were verified in {\sc Singular}~4.4.1 and Python/{\sc SymPy}; scripts are archived with this draft.

\begin{figure}[htbp]
\centering
\includegraphics[width=0.75\linewidth]{figures_ab/fig2_layer_descent.png}
\caption{The triangular layer structure: $92$ equations in five layers. Each layer is affine-linear in its unknowns given the solution above it.}
\label{ab:fig:layers}
\end{figure}

\subsection{Logical status}

\begin{remark}\label{ab:rem:conditional}
Theorem~\ref{ab:thm:main} does not depend on GGHV. Its proof uses Proposition~6.1 for $m=7$,
which is proved twice. On paper it follows from Riemann's existence theorem
(\S\ref{ab:sec:top}). In \S\ref{ab:sec:chartproof} the classification it provides is proved
algebraically, in the normalized chart, and that proof is formalized in Lean as
\texttt{chartClassification\_holds} (listing~\ref{ab:lst:proofchartclass}).
In Lean, \texttt{main\_theorem\_of\_chart} (listing~\ref{ab:lst:chart}) derives the NewtonNF2 form of
Theorem~\ref{ab:thm:main} from the explicit statement \texttt{ChartClassification}
(listing~\ref{ab:lst:def-chartclass}). \texttt{main\_theorem} (listing~\ref{ab:lst:unconditional})
combines it with \texttt{chartClassification\_holds} into a single kernel-checked
theorem independent of GGHV Proposition~4.3, and \texttt{\#print axioms} reports only \texttt{propext}, \texttt{Classical.choice}
and \texttt{Quot.sound}.
The $b_{12,24}=0$ part follows from \texttt{no\_completion\_K5\_PQ} together with \texttt{topLayerClassification\_of\_chart} and \texttt{chartClassification\_holds}; the $a_{8,16}=0$ part (Corollary~\ref{ab:cor:a816}) is \texttt{a816\_eq\_zero}, and the whole corollary is \texttt{lower\_edge\_rigidity} (listing~\ref{ab:lst:a816}); both are kernel-checked on the same three axioms and reuse the descent up to $t=0$. Outside Lean, an explicit certificate $a_{8,16}^2=\sum_k H_k e_k$ over $K_5$, stored as the $76$ cofactors of $1=\sum_k z^2H_k e_k-(1+a_{8,16}z)(a_{8,16}z-1)$ in \filename{scripts/a816\_certificate/a816\_lift.txt} and checked exactly by {\sc Singular} and by python-flint, proves $a_{8,16}=0$ without the descent (see~\S\ref{ab:sec:lower-edge}). Only Corollary~\ref{ab:cor:gghv} also assumes GGHV
Proposition~4.3.
\end{remark}

\subsection{Outline}

Section~\ref{ab:sec:source} source-gates all base objects against GGHV v1. Section~\ref{ab:sec:system} sets up the $72$-unknown system. Section~\ref{ab:sec:layers} gives the exact layer reduction. Section~\ref{ab:sec:normalization} fixes normalizations. Section~\ref{ab:sec:top} classifies the top layer. Section~\ref{ab:sec:descent} carries out the descent. Section~\ref{ab:sec:obstruction} proves the final obstruction. Section~\ref{ab:sec:lifting} lifts to characteristic zero. Section~\ref{ab:sec:conclusion} concludes; branch~(c) is eliminated in Part~II (Theorem~\ref{bc:thm:unconditional}). Section~\ref{ab:sec:attribution} covers prior work, attribution, and independence.

The proof is organized around four certificates. \emph{Certificate~I} (Newton reduction): $[P,Q]=\lambda x^2\Rightarrow E_5,\dots,E_1$ (\S\ref{ab:sec:layers}). \emph{Certificate~II} (top-layer completeness): $E_5\Rightarrow$ five torus orbits over $K_5$ (\S\ref{ab:sec:top}, \S\ref{ab:sec:chartproof}). \emph{Certificate~III} (descent obstruction): $E_4,E_3,E_2\Rightarrow$ the compatibility ideal equals $(t_1,t_2)^5$, hence $t=0$ (\S\ref{ab:sec:descent}, \S\ref{ab:sec:obstruction}). \emph{Certificate~IV} (vertex destruction): $t=0\Rightarrow B_0'=0\Rightarrow b_{12,24}=0$ (Corollary~\ref{ab:cor:t0}). The rigidity $a_{8,16}=0$ (Corollary~\ref{ab:cor:a816}) is a supplementary result; Theorem~\ref{ab:thm:main} does not use it. The dependency graph is
\[
\boxed{\text{Counterexample in case }(8,28)\xrightarrow{\text{GGHV 4.3}} \text{NewtonNF2},\ [P,Q]=\lambda x^2 \xrightarrow{\text{Thm.~1.1}} \text{contradiction}}
\]
with Theorem~\ref{ab:thm:main} itself independent of GGHV:
\begin{center}
\begin{tabular}{l|c}
\textbf{Result} & \textbf{Depends on GGHV?} \\ \midrule
Layer reduction & No \\
Top-layer classification & No \\
$K_5$ descent & No \\
$b_{12,24}=0$ & No \\
$a_{8,16}=0$ & No \\
Theorem~\ref{ab:thm:main} & No \\
Corollary~\ref{ab:cor:gghv} (Jacobian application) & Yes
\end{tabular}
\end{center}

\section{Source data}\label{ab:sec:source}

We verify every base object directly against version~1 of~\cite{GGHV} (arXiv:2204.14178v1), downloaded from \url{https://arxiv.org/pdf/2204.14178v1} and extracted with \texttt{pdftotext -layout}. No proof is drafted until the polygons, bracket, and $\varphi$ are confirmed character-by-character against GGHV's text.

\begin{proposition}[{\cite[Proposition~4.3]{GGHV}, source-gated}]\label{ab:prop:gghv43}
If there is a counterexample to the Jacobian conjecture in case~$(8,28)$, then there exist $P,Q\in L^{(1)}$ with $[P,Q]=x^2$ and one of:
\begin{itemize}[nosep]
\item[(1)] $N(P)=\{(0,0),(1,0),(8,14),(8,16),(0,8)\}$,\\
$N(Q)=\{(0,0),(2,1),(12,21),(12,24),(0,12)\}$;
\item[(2)] $N(P)=\{(0,0),(1,0),(8,14),(8,16)\}$,\\
$N(Q)=\{(0,0),(2,1),(12,21),(12,24)\}$.
\end{itemize}
\end{proposition}

\begin{proof}[Source locations]
The \texttt{pdftotext -layout} extraction of v1 yields:
\begin{itemize}[nosep]
\item Line~33: $[P,Q]:=P_xQ_y-P_yQ_x\in K^\times$ (bracket convention).
\item Line~492: ``Proposition~4.3 (Case $(8,28)$).''
\item Line~493: ``with $[P,Q]=x^2$'' (rendered as \texttt{x2} in the extraction).
\item Line~495: Case~(2) polygons\\
$N(P)=\{(0,0),(1,0),(8,14),(8,16)\}$, $N(Q)=\{(0,0),(2,1),(12,21),(12,24)\}$.
\item Lines~596--597: The morphism $\varphi(x)=x^{-1}$, $\varphi(y)=x^4y$; ``this is not an automorphism and the chain rule gives $[\varphi(P),\varphi(Q)]=-[P,Q]x^2$.''
\item Lines~599--600: ``in the cases a) and b) the Newton Polygons of $P$ and $Q$ become $N(P)=\{(0,0),(1,0),(8,14),(8,16)\}$, $N(Q)=\{(0,0),(2,1),(12,21),(12,24)\}$.''
\end{itemize}
A verification script (\filename{source\_gate\_verify.py}) asserts each of these lines automatically; it exits non-zero if any check fails.
\end{proof}

We work exclusively with case~(2) (branches~a,b).

\section{The polynomial system}\label{ab:sec:system}

\subsection{Lattice points}

The Newton polygons~\eqref{ab:eq:polygons} contain:
\begin{equation}
|N(P)\cap\mathbb{Z}^2|=25,\qquad |N(Q)\cap\mathbb{Z}^2|=47,
\end{equation}
verified by direct enumeration (cross-checked with Pick's theorem: $\operatorname{area}(N(P))=15$, boundary $18$, interior $7$). This gives $25+47=72$ unknown coefficients. The two constant terms do not appear in $[P,Q]$, leaving $70$ effective unknowns, but we retain all $72$ for the layer structure.

\subsection{The bracket equations}

Writing $P=\sum_{(i,j)\in N(P)}a_{ij}x^iy^j$ and $Q=\sum_{(i,j)\in N(Q)}b_{ij}x^iy^j$, the condition $[P,Q]=x^2$ expands to bilinear equations in the coefficients. Solving the full $72$-unknown bilinear system directly is impractical; we use the exact layer reduction of \S\ref{ab:sec:layers}. (Direct expansion remains feasible for \emph{verification} of the reduction, as in Proposition~\ref{ab:prop:layers}.)

\section{Exact layer reduction}\label{ab:sec:layers}

Set $u=xy^2$ and decompose into $y$-layers:
\begin{equation}\label{ab:eq:decomp}
P=\sum_{k} y^{-k}A_k(u),\qquad Q=\sum_{\ell} y^{-\ell}B_\ell(u),
\end{equation}
where $A_k,B_\ell\in\mathbb{C}[u]$ are one-variable polynomials whose coefficients are linear forms in the $a_{ij},b_{ij}$. The bracket $[P,Q]$ becomes a sum over $y$-layers, each a one-variable identity in $u$.

\begin{proposition}[Exact layer reduction]\label{ab:prop:layers}
The system $[P,Q]=x^2$ on the $72$ coefficients implies five one-variable polynomial identities $E_1,\dots,E_5$ (the converse is not needed):
\begin{center}
\begin{tabular}{c|c|l}
\textbf{Layer} & \textbf{Equations} & \textbf{Unknowns (given previous layers)} \\
\midrule
$E_5$ & $17$ & Top: $\alpha,\beta$ \\
$E_4$ & $18$ & $A_1,B_2$ given $A_2,B_3$ \\
$E_3$ & $19$ & $A_0,B_1$ given $A_1,B_2$ \\
$E_2$ & $19$ & $B_0$ given $A_0,B_1$ \\
$E_1$ & $19$ & Lowest coefficients \\
\midrule
Total & $92$ & \\
\end{tabular}
\end{center}
The layer reduction was verified as a polynomial identity (script \filename{C1\_layer\_reduction.sing}) with all $72$ coefficients treated as symbolic indeterminates; the residual is identically zero.
\end{proposition}

\begin{proof}[Computational verification]
Substituting~\eqref{ab:eq:decomp} into $[P,Q]-x^2$ and collecting by powers of $y$ yields the five identities $E_1,\dots,E_5$. Expanding each $E_i$ in the $72$ indeterminates and comparing coefficients with the direct bilinear expansion of $[P,Q]-x^2$ gives zero difference as a polynomial in $72$ variables. Script: \filename{C1\_layer\_reduction.sing} ({\sc Singular}).
\end{proof}

% BEGIN-MECHANICAL:lst:layers
\begin{leanlisting}{Mechanically extracted from \texttt{Jacobian/BranchAbLayers.lean} (sha256 \texttt{7c16141b4b126b3e}); statement verbatim, proof elided.}{ab:lst:layers}
theorem layers_of_jac (lam : K) (A₀ A₁ A₂ B₀ B₁ B₂ B₃ : Polynomial K)
    (p₀ p₁ p₂ q₀ q₁ q₂ q₃ : MvPolynomial (Fin 2) K)
    (hp₀ : X 1 ^ 0 * p₀ = evH A₀) (hp₁ : X 1 ^ 1 * p₁ = evH A₁) (hp₂ : X 1 ^ 2 * p₂ = evH A₂)
    (hq₀ : X 1 ^ 0 * q₀ = evH B₀) (hq₁ : X 1 ^ 1 * q₁ = evH B₁) (hq₂ : X 1 ^ 2 * q₂ = evH B₂)
    (hq₃ : X 1 ^ 3 * q₃ = evH B₃)
    (hJ : jac (p₀ + p₁ + p₂) (q₀ + q₁ + q₂ + q₃) = C lam * X 0 ^ 2) :
    layerTerm 2 3 A₂ B₃ = Polynomial.C lam * Polynomial.X ^ 2 /\
    layerTerm 2 2 A₂ B₂ + layerTerm 1 3 A₁ B₃ = 0 /\
    layerTerm 2 1 A₂ B₁ + layerTerm 1 2 A₁ B₂ + layerTerm 0 3 A₀ B₃ = 0 /\
    layerTerm 2 0 A₂ B₀ + layerTerm 1 1 A₁ B₁ + layerTerm 0 2 A₀ B₂ = 0 /\
    layerTerm 1 0 A₁ B₀ + layerTerm 0 1 A₀ B₁ = 0 := by
  -- proof body elided (see artifact)
\end{leanlisting}
% END-MECHANICAL:lst:layers

The basic definitions used above are printed verbatim here, since only a human reader (not the Lean kernel) checks definitions:
% BEGIN-MECHANICAL:lst:def-jac
\begin{leanlisting}{Mechanically extracted from \texttt{Jacobian/BranchAbLayers.lean} (sha256 \texttt{7c16141b4b126b3e}); definition printed verbatim in full.}{ab:lst:def-jac}
def jac (P Q : MvPolynomial (Fin 2) K) : MvPolynomial (Fin 2) K :=
  pderiv 0 P * pderiv 1 Q - pderiv 1 P * pderiv 0 Q
\end{leanlisting}
% END-MECHANICAL:lst:def-jac
% BEGIN-MECHANICAL:lst:def-uu
\begin{leanlisting}{Mechanically extracted from \texttt{Jacobian/BranchAbLayers.lean} (sha256 \texttt{7c16141b4b126b3e}); definition printed verbatim in full.}{ab:lst:def-uu}
def uu : MvPolynomial (Fin 2) K := X 0 * X 1 ^ 2
\end{leanlisting}
% END-MECHANICAL:lst:def-uu
% BEGIN-MECHANICAL:lst:def-evh
\begin{leanlisting}{Mechanically extracted from \texttt{Jacobian/BranchAbLayers.lean} (sha256 \texttt{7c16141b4b126b3e}); definition printed verbatim in full.}{ab:lst:def-evh}
def evH : Polynomial K →ₐ[K] MvPolynomial (Fin 2) K := Polynomial.aeval (uu (K := K))
\end{leanlisting}
% END-MECHANICAL:lst:def-evh
% BEGIN-MECHANICAL:lst:def-layerterm
\begin{leanlisting}{Mechanically extracted from \texttt{Jacobian/BranchAbLayers.lean} (sha256 \texttt{7c16141b4b126b3e}); definition printed verbatim in full.}{ab:lst:def-layerterm}
def layerTerm (a b : Nat) (A B : Polynomial K) : Polynomial K :=
  (a : Polynomial K) * A * Polynomial.derivative B - (b : Polynomial K) * Polynomial.derivative A * B
\end{leanlisting}
% END-MECHANICAL:lst:def-layerterm

The system is triangular (Figure~\ref{ab:fig:layers}): $E_5$ involves only the top-layer unknowns, and each $E_i$ ($i<5$) is affine-linear in its unknowns given a solution of the layers above. This is the key structural fact that makes the problem tractable. Conceptually, the Newton polygons induce a weighted grading, and the Jacobian equation respects that grading: the five $E_i$ are the homogeneous components of $[P,Q]=x^2$ with respect to it. The triangular structure is thus not a computational accident but a structural consequence of the grading, and the operator $\mathrm{LT}_{a,b}(A,B)=aAB'-bA'B$ is its homogeneous bracket.

\section{Normalization}\label{ab:sec:normalization}

The top-layer unknowns admit rescaling symmetries from the torus action. We fix, without loss of generality,
\begin{equation}\label{ab:eq:normalization}
a_{1,0}=a_{8,14}=1,
\end{equation}
leaving $a_{8,16}$ free. The remaining $1$-dimensional torus acts by the $Y$-weight grading used in the layer decomposition (\S\ref{ab:sec:layers}); the finite subgroup $\mu_7$ gives $7$ chart copies per orbit.

\begin{remark}[Lean chart normalization]
The Lean formalization uses chart coordinates $(a_1,\dots,a_7,b_0,\dots,b_9)$ with the torus normalization $a_7^3 b_0^2 = 1$ (see \texttt{ChartClassification} in \texttt{Jacobian/BranchAbChart.lean}). After fixing $\alpha_0 = 1$ and $\beta_{10} = 1$, the remaining scaling $u \mapsto \kappa u$ acts on $a_7$ with weight~$7$ and on $b_0$ with weight~$-10$, hence on $a_7^3 b_0^2$ with weight $21 - 20 = 1$. The condition $a_7^3 b_0^2 = 1$ therefore fixes $\kappa$ completely, with no $7$th-root ambiguity: the Lean chart meets each torus orbit \emph{exactly once}, giving $5$ points, not $35$. The \S\ref{ab:sec:normalization} chart ($a_{1,0} = a_{8,14} = 1$) meets each orbit $7$ times. The two charts are related by a torus element requiring a $7$th root.
\end{remark}

\emph{Vertex conditions used.}\label{ab:sec:lower-edge} The proof uses $(1,0)$ and $(2,1)$ (equivalently $\lambda\neq0$), $(8,14)$ (normalization and top-layer degree), and $(12,24)$ (the $t=0$ exclusion). The vertex $(8,16)$ is left free in \texttt{main\_theorem}. Once $E_1$ is included, $a_{8,16}=0$ is forced (Corollary~\ref{ab:cor:a816}). This gives a second contradiction, from the lower edge: \texttt{main\_theorem\_lower\_edge} (listing~\ref{ab:lst:a816-nf2}) refutes $\mathrm{NewtonNF2}$ without using the vertex $(12,24)$. It shares the descent up to $t=0$ with \texttt{main\_theorem}. The explicit certificate over $K_5$ (\filename{scripts/a816\_certificate/}, checked exactly by {\sc Singular} and by python-flint) reaches $a_{8,16}=0$ without the descent, outside Lean.

% BEGIN-MECHANICAL:lst:torus
\begin{leanlisting}{Mechanically extracted from \texttt{Jacobian/BranchAbTorus.lean} (sha256 \texttt{97918673d5c4dc09}); statement verbatim, proof elided.}{ab:lst:torus}
theorem layers_transport (α β κ : K) (A₀ A₁ A₂ B₀ B₁ B₂ B₃ : K[X])
    (E4 : layerTerm 2 2 A₂ B₂ + layerTerm 1 3 A₁ B₃ = 0)
    (E3 : layerTerm 2 1 A₂ B₁ + layerTerm 1 2 A₁ B₂ + layerTerm 0 3 A₀ B₃ = 0)
    (E2 : layerTerm 2 0 A₂ B₀ + layerTerm 1 1 A₁ B₁ + layerTerm 0 2 A₀ B₂ = 0) :
    layerTerm 2 2 (C α * sc κ A₂) (C β * sc κ B₂) + layerTerm 1 3 (C α * sc κ A₁) (C β * sc κ B₃) = 0 /\
    layerTerm 2 1 (C α * sc κ A₂) (C β * sc κ B₁) + layerTerm 1 2 (C α * sc κ A₁) (C β * sc κ B₂) +
      layerTerm 0 3 (C α * sc κ A₀) (C β * sc κ B₃) = 0 /\
    layerTerm 2 0 (C α * sc κ A₂) (C β * sc κ B₀) + layerTerm 1 1 (C α * sc κ A₁) (C β * sc κ B₁) +
      layerTerm 0 2 (C α * sc κ A₀) (C β * sc κ B₂) = 0 := by
  -- proof body elided (see artifact)
\end{leanlisting}
% END-MECHANICAL:lst:torus

\section{The top layer}\label{ab:sec:top}

\subsection{The Davenport--Stothers equation}

The top identity $E_5$ takes the form
\begin{equation}\label{ab:eq:e5}
\alpha\beta+u\,(2\alpha\beta'-3\alpha'\beta)=1,
\end{equation}
where $\alpha=A_2/u$ and $\beta=B_3/u^2$ are polynomials of degree $7$ and $10$ respectively, the degrees being forced by the vertices $(8,14)\in N(P)$ and $(12,21)\in N(Q)$.

Setting $u=s^2$ with $f(s)=\alpha(s^2)$ and $g(s)=s\beta(s^2)$, equation~\eqref{ab:eq:e5} becomes the Davenport--Stothers extremal equation~\cite{Stothers,Davenport}
\begin{equation}
2fg'-3f'g=c\neq0.
\end{equation}
Indeed $2fg'-3f'g=2(\alpha\beta+u(2\alpha\beta'-3\alpha'\beta))$. With $\deg f=14=2k$ and $\deg g=21=3k$ ($k=7$), we have $\deg(g^2-c\cdot f^3)=8=k+1$, which is exactly Davenport's bound: the pair is extremal.
with $f$ even and $g$ odd. Here $\deg_s f = 2m$ where $m=\deg_u\alpha$; the relation $1+2n-3m=0$ (with $n=\deg_u\beta$) forces $m$ odd. The cases $m=1,3,5,7$ admit solutions; $m=7$ is the relevant case for~\eqref{ab:eq:polygons}. The symmetric (parity-constrained) classification is computed here via elimination; we do not claim it follows from the literature.

\subsection{Classification}

\begin{proposition}[Top-layer solutions]\label{ab:prop:top-class}
The system $E_5$ has:
\begin{itemize}[nosep]
\item $m=3$: $3$ complex solutions; eliminant $3a_2^3-32$, irreducible over $\mathbb{Q}$;
\item $m=5$: $10$ complex solutions; eliminant $9a_4^{10}+37200a_4^5+95051008$, irreducible over $\mathbb{Q}$;
\item $m=7$: $35$ complex solutions; eliminant $\mathcal{W}(a_6)$ of degree $35$.
\end{itemize}
The $m=7$ case is the one relevant to~\eqref{ab:eq:polygons}.
\end{proposition}

\begin{proof}
We work over~$\mathbb{C}$: the Belyi classification is topological, hence valid over~$\mathbb{C}$, and the resulting count of solutions transfers to any field of characteristic~$0$ by base change (the equations are defined over~$\mathbb{Q}$). We normalize the third branch point to $c=1$ by the scaling $u\mapsto c\,u$, which preserves the torus action. We prove the counts in characteristic~$0$ via the Belyi-map classification. For the top layer, $E_5$ is $\alpha\beta+u(2\alpha\beta'-3\alpha'\beta)=1$ with $\deg\alpha=m$, $\deg\beta=(3m-1)/2$ ($m=3,5,7$, so $\deg\beta=4,7,10$). Put $\varphi=u\beta^2/\alpha^3$. Then $\varphi'=\beta E/\alpha^4=\beta/\alpha^4$, so $\varphi$ is a Belyi map: it ramifies only over $0$, $\infty$, and $c=\varphi(\infty)\neq0$. The ramification is forced as follows. At a root $r$ of $\alpha$, $E(r)=-3r\cdot\alpha'(r)\beta(r)=1$. So $\alpha$ has simple nonzero roots, none shared with $\beta$. At a root $s$ of $\beta$, $E(s)=2s\cdot\alpha(s)\beta'(s)=1$, so $\beta$ also has simple nonzero roots. Hence $\deg\varphi=3m$, and $\varphi'=\beta/\alpha^4$ vanishes in $\mathbb{C}$ only at the roots of $\beta$, which forces the passport
\[
\begin{array}{c|ccc}
m=3 & [2^4 1]\;(n=9) & [3^3] & [7\cdot 1^2]\\
m=5 & [2^7 1]\;(n=15) & [3^5] & [12\cdot 1^3]\\
m=7 & [2^{10}1]\;(n=21) & [3^7] & [17\cdot 1^4]
\end{array}
\]
By Riemann's existence theorem, covers with this passport correspond to $S_n$-conjugacy classes of transitive triples in $C_0\times C_1\times C_\infty$ with product~$1$ (here $n=3m$). The Frobenius character formula (computed exactly in \filename{belyi\_counts\_m357.py} via the Murnaghan--Nakayama rule) counts $N$ triples with $N=9!$, $2\cdot15!$, $5\cdot21!$ for $m=3,5,7$; since every such triple has trivial centralizer, each $S_n$-conjugacy class has $n!$ elements, giving $N/n!=1$, $2$, $5$ covers. Every triple counted is transitive, for general~$m$: a non-transitive triple would have an orbit on which $g_\infty$ acts trivially, so on that orbit $g_0=g_1^{-1}$; but $g_1$ has cycle type $3^m$ (hence fixes no point, and so does $g_1^{-1}$), while $g_0$ has cycle type $2^{(3m-1)/2}1$ and contains no $3$-cycle---a contradiction. Hence all counted triples are genuine covers. Each cover has trivial automorphism group: an automorphism $u\mapsto\kappa u$ acts freely on the $m$ poles and on the $(m+1)/2$ unramified points over $c$, so its order divides $\gcd(m,(m+1)/2)=1$.

Now suppose two solutions give isomorphic covers. The isomorphism fixes the unique unramified point over $0$ ($u=0$) and the unique maximally ramified point over $c$ ($u=\infty$). So it is $u\mapsto\kappa u$, and comparing zeros and poles shows the two solutions differ by the torus action. Hence distinct orbits give non-isomorphic covers, and there are at most $m\cdot(\#\text{covers})$ chart points: $3$, $10$, $35$.

For $m=7$, the exactly verified parametrization (script \filename{verify\_msolve\_param.py}, shipped with inputs \filename{e5\_m7.ms} and \filename{e5\_m7.out} in \filename{scripts/}) exhibits $35$ distinct points, so the classification is complete in characteristic~$0$. For $m=3$ and $m=5$ the Belyi bound is likewise sharp. In the normalization $\alpha_0=\beta_0=\alpha_m=1$ the system is zero-dimensional, of vector-space dimension~$3$ for $m=3$ and $10$ for $m=5$; hence the solution set is nonempty and has at most $d$ points, where $d=3$ resp.~$10$, by the dimension alone (independent of the Belyi chart). Eliminating all variables but $a_{m-1}$ gives, up to a nonzero rational factor, the stated eliminant. Each eliminant is irreducible and separable over~$\mathbb{Q}$, so the Galois group acts transitively on its $d$ distinct roots; since the system is defined over~$\mathbb{Q}$, the solution set is Galois-stable, and every root occurs as the $a_{m-1}$-coordinate of a solution. Thus there are at least $d$ distinct solutions, and with the upper bound of~$d$, exactly~$d$. All of this is checked by script \filename{e5\_exact\_counts\_m35.py} in \filename{scripts/}. Both classifications are also kernel-checked in Lean (\filename{lean/Jacobian/B26.lean}, over any field of characteristic~$0$): in this chart the system holds if and only if $T(a_{m-1})=0$ and the other unknowns are explicit functions of $a_{m-1}$ (\texttt{m3\_chart\_iff}: $8a_1=5a_2^2$; \texttt{m5\_chart\_iff}: $362a_4a_1=3a_4^5+13078$, $181a_4^2a_2=15a_4^5+32448$, $2172a_4^3a_3=861a_4^5+492128$; in both cases $\beta$ is an explicit polynomial in $a_1,\dots,a_{m-1}$), and $T$ has no repeated root (\texttt{m3\_T\_squarefree}, \texttt{m5\_T\_squarefree}). This gives exactly $d$ solutions over $\overline{\mathbb{Q}}$ without the Galois argument. The count is kernel-checked as well, over every algebraically closed field of characteristic~$0$ (\texttt{m3\_chart\_card}, \texttt{m5\_chart\_card} in \filename{lean/Jacobian/B26Count.lean}).

The $m=3$ eliminant $3a_2^3-32$ and $m=5$ eliminant $9a_4^{10}+37200a_4^5+95051008$ are verified irreducible over~$\mathbb{Q}$ by factorization (same script). Their irreducibility is also kernel-checked in Lean (\texttt{T3poly\_irreducible}, \texttt{T5poly\_irreducible} in \filename{lean/Jacobian/B26Irred.lean}). A rational root of $3a_2^3-32$ would satisfy $a_2^3=32/3$, but the $3$-adic valuation $-1$ of $32/3$ is not a multiple of~$3$; a cubic without a rational root is irreducible. For $m=5$ write the eliminant as $F(a_4^5)$ with $F(u)=9u^2+37200u+95051008=(3u+6200)^2+56611008$. This quadratic has no real root, so it is irreducible over~$\mathbb{Q}$. Let $u$ be a root of~$F$. The norm of $u$ from $\mathbb{Q}(u)$ to $\mathbb{Q}$ is $95051008/9=2^8\cdot13^5/3^2$, with $3$-adic valuation~$-2$. The norm of a fifth power is a fifth power, and $-2$ is not a multiple of~$5$, so $u$ is not a fifth power in $\mathbb{Q}(u)$. Hence $X^5-u$ is irreducible over $\mathbb{Q}(u)$, by the criterion for $X^p-a$ with $p$ prime. Let $\theta^5=u$. Then $\mathbb{Q}(\theta)$ contains $u=\theta^5$, so it equals $\mathbb{Q}(u)(\theta)$ and has degree $2\cdot5=10$ over~$\mathbb{Q}$. Since $\theta$ is a root of the eliminant, which has degree~$10$, the eliminant is irreducible. For $m=7$, the eliminating polynomial in $a_6$ has degree~$35$ (Gr\"obner elimination via {\sc Singular} \texttt{eliminate}).
\end{proof}

\begin{remark}[Status of \texttt{ChartClassification}]
\texttt{ChartClassification} (in \filename{Jacobian/BranchAbChart.lean}) encodes the $17$ equations of the $E_5$ chart system: the $16$ bilinear equations plus the torus normalization $a_7^3 b_0^2 = 1$. Up to v16 of the artifact it was an unproven hypothesis standing for the case $m=7$ of Proposition~6.1. It is now a theorem, \texttt{chartClassification\_holds} (\S\ref{ab:sec:chartproof}), so the implication $\texttt{ChartClassification}\ L \to \lnot\exists\dots$ of \texttt{main\_theorem\_of\_chart} is discharged and cannot be vacuous. The statement itself is unchanged. Its $17$ equations were verified against the \texttt{chart\_system\_Q.sing} Singular script and by the mechanical chart audit (see the computational archive). The generator of the proof in \S\ref{ab:sec:chartproof} reads this same statement from the Lean source.
\end{remark}

% BEGIN-MECHANICAL:lst:chart
\begin{leanlisting}{Mechanically extracted from \texttt{Jacobian/BranchAbChart.lean} (sha256 \texttt{15907ffeffb6d274}); statement verbatim, proof elided.}{ab:lst:chart}
theorem main_theorem_of_chart {L : Type*} [Field L] [CharZero L] (hc : ChartClassification L) :
    ¬ exists (P Q : MvPolynomial (Fin 2) L) (lam : L), lam ≠ 0 /\ NewtonNF2 P Q /\ jac P Q = MvPolynomial.C lam * MvPolynomial.X 0 ^ 2 :=
  -- proof term elided (see artifact)
\end{leanlisting}
% END-MECHANICAL:lst:chart

The two definitions on which \texttt{main\_theorem\_of\_chart} rests are printed verbatim in full:
% BEGIN-MECHANICAL:lst:def-mono
\begin{leanlisting}{Mechanically extracted from \texttt{Jacobian/BranchAbNewton.lean} (sha256 \texttt{fda6fce4c16abbe8}); definition printed verbatim in full.}{ab:lst:def-mono}
abbrev mono (i j : Nat) : Fin 2 →₀ Nat := Finsupp.single 0 i + Finsupp.single 1 j

@[simp] lemma mono_apply_zero (i j : Nat) : mono i j 0 = i := by simp [mono]
@[simp] lemma mono_apply_one (i j : Nat) : mono i j 1 = j := by simp [mono]
\end{leanlisting}
% END-MECHANICAL:lst:def-mono
% BEGIN-MECHANICAL:lst:def-innp
\begin{leanlisting}{Mechanically extracted from \texttt{Jacobian/BranchAbNewton.lean} (sha256 \texttt{fda6fce4c16abbe8}); definition printed verbatim in full.}{ab:lst:def-innp}
def inNP (m : Fin 2 →₀ Nat) : Prop := m 0 ≤ 8 /\ m 1 ≤ 2 * m 0 /\ 2 * m 0 ≤ m 1 + 2
\end{leanlisting}
% END-MECHANICAL:lst:def-innp
% BEGIN-MECHANICAL:lst:def-innq
\begin{leanlisting}{Mechanically extracted from \texttt{Jacobian/BranchAbNewton.lean} (sha256 \texttt{fda6fce4c16abbe8}); definition printed verbatim in full.}{ab:lst:def-innq}
def inNQ (m : Fin 2 →₀ Nat) : Prop := m 0 ≤ 12 /\ m 1 ≤ 2 * m 0 /\ 2 * m 0 ≤ m 1 + 3 /\ m 0 ≤ 2 * m 1
\end{leanlisting}
% END-MECHANICAL:lst:def-innq
% BEGIN-MECHANICAL:lst:def-newtonnf2
\begin{leanlisting}{Mechanically extracted from \texttt{Jacobian/BranchAbNewton.lean} (sha256 \texttt{fda6fce4c16abbe8}); definition printed verbatim in full.}{ab:lst:def-newtonnf2}
def NewtonNF2 (P Q : MvPolynomial (Fin 2) K) : Prop :=
  (forall m ∈ P.support, inNP m) /\ (forall m ∈ Q.support, inNQ m) /\
  P.coeff (mono 0 0) ≠ 0 /\ P.coeff (mono 1 0) ≠ 0 /\ P.coeff (mono 8 14) ≠ 0 /\ P.coeff (mono 8 16) ≠ 0 /\
  Q.coeff (mono 0 0) ≠ 0 /\ Q.coeff (mono 2 1) ≠ 0 /\ Q.coeff (mono 12 21) ≠ 0 /\ Q.coeff (mono 12 24) ≠ 0
\end{leanlisting}
% END-MECHANICAL:lst:def-newtonnf2
% BEGIN-MECHANICAL:lst:def-chartclass
\begin{leanlisting}{Mechanically extracted from \texttt{Jacobian/BranchAbChart.lean} (sha256 \texttt{15907ffeffb6d274}); definition printed verbatim in full.}{ab:lst:def-chartclass}
def ChartClassification (L : Type*) [Field L] : Prop :=
  forall a1 a2 a3 a4 a5 a6 a7 b0 b1 b2 b3 b4 b5 b6 b7 b8 b9 : L,
    (3 : L) * b1 + (-2 : L) * a1 * b0 = 0 →
    (5 : L) * b2 + (-5 : L) * a2 * b0 = 0 →
    (7 : L) * b3 + (2 : L) * a1 * b2 + (-3 : L) * a2 * b1 + (-8 : L) * a3 * b0 = 0 →
    (9 : L) * b4 + (4 : L) * a1 * b3 + (-1 : L) * a2 * b2 + (-6 : L) * a3 * b1 + (-11 : L) * a4 * b0 = 0 →
    (11 : L) * b5 + (6 : L) * a1 * b4 + (1 : L) * a2 * b3 + (-4 : L) * a3 * b2 + (-9 : L) * a4 * b1 + (-14 : L) * a5 * b0 = 0 →
    (13 : L) * b6 + (8 : L) * a1 * b5 + (3 : L) * a2 * b4 + (-2 : L) * a3 * b3 + (-7 : L) * a4 * b2 + (-12 : L) * a5 * b1 + (-17 : L) * a6 * b0 = 0 →
    (15 : L) * b7 + (10 : L) * a1 * b6 + (5 : L) * a2 * b5 + (-5 : L) * a4 * b3 + (-10 : L) * a5 * b2 + (-15 : L) * a6 * b1 + (-20 : L) * a7 * b0 = 0 →
    (17 : L) * b8 + (12 : L) * a1 * b7 + (7 : L) * a2 * b6 + (2 : L) * a3 * b5 + (-3 : L) * a4 * b4 + (-8 : L) * a5 * b3 + (-13 : L) * a6 * b2 + (-18 : L) * a7 * b1 = 0 →
    (19 : L) * b9 + (14 : L) * a1 * b8 + (9 : L) * a2 * b7 + (4 : L) * a3 * b6 + (-1 : L) * a4 * b5 + (-6 : L) * a5 * b4 + (-11 : L) * a6 * b3 + (-16 : L) * a7 * b2 = 0 →
    (21 : L) + (16 : L) * a1 * b9 + (11 : L) * a2 * b8 + (6 : L) * a3 * b7 + (1 : L) * a4 * b6 + (-4 : L) * a5 * b5 + (-9 : L) * a6 * b4 + (-14 : L) * a7 * b3 = 0 →
    (18 : L) * a1 + (13 : L) * a2 * b9 + (8 : L) * a3 * b8 + (3 : L) * a4 * b7 + (-2 : L) * a5 * b6 + (-7 : L) * a6 * b5 + (-12 : L) * a7 * b4 = 0 →
    (15 : L) * a2 + (10 : L) * a3 * b9 + (5 : L) * a4 * b8 + (-5 : L) * a6 * b6 + (-10 : L) * a7 * b5 = 0 →
    (12 : L) * a3 + (7 : L) * a4 * b9 + (2 : L) * a5 * b8 + (-3 : L) * a6 * b7 + (-8 : L) * a7 * b6 = 0 →
    (9 : L) * a4 + (4 : L) * a5 * b9 + (-1 : L) * a6 * b8 + (-6 : L) * a7 * b7 = 0 →
    (6 : L) * a5 + (1 : L) * a6 * b9 + (-4 : L) * a7 * b8 = 0 →
    (3 : L) * a6 + (-2 : L) * a7 * b9 = 0 →
    a7 ^ 3 * b0 ^ 2 = 1 →
    exists w : L, w ^ 5 - w ^ 4 + 3 * w ^ 3 + 3 * w ^ 2 + 26 = 0 /\
      a1 = (((512502980) / 1428159681 : L) + ((8080) / 109858437 : L) * w + ((877688) / 1428159681 : L) * w ^ (2 : Nat) + ((127760) / 84009393 : L) * w ^ (3 : Nat) + ((-531068) / 1428159681 : L) * w ^ (4 : Nat)) /\
      a2 = (((7293893720809472) / 226626674937113529 : L) + ((285855137792) / 5810940383002911 : L) * w + ((-47390789046272) / 226626674937113529 : L) * w ^ (2 : Nat) + ((-90487897751552) / 226626674937113529 : L) * w ^ (3 : Nat) + ((21860765739008) / 226626674937113529 : L) * w ^ (4 : Nat)) /\
      a3 = (((646429323397524292763648) / 323659079784278752637424249 : L) + ((-311599571772213100544) / 24896852291098365587494173 : L) * w + ((-13107803800410363756544) / 323659079784278752637424249 : L) * w ^ (2 : Nat) + ((-39360612513016248860672) / 323659079784278752637424249 : L) * w ^ (3 : Nat) + ((9886191189342570790912) / 323659079784278752637424249 : L) * w ^ (4 : Nat)) /\
      a4 = (((4014890803348499913288661860352) / 154078949379156364060580574704501523 : L) + ((-19247583898558620695724556288) / 3950742291773240104117450633448757 : L) * w + ((-147813452969167985031874871296) / 154078949379156364060580574704501523 : L) * w ^ (2 : Nat) + ((-2132931317809853248718661222400) / 154078949379156364060580574704501523 : L) * w ^ (3 : Nat) + ((580398578162396326286055178240) / 154078949379156364060580574704501523 : L) * w ^ (4 : Nat)) /\
      a5 = (((-47195336258536808324667608668861104128) / 73349781064717033648626206081592521450564721 : L) + ((-2993278732025095090465120575366889472) / 5642290851132079511432785083199424726966517 : L) * w + ((3250009871460478519137256780580519936) / 8149975673857448183180689564621391272284969 : L) * w ^ (2 : Nat) + ((-35972900451489262851635484544208994304) / 73349781064717033648626206081592521450564721 : L) * w ^ (3 : Nat) + ((4212495817303089511647840819054903296) / 24449927021572344549542068693864173816854907 : L) * w ^ (4 : Nat)) /\
      a6 = (((-35363777158142043633105970419962202288029696) / 24174276906186027491743446592090946632344039049203231 : L) + ((-24435887786194136948277618937091736054267904) / 1859559762014309807057188199391611279411079926861787 : L) * w + ((348725961520751776828947066063346243175710720) / 24174276906186027491743446592090946632344039049203231 : L) * w ^ (2 : Nat) + ((-8914342677872945930434815186362304837976064) / 1422016288599178087749614505417114507784943473482543 : L) * w ^ (3 : Nat) + ((67798985560990286741153010500606027934728192) / 24174276906186027491743446592090946632344039049203231 : L) * w ^ (4 : Nat)) /\
      a7 = (((342097029127255765817099097920847945090722114306048) / 885249425506264203827321815866695986293243258737477684336649 : L) + ((-5681054197220064630999879561744890228875177492480) / 22698703218109338559674918355556307340852391249678914982991 : L) * w + ((206885581446139393490715159158183040429422553333760) / 885249425506264203827321815866695986293243258737477684336649 : L) * w ^ (2 : Nat) + ((-63348258959870509813954058121926485026619479031808) / 885249425506264203827321815866695986293243258737477684336649 : L) * w ^ (3 : Nat) + ((25332256105915521578908408974616845866605607911424) / 885249425506264203827321815866695986293243258737477684336649 : L) * w ^ (4 : Nat)) /\
      b0 = (((2417876518990270596490221112276744402357446491679269234641011953332428638485407) / 81144154259988381448527095356464165463993157634428720446762385408 : L) + ((14711470286691815030110618242548669131976100022164901665557655244650732041478549) / 81144154259988381448527095356464165463993157634428720446762385408 : L) * w + ((617854596280970291332161410119563350111260243577025705635188990137876816776295) / 81144154259988381448527095356464165463993157634428720446762385408 : L) * w ^ (2 : Nat) + ((-3663785707199339942544458642105763338685260550327588794221990329866702931989049) / 81144154259988381448527095356464165463993157634428720446762385408 : L) * w ^ (3 : Nat) + ((283919616043133221256996475572763279716168016664606915395141666931558753453503) / 40572077129994190724263547678232082731996578817214360223381192704 : L) * w ^ (4 : Nat)) /\
      b1 = (((21312934521626399217781331192619502865459134337309977972690801492797458781) / 2064526619682179458796231817536743473030560696988314686717952 : L) + ((90786059190204626155643297641284354697393864232400683101193398378854556479) / 2064526619682179458796231817536743473030560696988314686717952 : L) * w + ((3676839220908931005614147091295150358755034032737886207680394068928684725) / 2064526619682179458796231817536743473030560696988314686717952 : L) * w ^ (2 : Nat) + ((-21087506112237214282702375786845007988504955169262767771079750788915150635) / 2064526619682179458796231817536743473030560696988314686717952 : L) * w ^ (3 : Nat) + ((1804471929942342272834364642391789791835513098528722294350400675635340349) / 1032263309841089729398115908768371736515280348494157343358976 : L) * w ^ (4 : Nat)) /\
      b2 = (((-144544033917986925185258501166756369390602635518529573940467345287) / 410368263186114060193506031816247923444210651466039296 : L) + ((575311520507585192568556318243232964279467466281194477346330819301) / 102592065796528515048376507954061980861052662866509824 : L) * w + ((54330691268897368779422216815165957896195264760328594263494900495) / 205184131593057030096753015908123961722105325733019648 : L) * w ^ (2 : Nat) + ((-174187880018379336306198263901059240047180472649226203058128140895) / 102592065796528515048376507954061980861052662866509824 : L) * w ^ (3 : Nat) + ((80166557140337989909132856542161474521909286203155803758626110033) / 410368263186114060193506031816247923444210651466039296 : L) * w ^ (4 : Nat)) /\
      b3 = (((-7425458405034303101010032795601866242474210402999635943131) / 20392339559214813347634883685912610051826122752 : L) + ((6744315542914009035390981237841878170110250243622814128087) / 20392339559214813347634883685912610051826122752 : L) * w + ((478440749235744236134872881006436261827393493416098817853) / 20392339559214813347634883685912610051826122752 : L) * w ^ (2 : Nat) + ((-3700344121490310017718636169454092008434178658425433227171) / 20392339559214813347634883685912610051826122752 : L) * w ^ (3 : Nat) + ((62233097938932572222702608499502859806885134430639402117) / 10196169779607406673817441842956305025913061376 : L) * w ^ (4 : Nat)) /\
      b4 = (((-196982124334537966834877980654129755738767580317247) / 4053408121472769419865599653882359709696 : L) + ((7356162939037843464175425549716484499362290432253) / 1013352030368192354966399913470589927424 : L) * w + ((3334214783216888283093389346861252319817155285863) / 2026704060736384709932799826941179854848 : L) * w ^ (2 : Nat) + ((-13592771536859666166850556543437368777444921613959) / 1013352030368192354966399913470589927424 : L) * w ^ (3 : Nat) + ((-1743167522426431101053174423220977563711561703271) / 4053408121472769419865599653882359709696 : L) * w ^ (4 : Nat)) /\
      b5 = (((-5945326279911459300427967186228236954680873) / 1611400923519540560385711240380416 : L) + ((-1207414796044598637778246641365546517002499) / 1611400923519540560385711240380416 : L) * w + ((147397906832220690705397488682848525927903) / 1611400923519540560385711240380416 : L) * w ^ (2 : Nat) + ((-1003046473955460175839195938091594613306305) / 1611400923519540560385711240380416 : L) * w ^ (3 : Nat) + ((-57989457976321180426899430535315371961481) / 805700461759770280192855620190208 : L) * w ^ (4 : Nat)) /\
      b6 = (((-40717564079347560855720292779927573) / 320299962217494673010720768 : L) + ((-4431956341187962499826675611271825) / 80074990554373668252680192 : L) * w + ((755747337967616823449366320710957) / 160149981108747336505360384 : L) * w ^ (2 : Nat) + ((-869471274892038676474107514591197) / 80074990554373668252680192 : L) * w ^ (3 : Nat) + ((-938838577777777016547324128345997) / 320299962217494673010720768 : L) * w ^ (4 : Nat)) /\
      b7 = (((-42258994952000509880862795) / 31833190703543877632 : L) + ((-36376188014455196253194745) / 31833190703543877632 : L) * w + ((5527319619223464764982093) / 31833190703543877632 : L) * w ^ (2 : Nat) + ((3622463246490253699729005) / 31833190703543877632 : L) * w ^ (3 : Nat) + ((-717039634870611328721643) / 15916595351771938816 : L) * w ^ (4 : Nat)) /\
      b8 = (((89529406222349457) / 6327518887936 : L) + ((-9772922162070171) / 1581879721984 : L) * w + ((8308630462897287) / 3163759443968 : L) * w ^ (2 : Nat) + ((8039403284956497) / 1581879721984 : L) * w ^ (3 : Nat) + ((-2811959945010567) / 6327518887936 : L) * w ^ (4 : Nat)) /\
      b9 = (((1461767229) / 5030912 : L) + ((82735263) / 5030912 : L) * w + ((34254981) / 5030912 : L) * w ^ (2 : Nat) + ((182339829) / 5030912 : L) * w ^ (3 : Nat) + ((-10377315) / 2515456 : L) * w ^ (4 : Nat))
\end{leanlisting}
% END-MECHANICAL:lst:def-chartclass

The top-layer classification statement and its witnesses:
% BEGIN-MECHANICAL:lst:def-toplayerclass
\begin{leanlisting}{Mechanically extracted from \texttt{Jacobian/BranchAbClassification.lean} (sha256 \texttt{29fdecd5a9fda336}); definition printed verbatim in full.}{ab:lst:def-toplayerclass}
def TopLayerClassification (L : Type*) [Field L] : Prop :=
  forall (A₂ B₃ : L[X]) (lam : L),
    (forall i, A₂.coeff i ≠ 0 → 1 ≤ i /\ i ≤ 8) → (forall i, B₃.coeff i ≠ 0 → 2 ≤ i /\ i ≤ 12) →
    A₂.coeff 1 ≠ 0 → A₂.coeff 8 ≠ 0 → B₃.coeff 2 ≠ 0 → B₃.coeff 12 ≠ 0 →
    layerTerm 2 3 A₂ B₃ = C lam * X ^ 2 →
    exists w ρ σ ε : L, w ^ 5 - w ^ 4 + 3 * w ^ 3 + 3 * w ^ 2 + 26 = 0 /\ ρ ≠ 0 /\ σ ≠ 0 /\ ε ≠ 0 /\
      (forall i, A₂.coeff i = ρ * ε ^ (i - 1) * topA w i) /\ (forall k, B₃.coeff k = σ * ε ^ (k - 2) * topB w k)
\end{leanlisting}
% END-MECHANICAL:lst:def-toplayerclass
% BEGIN-MECHANICAL:lst:def-topa
\begin{leanlisting}{Mechanically extracted from \texttt{Jacobian/BranchAbSharp.lean} (sha256 \texttt{9eb4eb229d2047f6}); definition printed verbatim in full.}{ab:lst:def-topa}
noncomputable def topA {L : Type*} [Field L] (w : L) : Nat → L
  | 1 => ((1 : L))
  | 2 => (((303) / 16 : L) + ((879) / 128 : L) * w + ((-389) / 128 : L) * w ^ (2 : Nat) + ((-51) / 128 : L) * w ^ (3 : Nat) + ((49) / 128 : L) * w ^ (4 : Nat))
  | 3 => (((315) / 16 : L) + ((693) / 8 : L) * w + ((-147) / 8 : L) * w ^ (2 : Nat) + ((-189) / 8 : L) * w ^ (3 : Nat) + ((147) / 16 : L) * w ^ (4 : Nat))
  | 4 => (((-79653) / 64 : L) + ((29463) / 64 : L) * w + ((7371) / 128 : L) * w ^ (2 : Nat) + ((-4767) / 16 : L) * w ^ (3 : Nat) + ((8001) / 128 : L) * w ^ (4 : Nat))
  | 5 => (((-33519) / 4 : L) + ((-6843) / 4 : L) * w + ((8181) / 8 : L) * w ^ (2 : Nat) + ((-1713) / 2 : L) * w ^ (3 : Nat) + ((-1305) / 8 : L) * w ^ (4 : Nat))
  | 6 => (((-8805) / 8 : L) + ((-76965) / 8 : L) * w + ((11523) / 16 : L) * w ^ (2 : Nat) + ((6927) / 4 : L) * w ^ (3 : Nat) + ((-18579) / 16 : L) * w ^ (4 : Nat))
  | 7 => (((389740) / 17 : L) + ((-51922) / 17 : L) * w + ((-48572) / 17 : L) * w ^ (2 : Nat) + ((63239) / 17 : L) * w ^ (3 : Nat) + ((-6475) / 17 : L) * w ^ (4 : Nat))
  | 8 => (((476437) / 68 : L) + ((314405) / 68 : L) * w + ((-174571) / 136 : L) * w ^ (2 : Nat) + ((1574) / 17 : L) * w ^ (3 : Nat) + ((82783) / 136 : L) * w ^ (4 : Nat))
  | _ => 0
\end{leanlisting}
% END-MECHANICAL:lst:def-topa
% BEGIN-MECHANICAL:lst:def-topb
\begin{leanlisting}{Mechanically extracted from \texttt{Jacobian/BranchAbSharp.lean} (sha256 \texttt{9eb4eb229d2047f6}); definition printed verbatim in full.}{ab:lst:def-topb}
noncomputable def topB {L : Type*} [Field L] (w : L) : Nat → L
  | 2 => ((1 : L))
  | 3 => (((101) / 8 : L) + ((293) / 64 : L) * w + ((-389) / 192 : L) * w ^ (2 : Nat) + ((-17) / 64 : L) * w ^ (3 : Nat) + ((49) / 192 : L) * w ^ (4 : Nat))
  | 4 => (((315) / 16 : L) + ((693) / 8 : L) * w + ((-147) / 8 : L) * w ^ (2 : Nat) + ((-189) / 8 : L) * w ^ (3 : Nat) + ((147) / 16 : L) * w ^ (4 : Nat))
  | 5 => (((-11379) / 8 : L) + ((4209) / 8 : L) * w + ((1053) / 16 : L) * w ^ (2 : Nat) + ((-681) / 2 : L) * w ^ (3 : Nat) + ((1143) / 16 : L) * w ^ (4 : Nat))
  | 6 => (((-695299) / 48 : L) + ((-48367) / 48 : L) * w + ((52679) / 32 : L) * w ^ (2 : Nat) + ((-45661) / 24 : L) * w ^ (3 : Nat) + ((-2211) / 32 : L) * w ^ (4 : Nat))
  | 7 => (((-309895) / 8 : L) + ((-251275) / 8 : L) * w + ((101481) / 16 : L) * w ^ (2 : Nat) + ((1811) / 4 : L) * w ^ (3 : Nat) + ((-60489) / 16 : L) * w ^ (4 : Nat))
  | 8 => (((1733348) / 13 : L) + ((-2368725) / 26 : L) * w + ((-184976) / 13 : L) * w ^ (2 : Nat) + ((1819939) / 52 : L) * w ^ (3 : Nat) + ((-569071) / 52 : L) * w ^ (4 : Nat))
  | 9 => (((18736247) / 39 : L) + ((1087133) / 13 : L) * w + ((-5656777) / 78 : L) * w ^ (2 : Nat) + ((1993384) / 39 : L) * w ^ (3 : Nat) + ((971365) / 78 : L) * w ^ (4 : Nat))
  | 10 => (((305098) / 17 : L) + ((5784066) / 17 : L) * w + ((-308335) / 17 : L) * w ^ (2 : Nat) + ((-950608) / 17 : L) * w ^ (3 : Nat) + ((754021) / 17 : L) * w ^ (4 : Nat))
  | 11 => (((-149721182) / 323 : L) + ((26391690) / 323 : L) * w + ((21560033) / 323 : L) * w ^ (2 : Nat) + ((-24378412) / 323 : L) * w ^ (3 : Nat) + ((2786935) / 323 : L) * w ^ (4 : Nat))
  | 12 => (((-91912444) / 969 : L) + ((-20715500) / 323 : L) * w + ((17474146) / 969 : L) * w ^ (2 : Nat) + ((-1623428) / 969 : L) * w ^ (3 : Nat) + ((-8783302) / 969 : L) * w ^ (4 : Nat))
  | _ => 0
\end{leanlisting}
% END-MECHANICAL:lst:def-topb

\subsection{Irreducibility and reality}

\begin{theorem}[Top-layer structure]\label{ab:thm:top-structure}
Let $\mathcal{W}\in\mathbb{Z}[a_6]$ be the degree-$35$ eliminating polynomial for $a_6$ at $m=7$. Then:
\begin{enumerate}[nosep]
\item[(i)] $\mathcal{W}$ is irreducible over $\mathbb{Q}$ ({\sc Singular} \texttt{factorize} over $\mathbb{Q}$ yields a nonzero constant times a single degree-$35$ factor);
\item[(ii)] $\mathcal{W}(a_6)=V(a_6^7)$ for a quintic $V\in\mathbb{Z}[S]$;
\item[(iii)] $V$ is separable (nonzero discriminant) and has exactly one real root, $S\approx 3208.\allowbreak612478738048$;
\item[(iv)] $\mathcal{W}$ is separable modulo $101$
($\gcd(\mathcal{W},\mathcal{W}')=1$ in $\mathbb{F}_{101}[a_6]$), so the $35$ complex solutions are distinct.
\end{enumerate}
Consequently the $35$ solutions form a single Galois orbit over $\mathbb{Q}$, comprising one real solution and $17$ complex-conjugate pairs. See Figure~\ref{ab:fig:top}. The map $a_6\mapsto a_6^7$ sends the $35$ roots of $\mathcal{W}$ onto the $5$ roots of $V$, $7$ to $1$, and both sets are single Galois orbits.
\end{theorem}

\begin{figure}[htbp]
\centering
\includegraphics[width=0.85\linewidth]{figures_ab/fig3_top_layer.png}
\caption{Top: the five roots of the quintic $V(S)$ in the complex $S$-plane, computed from the integer coefficients; the real root $S\approx3208.61$ is marked with a diamond. Bottom: all $35$ roots of $\mathcal{W}(a_6)$ in $\mathbb{C}$, colored by $S$-orbit; each $S\neq0$ lifts to $7$ values of $a_6$ via $7$th roots of unity. The dashed circle marks the real $S$-orbit. The $35$ roots form a single Galois orbit over $\mathbb{Q}$.}
\label{ab:fig:top}
\end{figure}

\begin{proof}
(i) Direct factorization in {\sc Singular} over $\mathbb{Q}$. (ii) By inspection: every exponent of $\mathcal{W}$ is divisible by $7$. (iii) Discriminant computed exactly in {\sc SymPy} (nonzero integer); real-root count via Sturm sequence gives exactly one real $S$. (iv) Direct $\gcd$ computation in $\mathbb{F}_{101}[a_6]$.
\end{proof}

\begin{remark}[Characteristic-zero/modular consistency]
Reducing the integer quintic $V$ modulo $101$ gives $58$ times the quintic obtained directly from the mod-$101$ Gr\"obner computation (equivalently, the mod-$101$ quintic is $54$ times $V$), confirming consistency between the characteristic-zero and modular routes.
\end{remark}


\subsection{An algebraic proof of the top-layer classification for $m=7$, and its formalization}\label{ab:sec:chartproof}

The proof of Proposition~\ref{ab:prop:top-class} above rests on Riemann's existence theorem, which is not in Mathlib (at least not in the version pinned by the artifact). For $m=7$ we now give a second proof of the classification it provides, in the normalized chart of \S\ref{ab:sec:normalization}: every solution is one of the points $(P_j(w),Q_k(w))$ with $R(w)=0$, the five $K_5$-conjugates. This is the form in which the main theorem uses Proposition~\ref{ab:prop:top-class}. The eliminant $\mathcal{W}$ and the cases $m=3,5$ are not reproved here. The proof is purely algebraic and valid over every field of characteristic~$0$, and it is formalized: Theorem~\ref{ab:thm:chartclass} is the Lean theorem \texttt{chartClassification\_holds} (listing~\ref{ab:lst:proofchartclass}). We work in the Lean chart of \S\ref{ab:sec:normalization}: $\alpha=1+a_1u+\dots+a_7u^7$, $\beta=b_0+b_1u+\dots+b_9u^9+u^{10}$, and $E_n=\sum_{i+k=n}(1+2k-3i)\,a_ib_k$ ($1\le n\le 16$) is the coefficient of $u^n$ in $\alpha\beta+u(2\alpha\beta'-3\alpha'\beta)$.

\begin{theorem}[\texttt{chartClassification\_holds}]\label{ab:thm:chartclass}
Let $L$ be a field of characteristic~$0$. Suppose $(a_1,\dots,a_7,b_0,\dots,b_9)\in L^{17}$ satisfies $E_1=\dots=E_{16}=0$ and $a_7^3b_0^2=1$. Then $(a_j,b_k)=(P_j(w),Q_k(w))$ for a root $w\in L$ of $R$. Here $P_j,Q_k\in\mathbb{Q}[w]$ are the coordinates in listing~\ref{ab:lst:def-chartclass}.
\end{theorem}

\begin{proof}
\emph{Step 1: reversed coordinates.} Since $a_7^3b_0^2=1$, we have $a_7\neq0$. Put $y_i=a_{7-i}/a_7$ for $1\le i\le 6$ and $y_7=1/a_7$, and give $y_i$ the weight~$i$. Define $S_k\in\mathbb{Q}[y_1,\dots,y_7]$, weighted homogeneous of weight~$k$, by
\[
(1+y_1v+\dots+y_7v^7)^{3/2}=\sum_{k\ge0}S_k\,v^k ,
\]
or, equivalently, by $S_0=1$ and $2k\,S_k=\sum_{i=1}^{7}(5i-2k)\,y_i\,S_{k-i}$. Write $\tilde\beta_k=b_{10-k}$ for $0\le k\le 10$ (so $\tilde\beta_0=1$) and $\tilde\beta_k=0$ for $k>10$; in both recursions, terms with a negative index are zero. Reindexing $E_{17-k}$ by $i\mapsto 7-i$, $k\mapsto 10-k$ gives
\[
E_{17-k}/a_7=-2k\,\tilde\beta_k+\sum_{i=1}^{7}(5i-2k)\,y_i\,\tilde\beta_{k-i}\qquad(1\le k\le 16).
\]
This is the recursion for $S_k$. Induction on $k$ gives $b_{10-k}=S_k(y)$ for $1\le k\le 10$, and
\begin{equation}\label{ab:eq:Ssystem}
S_{11}(y)=\dots=S_{16}(y)=0,\qquad S_{10}(y)^2=y_7^3 .
\end{equation}
The last equation holds because $b_0=S_{10}(y)$ and $a_7^3b_0^2=1$. In the language of \S\ref{ab:sec:top}, the reversed $\beta$ is the truncation of the power-series square root of the reversed $\alpha^3/a_7^3$. This is the Davenport--Stothers condition.

\emph{Step 2: the orbit relations.} Consider the six weighted-homogeneous relations
\begin{align*}
y_4&=\tfrac{2431}{18240}y_1^4-\tfrac{1027}{1710}y_1^2y_2+\tfrac{143}{380}y_2^2+\tfrac{41}{60}y_1y_3,\\
y_5&=\tfrac{1154351}{25313472}y_1^5-\tfrac{3833107}{18985104}y_1^3y_2+\tfrac{6512}{43947}y_1y_2^2+\tfrac{24449}{166536}y_1^2y_3,\\
y_6&=\tfrac{4206749569}{4242031637760}y_1^6+\tfrac{8589189473}{1590761864160}y_1^4y_2-\tfrac{440013827}{9819517680}y_1^2y_2^2\\&\qquad+\tfrac{7454351}{734423760}y_1^3y_3+\tfrac{253}{5054}y_2^3,\\
y_7&=-\tfrac{4539711243127}{16353031963564800}y_1^7+\tfrac{13373800868041}{6132386986336800}y_1^5y_2-\tfrac{214713769679}{37854240656400}y_1^3y_2^2\\&\qquad+\tfrac{81661753}{707800898700}y_1^4y_3+\tfrac{92333}{19483170}y_1y_2^3,\\
y_2y_3&=\tfrac{3738317}{46630080}y_1^5-\tfrac{16882087}{34972560}y_1^3y_2+\tfrac{115021}{161910}y_1y_2^2+\tfrac{2277881}{5828760}y_1^2y_3,\\
y_3^2&=\tfrac{740761013029}{10605079094400}y_1^6-\tfrac{754440973541}{1988452330200}y_1^4y_2+\tfrac{3827878261}{8182931400}y_1^2y_2^2\\&\qquad+\tfrac{137816863}{918029700}y_1^3y_3+\tfrac{289}{12635}y_2^3 .
\end{align*}
Write $g=0$ for any of them. Each satisfies
\[
y_7^3\cdot g\in(S_{11},\dots,S_{16})\subset\mathbb{Q}[y_1,\dots,y_7],
\]
with explicit cofactors of $341$ to $694$ terms, whose numerators and denominators have at most $42$ digits. Since $y_7\neq0$, every solution of~\eqref{ab:eq:Ssystem} satisfies all six relations.

The relations were found by exact interpolation on the $K_5$ point, as all weighted-homogeneous forms of weight at most~$7$ that vanish on the five orbits. The cofactors come from exact weighted-homogeneous linear algebra over $\mathbb{Q}$. Let $(g)$ be the ideal of the six relations. A Gr\"obner computation over $\mathbb{Q}$ (\filename{saturation\_check.py}) shows that $(S_{11},\dots,S_{16})\subseteq(g)$ and $(g):y_7=(g)$. With the certificates this gives
\[
(S_{11},\dots,S_{16}):y_7^\infty=(S_{11},\dots,S_{16}):y_7^3=(g).
\]
The ideal $(g)$ has Krull dimension~$1$ and needs all six generators, so it is a weighted complete intersection. It has exactly five solutions in the chart $y_1=1$ and none with $y_1=0$ other than $y=0$, in agreement with weighted B\'ezout, $4\cdot5\cdot5\cdot6\cdot6\cdot7/7!=5$. These facts explain the choice of relations but are not used in the proof.

\emph{Step 3: a resultant.} The relation for $y_7$ shows $y_7\in y_1\,\mathbb{Q}[y_1,y_2,y_3]$, so $y_1\neq0$. Put $t=y_2/y_1^2$ and $s=y_3/y_1^3$. The relation for $y_2y_3$ becomes $(t-c)\,s=N(t)$ with $c=\tfrac{2277881}{5828760}$ and $\deg N=2$, and the relation for $y_3^2$ becomes monic quadratic in~$s$. Their Sylvester resultant in $s$ is $\kappa\,m(t)$ with $\kappa\in\mathbb{Q}^\times$ and
\[
\begin{aligned}
m(t)={}&287548593020928\,t^5-688401965085696\,t^4+640652914818432\,t^3\\
&-292066554895024\,t^2+65563255857792\,t-5817852446211,
\end{aligned}
\]
which is irreducible over $\mathbb{Q}$. Hence $t-c$ is invertible modulo~$m$, and $s=q_3(t)$. The first four relations then give $y_i/y_1^i=q_i(t)$ for $4\le i\le 7$, with $q_i\in\mathbb{Q}[t]$ reduced modulo~$m$.

\emph{Step 4: the torus.} Let $z=1/y_1$. By weighted homogeneity,
\[
z^{10}S_{10}(y)=S_{10}\bigl(1,t,q_3(t),\dots,q_7(t)\bigr)\equiv\sigma(t)\pmod m\qquad\text{and}\qquad z^7y_7=q_7(t).
\] Multiplying $S_{10}^2=y_7^3$ by $z^{21}$ gives $z\,\sigma(t)^2=q_7(t)^3$, that is,
\[
y_1=\sigma(t)^2\,q_7(t)^{-3}\bmod m(t).
\]
This fixes the torus without a seventh root, as in \S\ref{ab:sec:normalization}, and makes every $y_i$ a polynomial in $t$ modulo~$m$. The substitution $w=\omega(t)$, with $\omega\in\mathbb{Q}[t]$ of degree~$4$, induces an isomorphism $\mathbb{Q}[t]/(m)\cong\mathbb{Q}[w]/(R)=K_5$ and gives $R(w)=0$ and $y_i=Y_i(w)$.

\emph{Step 5: back to the chart.} The formulas $a_7=1/y_7$ and $a_{7-i}=y_i/y_7$, together with the recursion of Step~1 for $b_{10-k}$, give $a_j=P_j(w)$ and $b_k=Q_k(w)$.
\end{proof}

\begin{lemma}[Characteristic-zero transfer]\label{ab:lem:chart-transfer}
Let $F$ be a field of characteristic~$0$. Every $F$-valued solution of the normalized $m=7$ chart system is obtained by specializing the $K_5$-point at a root $w\in F$ of $R$; its five specializations at the five roots of $R$ (over a splitting field of $R$) are the five solutions.
\end{lemma}
\begin{proof}
Steps~1--5 are proved algebraically over $\mathbb{Q}$ and hence hold over every characteristic-zero field by base change. The Lean formalization quantifies over all $L$ with \texttt{CharZero}. The five solutions are the $K_5$-conjugates $(P_j(w),Q_k(w))$ with $R(w)=0$.
\end{proof}

\begin{remark}[Roles of the two classifications]
The $m=7$ case of Proposition~\ref{ab:prop:top-class} is proved twice. The Belyi route combines the geometric upper bound (at most $35$ chart points over $\mathbb{C}$) with the $35$ exhibited points from the msolve parametrization, giving a complete classification in characteristic~$0$; it depends on Riemann's existence theorem and the computed parametrization. The algebraic certificate of this section gives an independent classification: every solution in characteristic~$0$ is one of the five $K_5$-orbits (Lemma~\ref{ab:lem:chart-transfer}). Lean's \texttt{main\_theorem} uses only the algebraic route.
\end{remark}

\begin{remark}[Formalization]\label{ab:rem:chartproof-lean}
The proof is formalized in \filename{lean/Jacobian/ChartProof/}: $12$ modules, all generated by \filename{lean/certgen/chartproof/} except the hand-written \filename{Reflect.lean}.
\begin{itemize}[nosep]
\item \emph{Steps.} There are $111$ steps. In $105$ of them a polynomial $g$ vanishes because $g=\sum_i c_ie_i$ for polynomials $e_i$ already known to vanish. In the other six, one for each relation of Step~2, the identity $y_7^3\,g=e$ holds with $e$ already known to vanish, and $y_7^3$ is cancelled using $y_7\neq0$.
\item \emph{Kernel checking.} Every identity is checked by the Lean kernel. \texttt{lc\_zero} (listing~\ref{ab:lst:lczero}) and \texttt{eq\_of\_toPolyK} reduce it to an equality of normal forms. These are computed by the verified normalizer of \texttt{Lean.Grind.CommRing} from Lean core, with a kernel-friendly multiplication defined by recursors (\texttt{toPolyK}, soundness theorem \texttt{denote\_toPolyK}), and \texttt{decide +kernel} evaluates them.
\item \emph{Axioms.} No \texttt{native\_decide} is used. \texttt{\#print axioms} for \texttt{chartClassification\_holds} and for \texttt{main\_theorem} reports only \texttt{propext}, \texttt{Classical.choice} and \texttt{Quot.sound}.
\item \emph{Batching.} The certificates of Step~2 have between $27{,}747$ and $56{,}999$ monomial products each. They are split into batches of at most $9{,}000$ products, one declaration per batch.
\item \emph{Independent checks.} Every certificate was also verified by exact SymPy expansion before it was written, and the generator reproduces all eleven generated modules byte-identically. A second checker, written independently of the generator (\filename{independent\_check.py}), parses the Lean text and re-verifies all $111$ identities over $\mathbb{Z}$. This checker verifies the arithmetic of the certificates only; it does not re-prove the mathematical statements. The Lean kernel remains solely responsible for verifying that these certificates prove the stated formal propositions. Six controls behave as expected: four one-integer perturbations are rejected by \texttt{decide}, a planted \texttt{sorry} is flagged, and an unmodified module is accepted.
\item \emph{Cost.} The build takes about $25$ minutes sequentially and needs at most $5.4$\,GB of memory for any one module.
\end{itemize}
\end{remark}

\begin{remark}[Why these coordinates]
Earlier formulations failed to yield certificates of practical size because of the degenerate Davenport--Stothers components of lower degree (the solutions for $m=1,3,5$). With $a_0=b_0=1$ and $b$ eliminated, these components force saturation by $a_7^7$ or $a_7^8$, so the certificates have weight~$61$. In the chart $a_0=b_{10}=1$, the sub-case $a_1=0$ alone needed $503{,}254$ cofactor terms modulo a prime. Replayed over~$\mathbb{Q}$, its Buchberger trace reached coefficients of $2{,}722$ digits after $929$ of its $2{,}065$ elements (\filename{lean/CLASSIFICATION\_STATUS.md}, \S5). In the reversed coordinates the same components lie on $y_7=0$, and $y_7^3$ suffices, at weight at most~$28$. The five orbits then form a weighted complete intersection that can be written down explicitly. The six certificates of Step~2 have $341$ to $694$ cofactor terms each over~$\mathbb{Q}$, roughly a thousand times fewer than the modular lift for $a_1=0$ alone.
\end{remark}

% keep the caption of the next listing on the same page as its code
\par\vskip 0pt plus 8\baselineskip\penalty-200\vskip 0pt plus -8\baselineskip\relax
% BEGIN-MECHANICAL:lst:proofchartclass
\begin{leanlisting}{Mechanically extracted from \texttt{Jacobian/ChartProof/Final.lean} (sha256 \texttt{638b644ffb58cb8b}); statement verbatim, proof elided.}{ab:lst:proofchartclass}
theorem chartClassification_holds (L : Type*) [Field L] [CharZero L] : ChartClassification L := by
  -- proof body elided (see artifact)
\end{leanlisting}
% END-MECHANICAL:lst:proofchartclass
% BEGIN-MECHANICAL:lst:unconditional
\begin{leanlisting}{Mechanically extracted from \texttt{Jacobian/ChartProof/Final.lean} (sha256 \texttt{638b644ffb58cb8b}); statement verbatim, proof elided.}{ab:lst:unconditional}
theorem main_theorem {L : Type*} [Field L] [CharZero L] :
    ¬ exists (P Q : MvPolynomial (Fin 2) L) (lam : L), lam ≠ 0 /\ NewtonNF2 P Q /\ jac P Q = MvPolynomial.C lam * MvPolynomial.X 0 ^ 2 :=
  -- proof term elided (see artifact)
\end{leanlisting}
% END-MECHANICAL:lst:unconditional
% BEGIN-MECHANICAL:lst:lczero
\begin{leanlisting}{Mechanically extracted from \texttt{Jacobian/ChartProof/Reflect.lean} (sha256 \texttt{9e3ea428ffe7236b}); statement verbatim, proof elided.}{ab:lst:lczero}
theorem lc_zero (ctx : Context α) (g : Expr) (l : List (Expr × Expr))
    (hcert : Poly.beq' (toPolyK (.sub g (lcExpr l))) (.num 0) = true) (h : AllZero ctx l) :
    g.denote ctx = 0 := by
  -- proof body elided (see artifact)
\end{leanlisting}
% END-MECHANICAL:lst:lczero

\section{Descent}\label{ab:sec:descent}

Fix a top-layer solution $(\alpha,\beta)$. The descent solves $E_4,E_3,E_2$ in turn. We use the scaled variables $A_2=u\alpha$ and $B_3=u^2\beta$.

\subsection{\texorpdfstring{Layer $E_4$}{Layer E4}}

\begin{proposition}[$E_4$ structure]\label{ab:prop:e4}
$E_4$ is a linear system in $19$ unknowns ($A_1$ on $u^1,\dots,u^8$; $B_2$ on $u^2,\dots,u^{12}$) with $18$ equations, rank $17$, and kernel dimension $2$. The kernel is parametrized by $t=(t_1,t_2)$, giving $(A_1(t),B_2(t))$ linear in $t$. Computed exactly in characteristic zero over the quintic field $K_5=\mathbb{Q}[w]/(w^5-w^4+3w^3+3w^2+26)$ (script \filename{exact\_ranks\_K5.py}); the modular evaluation at $p=101$, $a_6=55$ is an independent cross-check.
\end{proposition}

\begin{proof}
Rank computation on the $18\times19$ coefficient matrix, with residual asserts verifying the kernel. Scripts: \filename{exact\_ranks\_K5.py} (exact rank and kernel over $K_5$) and \filename{C5\_E2\_obstruction\_corrected.py} (which verifies $E_4$ with $E_4$-kernel residuals zero; mod-$101$ cross-check).
\end{proof}

\subsection{\texorpdfstring{Layer $E_3$}{Layer E3}}

\begin{proposition}[$E_3$ structure and unconditional solvability]\label{ab:prop:e3}
$E_3$ is a $19\times20$ linear system of rank $18$ with $2$-dimensional kernel. Moreover, $E_3$ is solvable for \emph{every} direction $t=(t_1,t_2)$: writing the system as $M(t)x=K(t)$ with $M$ constant and $K$ quadratic in $t$, the left nullvector $v$ (with $v^{\mathrm{T}} M=0$) satisfies
\begin{equation}
v^{\mathrm{T}} K_{11}\equiv0,\qquad v^{\mathrm{T}} K_{22}\equiv0,\qquad v^{\mathrm{T}} K_{12}\equiv0
\end{equation}
identically as polynomials in $t$, where $K=t_1^2K_{11}+t_1t_2K_{12}+t_2^2K_{22}$. Hence no direction is obstructed at $E_3$. All computations are exact over $K_5$; see Lemma~\ref{ab:lem:lifting}(i).
\end{proposition}

\begin{proof}
The $19\times20$ matrix $M$ has rank $18$; its left nullspace is $1$-dimensional, spanned by $v$. Decomposing the inhomogeneous part and computing $v^{\mathrm{T}} K_{ij}$ symbolically gives zero for all $i,j$. Script: \filename{C5\_E2\_obstruction\_corrected.py} (which verifies $E_3$ with residual asserts).
\end{proof}

The general solution of $E_3$ is thus the $2$-parameter affine family
\begin{equation}
x(t,s)=t_1^2x_{11}+t_1t_2x_{12}+t_2^2x_{22}+s_1w_1+s_2w_2,
\end{equation}
encoding $(A_0,B_1)$ in terms of $(t,s)=(t_1,t_2,s_1,s_2)\in\mathbb{C}^4$.

\subsection{\texorpdfstring{Layer $E_2$}{Layer E2}}

\begin{proposition}[$E_2$ injectivity]\label{ab:prop:e2}
The homogeneous operator of $E_2$ in the unknown $B_0'$ is a $19\times12$ matrix of rank $12$, hence injective. Its left nullspace has dimension $19-12=7$. Exact over $K_5$; see Lemma~\ref{ab:lem:lifting}(i)--(ii).
\end{proposition}

Therefore $B_0'$ (and hence $B_0$ up to an additive constant) is uniquely fixed by $(A_0,B_1,A_1,B_2)$. The seven left-nullvector pairings then give seven compatibility conditions:
\begin{equation}\label{ab:eq:seven}
W\cdot\mathrm{rhs}(t,s)=0,
\end{equation}
where $W$ is $7\times19$ and $\mathrm{rhs}(t,s)$ is the $E_2$ inhomogeneous term. These seven conditions are the \emph{only} remaining obstructions; $E_1$ is never needed.

% BEGIN-MECHANICAL:lst:descent
\begin{leanlisting}{Mechanically extracted from \texttt{Jacobian/Descent/Main.lean} (sha256 \texttt{e18a3a9e531df6d6}); statement verbatim, proof elided.}{ab:lst:descent}
theorem descent_K5 {L : Type*} [Field L] [CharZero L]
    (w : L)
    (hw : w ^ 5 - w ^ 4 + 3 * w ^ 3 + 3 * w ^ 2 + 26 = 0)
    (a_1_0 a_2_2 a_3_4 a_4_6 a_5_8 a_6_10 a_7_12 a_8_14 b_2_1 b_3_3 b_4_5 b_5_7 b_6_9 b_7_11 b_8_13 b_9_15 b_10_17 b_11_19 b_12_21 : L)
    (h_a_1_0 : a_1_0 = ((1 : L)))
    -- ... (18 h4_*, 19 h3_*, 19 h2_*, 8 h_a_*, 11 h_b_*, 51 lower-layer unknowns in one binder)
    (hv : b_12_24 ≠ 0) :
    False := by
  -- proof body elided (see artifact)
\end{leanlisting}
% END-MECHANICAL:lst:descent

\section{The final obstruction}\label{ab:sec:obstruction}

We use the homogeneous formulation: $E_3$ with $A_0$ on $u^1,\dots,u^8$ ($a_{8,16}$ free, no constant term) and $B_1$ on $u^1,\dots,u^{12}$. This gives the genuine $2$-dimensional kernel with two independent directions $s=(s_1,s_2)$, both moving actual coefficients. (An earlier formulation froze $a_{8,16}=1$ and included the constant term as an unknown; the constant is killed by $A_0'$ and contributes only a trivial direction.)

\begin{remark}[Free parameter]
\S\ref{ab:sec:normalization} leaves $a_{8,16}$ free from the start. The $E_2$/$E_3$ analysis therefore applies to every slice, including $a_{8,16}=1$, without a superset argument.
\end{remark}

\subsection{Explicit form of the seven conditions}

Expanding~\eqref{ab:eq:seven} symbolically in $(t_1,t_2,s_1,s_2)$ gives:

\begin{lemma}\label{ab:lem:seven-explicit}
Each of the seven conditions has the form
\begin{equation}
b_i(t_1,t_2)+s_1\cdot M_i(t_1,t_2)+s_2\cdot L_i(t_1,t_2)=0,\qquad i=1,\dots,7,
\end{equation}
where $b_i$ is homogeneous cubic and $M_i,L_i$ are linear in $(t_1,t_2)$. Both $s_1$ and $s_2$ appear genuinely.
\end{lemma}

\begin{proof}
Exact expansion over $K_5$ in python-flint (script \filename{exact\_obstruction\_K5.py}) of $W\cdot\mathrm{rhs}(t,s)$ with the homogeneous $E_3$ parametrization $x(t,s)=t_1^2x_{11}+t_1t_2x_{12}+t_2^2x_{22}+s_1w_1+s_2w_2$ substituted; the {\sc SymPy} expansion over $\mathbb{F}_{101}$ (script \filename{C5\_E2\_obstruction\_corrected.py}) is a cross-check.
\end{proof}

\subsection{\texorpdfstring{Elimination of $(s_1,s_2)$}{Elimination of (s1,s2)}}

For fixed $t\neq0$, the seven equations are affine-linear in $(s_1,s_2)$. A solution can exist only if the $7\times3$ matrix $[b(t)\mid M(t)\mid L(t)]$ has rank at most $2$, i.e.\ only if all $3\times3$ minors vanish:
\begin{equation}\label{ab:eq:minors}
\det\begin{pmatrix}
b_i & M_i & L_i \\ b_j & M_j & L_j \\ b_k & M_k & L_k
\end{pmatrix}=0,\qquad 1\leq i<j<k\leq7.
\end{equation}
There are $\binom{7}{3}=35$ such minors, each a homogeneous binary quintic in $(t_1,t_2)$ (cubic $\times$ linear $\times$ linear).

\begin{theorem}[No nonzero direction]\label{ab:thm:no-direction}
The $35$ quintics~\eqref{ab:eq:minors} span the degree-five homogeneous component $K_5[t_1,t_2]_5$ and generate the ideal $(t_1,t_2)^5$, i.e.\ they have Gr\"obner basis
\begin{equation}
(t_1,t_2)^5 = \left\{\, t_1^5,\; t_1^4 t_2,\; t_1^3 t_2^2,\; t_1^2 t_2^3,\; t_1 t_2^4,\; t_2^5 \,\right\}
\end{equation}
exactly over $K_5$ (the $35\times6$ coefficient matrix has rank $6$, spanning all binary quintics). Hence their common projective zero locus is empty and their only common affine zero in $\overline{K_5}^2$ is $(0,0)$: there is \emph{no} nonzero $t$ admitting $(s_1,s_2)$ satisfying the seven $E_2$ conditions. The mod-$101$ Gr\"obner computation (lex order) is an independent cross-check.
\end{theorem}

\begin{proof}
Exact computation over $K_5$: the $35$ minors are homogeneous quintics, and their $35\times6$ coefficient matrix has rank~$6$, so
\[
\operatorname{span}_{K_5}\{M_{ijk}\}=K_5[t_1,t_2]_5,
\]
the full six-dimensional space of binary quintics. Hence every degree-five monomial lies in the ideal $I$ they generate, i.e.\ $(t_1,t_2)^5\subseteq I$; conversely each generator lies in $(t_1,t_2)^5$, so $I=(t_1,t_2)^5$. The radical is $(t_1,t_2)$, so the affine variety over the algebraic closure is $\{(0,0)\}$. Script: \filename{exact\_obstruction\_K5.py}, which also writes the $35\times6$ matrix and a $6\times6$ nonzero minor to \filename{logs/k5\_minor\_certificate.json} as a machine-readable certificate. The {\sc SymPy} Gr\"obner basis over $GF(101)$ (script \filename{C5\_E2\_obstruction\_corrected.py}, saved, rerunnable; log archived) is a cross-check.
\end{proof}

\begin{remark}[Interpretation]
The top layer admits five characteristic-zero orbits, but none extends to a full solution: the $E_4$ kernel gives a two-dimensional space of directions $t\in K^2$, the $E_3$ layer imposes no further restriction, and then the seven $E_2$ compatibility conditions force $t=0$ via the $35$ quintic minors. The remaining case $t=0$ is ruled out separately (Corollary~\ref{ab:cor:t0}). In other words, the top-layer solutions exist in isolation, but none extends to a Newton-normal-form pair with $[P,Q]=\lambda x^2$.
\end{remark}

\subsection{\texorpdfstring{The case $t=0$}{The case t=0}}

\begin{lemma}[Collapse at $t=0$]\label{ab:lem:t0-collapse}
At $t=(0,0)$ the $E_2$ layer equation reduces to
\[
2A_2B_0'=0.
\]
\end{lemma}
\begin{proof}
Because the target bracket identity $[P,Q]=x^2=u^2y^{-4}$ contains only a $y^{-4}$ component, the layer equation $E_4$ at $y^{-3}$ is homogeneous:
\[
2(A_2B_2'-A_2'B_2)+(A_1B_3'-3A_1'B_3)=0.
\]
Consequently $t=(0,0)$ forces $(A_1,B_2)\equiv(0,0)$. In layer $E_3$ the quadratic driving term $K(t)$ vanishes identically at $t=0$, leaving only the homogeneous kernel. In layer $E_2$ the terms $(A_1B_1'-A_1'B_1)-2A_0'B_2$ then vanish identically, leaving $2A_2B_0'=0$.
\end{proof}

\begin{corollary}[$t=0$ impossible]\label{ab:cor:t0}
By Lemma~\ref{ab:lem:t0-collapse}, at $t=0$ we have $2A_2B_0'=0$. Since $A_2=u\alpha\neq0$, this requires $B_0'=0$, forcing $B_0(u)$ to be constant. Hence $b_{12,24}=0$; in particular, $(12,24)$ cannot be a vertex of $N(Q)$.
\end{corollary}

\begin{remark}
The Lean formalization \texttt{e2\_t0} (listing~\ref{ab:lst:t0}) proves $b_{12,24}=0$ via a \emph{different route}: a left-inverse certificate on the reduced $E_2$ equations with $t=0$ supplied as hypotheses, rather than the $2A_2B_0'=0$ argument above. Both routes reach the same conclusion.
\end{remark}

% BEGIN-MECHANICAL:lst:t0
\begin{leanlisting}{Mechanically extracted from \texttt{Jacobian/Descent/E2/T0.lean} (sha256 \texttt{3559bcbdbff42d3c}); statement verbatim, proof elided.}{ab:lst:t0}
theorem e2_t0 {L : Type*} [Field L] [CharZero L]
    (w : L)
    (hw : w ^ 5 - w ^ 4 + 3 * w ^ 3 + 3 * w ^ 2 + 26 = 0)
    (b_11_20 b_12_22 b_11_21 b_12_23 b_1_2 b_2_4 b_3_6 b_4_8 b_5_10 b_6_12 b_7_14 b_8_16 b_9_18 b_10_20 b_11_22 b_12_24 : L)
    -- 19 reduced E2 equations (r2_1_1) ... (r2_19_37), hypotheses elided for brevity
    (ht1 : b_11_20 = 0)
    (ht2 : b_12_22 = 0) :
    b_12_24 = 0 := by
  -- proof body elided (see artifact)
\end{leanlisting}
% END-MECHANICAL:lst:t0

\subsection{Proof of the main theorem}

\begin{proof}[Proof of Theorem~\ref{ab:thm:main}]
By Proposition~\ref{ab:prop:layers}, the layers satisfy $E_5,\dots,E_1$. Since $a_{1,0}\cdot a_{8,14}\cdot b_{2,1}\cdot b_{12,21}\neq 0$, Proposition~6.1 puts $(A_2,B_3)$ in the torus orbit of the $K_5$ point for some root $w$ of $R$. By torus transport (Listing~\ref{ab:lst:torus}) we may take it to be that point. Propositions~\ref{ab:prop:e4}--\ref{ab:prop:e2} and Theorem~\ref{ab:thm:no-direction} force $t=0$. At $t=0$, Lemma~\ref{ab:lem:t0-collapse} and Corollary~\ref{ab:cor:t0} give $B_0'=0$, so $b_{12,24}=0$ (in Lean, \texttt{e2\_t0}, Listing~\ref{ab:lst:t0}). The last sentence of the theorem follows because $\mathrm{NewtonNF2}$ requires $b_{12,24}\neq 0$.
\end{proof}

\begin{proof}[Proof of Corollary~\ref{ab:cor:a816}]
By the proof of Theorem~\ref{ab:thm:main} we may take the top layer to be the $K_5$ point, where $E_4,E_3,E_2$ force $t=0$; torus transport multiplies each coefficient by a nonzero scalar. At $t=0$ we have $A_1=B_2=0$ (Proposition~\ref{ab:prop:e4}). By Proposition~\ref{ab:prop:e3}, the coefficients of $A_0$ other than $a_{0,0}$ and the coefficients of $B_1$ are $K_5$-linear forms in $s=(b_{11,21},b_{12,23})$, the two free unknowns of $E_3$. In particular $a_{8,16}=\kappa\,b_{12,23}$ with $\kappa\in K_5^{\times}$, and $b_{11,21}$ occurs in $a_{1,2}$ with a nonzero coefficient. Since $A_1=0$, the layer $E_1$, $A_1B_0'-A_0'B_1=0$, reduces to $A_0'B_1=0$ in $K[u]$. If $B_1=0$, then $s=0$. If $A_0'=0$, then $a_{1,2}=a_{8,16}=0$, so again $s=0$. Hence $A_0=a_{0,0}$ and $B_1=0$; in particular $a_{8,16}=0$. Finally $B_0$ is constant by Corollary~\ref{ab:cor:t0}. In Lean (\texttt{lower\_edge\_rigidity}, generated by \filename{lean/certgen/gen\_a816.py}), the $E_1$ step uses the coefficients of $x^{17}y^{34}$, $x^{18}y^{36}$ and $x^{19}y^{38}$ in $E_1$ instead. At $t=0$ they are quadratic forms in $s$ that span all three quadratic monomials, and explicit $K_5$ multipliers give $b_{12,23}^2=0$ and $b_{11,21}^2=0$. The last one alone gives $a_{8,16}=0$: it is $12\,a_{8,15}b_{12,24}-8\,a_{8,16}b_{12,23}=-8\kappa\,b_{12,23}^2$.

A second proof of $a_{8,16}=0$, outside Lean and independent of the descent, is an explicit certificate. At the $K_5$ point, let $e_1,\dots,e_{75}$ be the coefficients of $[P,Q]-x^2$ in the layers $E_4,\dots,E_1$, as built by \filename{scripts/a816\_full.sing}; their unknowns are the $51$ lower coefficients of $P$ and $Q$ that occur. There is an explicit identity
\[
a_{8,16}^2=\sum_{k=1}^{75}H_k\,e_k,\qquad H_k\in K_5[a,b],
\]
stored in Rabinowitsch form, $1=\sum_k z^2H_k\,e_k-(1+a_{8,16}z)(a_{8,16}z-1)$, as the $76$ cofactors of \filename{scripts/a816\_certificate/a816\_lift.txt} ($3464$ terms, coefficients of at most $493$ digits). It is built from the layer structure. Give the coefficients of $A_1,B_2$ depth~$1$, those of $A_0,B_1$ depth~$2$ and those of $B_0$ depth~$3$; each equation of $E_{5-d}$ is then homogeneous of depth~$d$ and linear in the unknowns of depth~$d$. Row reduction of these linear parts in $E_4$, $E_3$, $E_2$ (ranks $17$, $18$, $12$ over $K_5$) expresses every unknown, modulo the layer ideal, as a polynomial in four free ones, $\tau=(b_{11,20},b_{12,22})$ of depth~$1$ and $\sigma=(a_{8,16},b_{11,21})$ of depth~$2$. The $32$ products of depth~$4$ formed from the reduced generators span all $14$ monomials of depth~$4$ in $K_5[\tau,\sigma]$, and lifting back gives the $H_k$. The identity is checked exactly by {\sc Singular} over $\mathbb{Q}(w)$ and independently by python-flint, which rebuilds $[P,Q]$ from $P$ and $Q$ (\filename{scripts/a816\_certificate/verify\_bundle.sh}). Hence $a_{8,16}=0$ at the $K_5$ point once $E_1$ is included. Torus transport multiplies $a_{8,16}$ by a nonzero scalar, so $a_{8,16}=0$ on the whole orbit.
\end{proof}

\begin{remark}[Further rigidity, outside the certificate]\label{ab:rem:full-rigidity}
The same span statement puts every monomial of depth at least~$4$ in $\tau,\sigma$ into the layer ideal, so each of the $51$ unknowns (the non-constant lower coefficients of $P$ and $Q$) is nilpotent modulo it: at the $K_5$ point the only solution of $E_4,\dots,E_1$ is $P=P_2+\mathrm{const}$, $Q=Q_3+\mathrm{const}$. Explicitly, let $x$ be an unknown of depth~$d$, and let $p_x\in K_5[\tau,\sigma]$ be the polynomial that the row reduction gives, with $x\equiv p_x$ modulo the layer ideal ($p_x=x$ for a free unknown). It is homogeneous of depth~$d$ with no constant term. So $x^{\lceil 4/d\rceil}\equiv p_x^{\lceil 4/d\rceil}$, and every monomial of $p_x^{\lceil 4/d\rceil}$ is divisible by a monomial of depth exactly~$4$. The certificates behind this are written out and checked exactly by two independent programs (\filename{scripts/a816\_rigidity/}): the $47$ identities $x-p_x=\sum_kD_{x,k}e_k$ for the pivot unknowns, and the $14$ identities $m=\sum_kH_{m,k}e_k$ for the monomials~$m$ of depth~$4$. The zero-set form of this statement, $P=P_2+\mathrm{const}$ and $Q=Q_3+\mathrm{const}$, is the Lean theorem \texttt{lower\_edge\_rigidity} (Corollary~\ref{ab:cor:a816}); the nilpotency modulo the layer ideal is not part of the Lean formalization. Neither is used elsewhere.
\end{remark}

% BEGIN-MECHANICAL:lst:a816
\begin{leanlisting}{Mechanically extracted from \texttt{Jacobian/A816/Final.lean} (sha256 \texttt{b7f3ed6341617a44}); statement verbatim, proof elided.}{ab:lst:a816}
theorem lower_edge_rigidity {L : Type*} [Field L] [CharZero L] (P Q : MvPolynomial (Fin 2) L) (lam : L)
    (hP : forall m ∈ P.support, inNP m) (hQ : forall m ∈ Q.support, inNQ m)
    (v10 : P.coeff (mono 1 0) ≠ 0) (v814 : P.coeff (mono 8 14) ≠ 0)
    (w21 : Q.coeff (mono 2 1) ≠ 0) (w1221 : Q.coeff (mono 12 21) ≠ 0)
    (hJ : jac P Q = MvPolynomial.C lam * MvPolynomial.X 0 ^ 2) :
    (forall i j, 2 * i ≤ j + 1 → (i, j) ≠ (0, 0) → P.coeff (mono i j) = 0) /\
      (forall k l, 2 * k ≤ l + 2 → (k, l) ≠ (0, 0) → Q.coeff (mono k l) = 0) := by
  -- proof body elided (see artifact)
\end{leanlisting}
% END-MECHANICAL:lst:a816
% BEGIN-MECHANICAL:lst:a816-nf2
\begin{leanlisting}{Mechanically extracted from \texttt{Jacobian/A816/Final.lean} (sha256 \texttt{b7f3ed6341617a44}); statement verbatim, proof elided.}{ab:lst:a816-nf2}
theorem main_theorem_lower_edge {L : Type*} [Field L] [CharZero L] :
    ¬ exists (P Q : MvPolynomial (Fin 2) L) (lam : L), lam ≠ 0 /\ NewtonNF2 P Q /\ jac P Q = MvPolynomial.C lam * MvPolynomial.X 0 ^ 2 := by
  -- proof body elided (see artifact)
\end{leanlisting}
% END-MECHANICAL:lst:a816-nf2

\begin{proof}[Proof of Corollary~\ref{ab:cor:gghv}]
Proposition~4.3 of~\cite{GGHV}, then Theorem~\ref{ab:thm:main}.
\end{proof}

% BEGIN-MECHANICAL:lst:main
\begin{leanlisting}{Mechanically extracted from \texttt{Jacobian/BranchAbFinal.lean} (sha256 \texttt{cd4e207adba5cb7b}); statement verbatim, proof elided.}{ab:lst:main}
theorem no_completion_K5 {L : Type*} [Field L] [CharZero L] (w : L)
    (hw : w ^ 5 - w ^ 4 + 3 * w ^ 3 + 3 * w ^ 2 + 26 = 0) (lam : L)
    (A₀ A₁ A₂ B₀ B₁ B₂ B₃ : Polynomial L) (p₀ p₁ p₂ q₀ q₁ q₂ q₃ : MvPolynomial (Fin 2) L)
    (hp₀ : X 1 ^ 0 * p₀ = evH A₀) (hp₁ : X 1 ^ 1 * p₁ = evH A₁) (hp₂ : X 1 ^ 2 * p₂ = evH A₂)
    (hq₀ : X 1 ^ 0 * q₀ = evH B₀) (hq₁ : X 1 ^ 1 * q₁ = evH B₁) (hq₂ : X 1 ^ 2 * q₂ = evH B₂)
    (hq₃ : X 1 ^ 3 * q₃ = evH B₃)
    (hJ : jac (p₀ + p₁ + p₂) (q₀ + q₁ + q₂ + q₃) = C lam * X 0 ^ 2)
    (sA₀ : forall i, A₀.coeff i ≠ 0 → 0 ≤ i /\ i ≤ 8)
    (sA₁ : forall i, A₁.coeff i ≠ 0 → 1 ≤ i /\ i ≤ 8)
    (sA₂ : forall i, A₂.coeff i ≠ 0 → 1 ≤ i /\ i ≤ 8)
    (sB₀ : forall i, B₀.coeff i ≠ 0 → 0 ≤ i /\ i ≤ 12)
    (sB₁ : forall i, B₁.coeff i ≠ 0 → 1 ≤ i /\ i ≤ 12)
    (sB₂ : forall i, B₂.coeff i ≠ 0 → 2 ≤ i /\ i ≤ 12)
    (sB₃ : forall i, B₃.coeff i ≠ 0 → 2 ≤ i /\ i ≤ 12)
    (ρ σ ε : L) (hρ : ρ ≠ 0) (hσ : σ ≠ 0) (hε : ε ≠ 0)
    (t_a_1_0 : (A₂.coeff 1) = ρ * ε ^ (0 : Nat) * ((1 : L)))
    (t_a_2_2 : (A₂.coeff 2) = ρ * ε ^ (1 : Nat) * (((303) / 16 : L) + ((879) / 128 : L) * w + ((-389) / 128 : L) * w ^ (2 : Nat) + ((-51) / 128 : L) * w ^ (3 : Nat) + ((49) / 128 : L) * w ^ (4 : Nat)))
    (t_a_3_4 : (A₂.coeff 3) = ρ * ε ^ (2 : Nat) * (((315) / 16 : L) + ((693) / 8 : L) * w + ((-147) / 8 : L) * w ^ (2 : Nat) + ((-189) / 8 : L) * w ^ (3 : Nat) + ((147) / 16 : L) * w ^ (4 : Nat)))
    (t_a_4_6 : (A₂.coeff 4) = ρ * ε ^ (3 : Nat) * (((-79653) / 64 : L) + ((29463) / 64 : L) * w + ((7371) / 128 : L) * w ^ (2 : Nat) + ((-4767) / 16 : L) * w ^ (3 : Nat) + ((8001) / 128 : L) * w ^ (4 : Nat)))
    (t_a_5_8 : (A₂.coeff 5) = ρ * ε ^ (4 : Nat) * (((-33519) / 4 : L) + ((-6843) / 4 : L) * w + ((8181) / 8 : L) * w ^ (2 : Nat) + ((-1713) / 2 : L) * w ^ (3 : Nat) + ((-1305) / 8 : L) * w ^ (4 : Nat)))
    (t_a_6_10 : (A₂.coeff 6) = ρ * ε ^ (5 : Nat) * (((-8805) / 8 : L) + ((-76965) / 8 : L) * w + ((11523) / 16 : L) * w ^ (2 : Nat) + ((6927) / 4 : L) * w ^ (3 : Nat) + ((-18579) / 16 : L) * w ^ (4 : Nat)))
    (t_a_7_12 : (A₂.coeff 7) = ρ * ε ^ (6 : Nat) * (((389740) / 17 : L) + ((-51922) / 17 : L) * w + ((-48572) / 17 : L) * w ^ (2 : Nat) + ((63239) / 17 : L) * w ^ (3 : Nat) + ((-6475) / 17 : L) * w ^ (4 : Nat)))
    (t_a_8_14 : (A₂.coeff 8) = ρ * ε ^ (7 : Nat) * (((476437) / 68 : L) + ((314405) / 68 : L) * w + ((-174571) / 136 : L) * w ^ (2 : Nat) + ((1574) / 17 : L) * w ^ (3 : Nat) + ((82783) / 136 : L) * w ^ (4 : Nat)))
    (t_b_2_1 : (B₃.coeff 2) = σ * ε ^ (0 : Nat) * ((1 : L)))
    (t_b_3_3 : (B₃.coeff 3) = σ * ε ^ (1 : Nat) * (((101) / 8 : L) + ((293) / 64 : L) * w + ((-389) / 192 : L) * w ^ (2 : Nat) + ((-17) / 64 : L) * w ^ (3 : Nat) + ((49) / 192 : L) * w ^ (4 : Nat)))
    (t_b_4_5 : (B₃.coeff 4) = σ * ε ^ (2 : Nat) * (((315) / 16 : L) + ((693) / 8 : L) * w + ((-147) / 8 : L) * w ^ (2 : Nat) + ((-189) / 8 : L) * w ^ (3 : Nat) + ((147) / 16 : L) * w ^ (4 : Nat)))
    (t_b_5_7 : (B₃.coeff 5) = σ * ε ^ (3 : Nat) * (((-11379) / 8 : L) + ((4209) / 8 : L) * w + ((1053) / 16 : L) * w ^ (2 : Nat) + ((-681) / 2 : L) * w ^ (3 : Nat) + ((1143) / 16 : L) * w ^ (4 : Nat)))
    (t_b_6_9 : (B₃.coeff 6) = σ * ε ^ (4 : Nat) * (((-695299) / 48 : L) + ((-48367) / 48 : L) * w + ((52679) / 32 : L) * w ^ (2 : Nat) + ((-45661) / 24 : L) * w ^ (3 : Nat) + ((-2211) / 32 : L) * w ^ (4 : Nat)))
    (t_b_7_11 : (B₃.coeff 7) = σ * ε ^ (5 : Nat) * (((-309895) / 8 : L) + ((-251275) / 8 : L) * w + ((101481) / 16 : L) * w ^ (2 : Nat) + ((1811) / 4 : L) * w ^ (3 : Nat) + ((-60489) / 16 : L) * w ^ (4 : Nat)))
    (t_b_8_13 : (B₃.coeff 8) = σ * ε ^ (6 : Nat) * (((1733348) / 13 : L) + ((-2368725) / 26 : L) * w + ((-184976) / 13 : L) * w ^ (2 : Nat) + ((1819939) / 52 : L) * w ^ (3 : Nat) + ((-569071) / 52 : L) * w ^ (4 : Nat)))
    (t_b_9_15 : (B₃.coeff 9) = σ * ε ^ (7 : Nat) * (((18736247) / 39 : L) + ((1087133) / 13 : L) * w + ((-5656777) / 78 : L) * w ^ (2 : Nat) + ((1993384) / 39 : L) * w ^ (3 : Nat) + ((971365) / 78 : L) * w ^ (4 : Nat)))
    (t_b_10_17 : (B₃.coeff 10) = σ * ε ^ (8 : Nat) * (((305098) / 17 : L) + ((5784066) / 17 : L) * w + ((-308335) / 17 : L) * w ^ (2 : Nat) + ((-950608) / 17 : L) * w ^ (3 : Nat) + ((754021) / 17 : L) * w ^ (4 : Nat)))
    (t_b_11_19 : (B₃.coeff 11) = σ * ε ^ (9 : Nat) * (((-149721182) / 323 : L) + ((26391690) / 323 : L) * w + ((21560033) / 323 : L) * w ^ (2 : Nat) + ((-24378412) / 323 : L) * w ^ (3 : Nat) + ((2786935) / 323 : L) * w ^ (4 : Nat)))
    (t_b_12_21 : (B₃.coeff 12) = σ * ε ^ (10 : Nat) * (((-91912444) / 969 : L) + ((-20715500) / 323 : L) * w + ((17474146) / 969 : L) * w ^ (2 : Nat) + ((-1623428) / 969 : L) * w ^ (3 : Nat) + ((-8783302) / 969 : L) * w ^ (4 : Nat)))
    (hv : B₀.coeff 12 ≠ 0) : False := by
  -- proof body elided (see artifact)
\end{leanlisting}
% END-MECHANICAL:lst:main

\section{Exact characteristic-\texorpdfstring{$0$}{0} descent}\label{ab:sec:lifting}

The proof is exact over $K_5$ (items (i)--(ii) below).

\begin{lemma}[Base change for the $K_5$ computations]\label{ab:lem:basechange}
$R(w)=w^5-w^4+3w^3+3w^2+26$ is irreducible and separable over $\mathbb{Q}$, so $K_5=\mathbb{Q}[w]/(R)$ is a number field. The coefficients of the $E_4,E_3,E_2$ matrices and of the seven compatibility polynomials lie in $K_5$. Exact rank identities and exact ideal identities over $K_5$ therefore remain valid under every characteristic-zero field extension.
\end{lemma}

\begin{proof}
$R$ is irreducible modulo $67$ (a single quintic factor), hence irreducible over $\mathbb{Q}$, and separable since its discriminant $784{,}822{,}272=2^{12}\cdot3\cdot13\cdot17^3$ is nonzero. Hence $K_5$ is a field. Rank of a matrix and equality of ideals are preserved under field extension: rank is the size of the largest nonzero minor, and ideal membership is witnessed by explicit polynomial combinations, both of which persist after extending scalars.
\end{proof}

\begin{lemma}[Exact obstruction]\label{ab:lem:lifting}
Let $\rho\in\overline{\mathbb{Q}}$ be any root of $\mathcal{W}$, and put $K=\mathbb{Q}(\rho)$. The $E_2$ obstruction---no nonzero $(t_1,t_2)$ satisfies the seven compatibility conditions---holds over $K$.
\end{lemma}

\begin{proof}
By Lemma~\ref{ab:lem:basechange}, the exact rank and ideal identities over $K_5$ persist under every characteristic-zero field extension. Since $K_5\cong\mathbb{Q}(\rho^7)$ for every root $\rho$ of $\mathcal{W}$ (the $S=a_6^7$ parametrization; Remark~\ref{ab:rem:import}), the obstruction holds over $\mathbb{Q}(\rho)$, hence over every characteristic-zero field containing such a root. The ingredients are the exact characteristic-$0$ computations (scripts \filename{exact\_ranks\_K5.py}, \filename{exact\_obstruction\_K5.py}):
\begin{enumerate}[nosep]
\item[(i)] \emph{Exact characteristic-$0$ ranks.} The ranks are computed exactly over the quintic field $K_5=\mathbb{Q}[w]/(w^5-w^4+3w^3+3w^2+26)$ (script \filename{exact\_ranks\_K5.py}): $E_4=17$, $E_3=18$, $E_2=12$.
\item[(ii)] \emph{Exact descent over $K_5$.} The entire $E_4\to E_3\to E_2\to35$-minor pipeline is redone exactly over $K_5$ (script \filename{exact\_obstruction\_K5.py}, $\sim$1~s): $E_3$ is solvable for every $t$, the left nullity of $E_2$ is $7$, and the $35$ minors span \emph{all} binary quintics (the $35\times6$ coefficient matrix has rank~$6$ over $K_5$). Hence the only common zero is $t=0$. End-to-end residuals against the original bracket equations vanish exactly. The $t=0$ case is excluded by the vertex $(12,24)$ argument, which is valid in every characteristic other than $2$ and $3$ (it uses $2\ne0$ and $12\ne0$; see \S8.3). The Lean formalization (\texttt{no\_completion\_K5}, listing~\ref{ab:lst:main}, via \texttt{main\_theorem\_of\_chart}, listing~\ref{ab:lst:chart}) quantifies over all roots of $R$ and all torus parameters, covering all top-layer solutions by Proposition~6.1, whose $m=7$ classification is formalized (\texttt{chartClassification\_holds}, \S\ref{ab:sec:chartproof}).
\end{enumerate}
\end{proof}

\begin{remark}[On the derivation of $\mathcal{W}$]\label{ab:rem:import}
The polynomial $\mathcal{W}$ was originally imported from the reference data (see \filename{W\_IMPORT\_DISCLOSURE.md}). It is now \emph{independently derived}: the minimal polynomial of $a_6^7$, computed from the exactly verified top-layer point over $K_5$, equals the quintic $V$ (hence $\mathcal{W}(a_6)=V(a_6^7)$) coefficient for coefficient (script \texttt{compare\_V.gp}). The import caveat in \filename{W\_IMPORT\_DISCLOSURE.md} is retired for $\mathcal{W}$ itself. The mod-$101$ reduction matches the independently computed elimination up to scalar $54$ as a cross-check.
\end{remark}

\section{Conclusion}\label{ab:sec:conclusion}

We have eliminated branches~(a) and~(b) of the GGHV $(72,108)$ candidate---case~(2) of~\cite[Proposition~4.3]{GGHV}---by a computational audit:
\begin{itemize}[nosep]
\item source-gated base objects (polygons, bracket, and $\varphi$ verified against GGHV v1);
\item exact layer reduction ($92$ equations, symbolic identity);
\item top-layer determination (degree-$35$ eliminant, one real solution);
\item descent with explicit ranks and unconditional $E_3$ solvability;
\item minor-span obstruction ($(t_1,t_2)^5$) closing the $E_2$ conditions;
\item exact $K_5$-descent (Lean \texttt{no\_completion\_K5} via \texttt{main\_theorem\_of\_chart}) covering all roots of $R$ and all torus parameters;
\item an algebraic proof of the $m=7$ top-layer classification (Proposition~6.1 in the normalized chart), using reversed coordinates, six orbit relations and a Sylvester resultant (\S\ref{ab:sec:chartproof}). It is formalized in Lean, so the NewtonNF2 form of Theorem~\ref{ab:thm:main} is the kernel-checked theorem \texttt{main\_theorem}, independent of GGHV Proposition~4.3;
\item lower-edge rigidity (Corollary~\ref{ab:cor:a816}): $P=P_2+a_{0,0}$ and $Q=Q_3+b_{0,0}$, kernel-checked as \texttt{lower\_edge\_rigidity}, so that the vertex $(8,16)$ refutes $\mathrm{NewtonNF2}$ as well (\texttt{main\_theorem\_lower\_edge}).
\end{itemize}

\section{Prior work, attribution, and independence}\label{ab:sec:attribution}

\emph{Prior work.} F.~Santiba\~nez-Leal's CAOS records EXP-061 and EXP-065 treat exactly these cases with sampled certificates, later downgraded by his own audit. We cite them here for completeness.

\emph{Attribution.} The layer reduction strategy ($u=xy^2$, five triangular layers), the identification of the Davenport--Stothers top layer, the descent architecture, and the modular certificate approach originate in the initial computational package (\filename{branch\_ab\_elimination.zip}, MD5 \texttt{0cfa49d0\dots}). Its full-formulation certificate at $p=10^9+9$ and $p=10^9+21$, with lifting argument, is a second, independent route to the same obstruction (see the Appendix). This work reimplements the reduction independently in {\sc Singular} (C1) and recomputes the descent, but the overall strategy is not original to us.

\emph{Characteristic-$0$ core.} The characteristic-zero computations---the Belyi counts (\filename{belyi\_counts\_m357.py}), the exact top-layer point over $K_5$ (\filename{e5\_exact\_K5.json}), the derivation of $\mathcal{W}$ (\texttt{compare\_V.gp}), the exact ranks (\filename{exact\_ranks\_K5.py}), the exact $E_4\to E_3\to E_2$ descent (\filename{exact\_obstruction\_K5.py}), and the msolve-parametrization check (\filename{verify\_msolve\_param.py})---are documented in the scripts (see their attribution headers). Tools: python-flint~0.9.0 (exact $K_5$ arithmetic), sympy~1.14.0, PARI/GP~2.15.4, Singular~4.3.2. The independence case rests on the artifacts, not on who ran them: the Lean kernel checks every formalized proof; the Singular Gr\"obner computations (\filename{a816\_full.sing}, \filename{chart\_system\_Q.sing}) can be rerun from the archive; the $a_{8,16}$ certificate (\filename{scripts/a816\_certificate/}) is checked by two independent programs; and the Belyi counts were verified by an independent Murnaghan--Nakayama implementation.

\emph{Independent AI-assisted computational audit (completed 2026-09-25).} An independent AI-assisted computational audit reached the same conclusion by an independent route without access to the present scripts. The auditing agent received only a short brief---the degree pair $(72,108)$ and the instruction to work from GGHV Proposition~4.3---plus GGHV v1 itself. It independently produced a proof using weight grading, a Belyi-map count (at most $35$ complex points via Riemann's existence theorem and the Frobenius character formula), exact modular certificates at $p=32003$ and $p=536870909$ (Singular, Macaulay2, msolve), and Hensel lifting, reaching the same nonexistence conclusion and the stronger lower-edge rigidity $a_{8,16}=b_{12,24}=0$. Materials: \href{https://github.com/brandonmccraryresearch-cloud/2D-Jacobian-Conjecture-Workspace/tree/main/tier1_blind_gghv43_nf2}{\textcolor{magenta}{\filename{tier1\_blind\_gghv43\_nf2/}}}~\cite{ClaudeBlind}, committed as \texttt{37c73b67} on 2026-09-25. The evidential weight rests on the reproduced artifacts and computations, not on the identity of the auditing system.

\begin{remark}[What remains]
The application of Theorem~\ref{ab:thm:main} to the Jacobian conjecture, Corollary~\ref{ab:cor:gghv}, is conditional on~\cite[Proposition~4.3]{GGHV} (Remark~\ref{ab:rem:conditional}). Branch~(c) (case~(1) of Proposition~\ref{ab:prop:gghv43}) is eliminated in Part~II, unconditionally in Lean (Theorem~\ref{bc:thm:unconditional}). What remains conditional, for any application to the Jacobian conjecture, is only the correctness and exhaustiveness of~\cite[Proposition~4.3]{GGHV} itself, which this paper does not re-derive.
\end{remark}

\part{\texorpdfstring{Elimination of the GGHV Branch-(c) Normal Form at Degree Pair (72,108)}{Elimination of the GGHV Branch-(c) Normal Form at Degree Pair (72,108)}}

\section{Introduction}

\subsection{Context and prior work}

The planar Jacobian conjecture asserts that a polynomial map $(P,Q):\mathbb{C}^2\to\mathbb{C}^2$ with $[P,Q]:=P_xQ_y-P_yQ_x\in\mathbb{C}^\times$ is invertible. In~\cite{GGHV}, Guccione, Guccione, Horruitiner, and Valqui raise the lower bound on $\max\{\deg P,\deg Q\}$ for a hypothetical counterexample from $100$ (Moh~\cite{Moh}) to $108$, leaving open only the degree pair $(72,108)$ and its symmetric $(108,72)$.

For the degree pair $(72,108)$, \cite{GGHV} (an unrefereed preprint, arXiv:2204.14178v1) identifies two cases. The corner $(9,27)$ is discarded in \cite[\S5]{GGHV}; the open case is the corner $(8,28)$ with $(m,n)=(3,2)$. For this case, \cite[Proposition~4.3]{GGHV} shows that any counterexample gives rise, after an explicit (non-invertible) morphism $\varphi$, to one of two Newton polygon configurations. Case~(2), branches~(a) and~(b), is eliminated in~\cite{BranchABv19}. Case~(1), the subject of this paper, is branch~(c):
\begin{equation}\label{bc:eq:polygons}
\begin{aligned}
N(P)&=\operatorname{conv}\{(0,0),(1,0),(8,14),(8,16),(0,8)\},\\
N(Q)&=\operatorname{conv}\{(0,0),(2,1),(12,21),(12,24),(0,12)\}.
\end{aligned}
\end{equation}
with $[P,Q]=\lambda x^2$ in post-$\varphi$ coordinates and $\lambda\ne0$. The additional vertices $(0,8)$ and $(0,12)$ distinguish branch~(c) from branches~(a,b). See Figure~\ref{bc:fig:polygons} for the polygon shapes.

\begin{figure}[htbp]
\centering
\includegraphics[width=0.95\textwidth]{figures_c/fig1_newton_polygons.pdf}
\caption{Branch-(c) Newton polygons, rendered from the exact vertex data of~\eqref{bc:eq:polygons} (source-gated at GGHV v1 lines~494 and~617--618; see the validation note). In each panel the branch-(a,b) sub-polygon is drawn in the lighter shade underneath; the extra branch-(c) vertices $(0,8)$ and $(0,12)$ are highlighted in the accent color. All lattice points are marked ($61$ for $N(P)$, $125$ for $N(Q)$); the ten vertex coefficients required nonzero by \texttt{NewtonNFc} are labelled. Rendered by \texttt{figures/fig1\_newton\_polygons.py} from the vertex lists.}
\label{bc:fig:polygons}
\end{figure}

Santib\'a\~nez-Leal has published machine evidence related to the $(72,108)$ candidate~\cite{Santibanez,SantibanezPaper}, explicitly not claiming elimination. The candidate itself is documented in~\cite{JC125note}. Independent non-peer-reviewed treatments of the GGHV configuration in 2026 include Helali's exact characteristic-zero certificate~\cite{Helali} and the structural analysis of wstrinz~\cite{wstrinz}; none is peer-reviewed, and the parity-constrained Davenport--Stothers equation relevant to the top layer remains unclassified.

\subsection{Main result}

\begin{definition}[Branch-(c) Newton normal form]\label{bc:def:newtonnfc}
Let $L$ be a field of characteristic~$0$ and $P,Q\in L[x,y]$. \texttt{NewtonNFc}$(P,Q)$ holds when $\operatorname{supp}P\subset N(P)$ and $\operatorname{supp}Q\subset N(Q)$ with $N(P),N(Q)$ as in~\eqref{bc:eq:polygons}, and the ten vertex coefficients (at $(0,0),(1,0),(8,14),(8,16),(0,8)$ of $P$ and $(0,0),(2,1),(12,21),(12,24),(0,12)$ of $Q$) are all nonzero. In Lean this is \texttt{NewtonNFc} (\filename{Jacobian/BranchC/DescentClaim.lean}); see listing~\ref{bc:lst:newtonnfc}.
\end{definition}

\begin{theorem}[Elimination of the GGHV branch-(c) normal form, conditional form]\label{bc:thm:main}
Let $L$ be a field of characteristic~$0$ satisfying \texttt{ChartEmptyC\_T1ne0} (listing~\ref{bc:lst:chartempty}): the twelve chart conditions of \S\ref{bc:sec:conditions} have no common zero with $b_{12,22}=1$ on the stratum $b_{11,20}\ne0$. Then there are no $P,Q\in L[x,y]$ and no $\lambda\ne0$ with \texttt{NewtonNFc}$(P,Q)$ and $[P,Q]=\lambda x^2$.
\end{theorem}

The hypothesis \texttt{ChartEmptyC\_T1ne0} is a precise Lean \texttt{Prop}. It was first established outside Lean in \S\ref{bc:sec:wholechart}, at epistemic grade~B (see \S\ref{bc:sec:epistemic}); since November 4, 2026 it is proved inside Lean with no hypotheses (\texttt{challenge\_target}), and Theorem~\ref{bc:thm:unconditional} discharges it, upgrading this conditional form to an unconditional elimination.

\begin{corollary}[Unconditional elimination (Lean certificate)]\label{bc:cor:unconditional}
\texttt{ChartEmptyC\_T1ne0} holds for every field $L$ of characteristic~$0$ with no hypotheses (\texttt{challenge\_target}, Theorem~\ref{bc:thm:unconditional}), and hence no $P,Q,\lambda$ as in Theorem~\ref{bc:thm:main} exist. The rank route of \S\ref{bc:sec:wholechart} (the rank lemma, Lemma~\ref{bc:lem:rank}, instantiated at two primes by two independent implementations, with negative controls) supplies the same conclusion independently at grade~B.
\end{corollary}

\begin{corollary}[Jacobian application]\label{bc:cor:gghv}
Assume \cite[Proposition~4.3]{GGHV} is correct and exhaustive. Then no counterexample to the planar Jacobian conjecture of degree pair $(72,108)$ arises from branch~(c).
\end{corollary}

\begin{theorem}[Unconditional elimination of branch~(c) in Lean]\label{bc:thm:unconditional}
Let $L$ be a field of characteristic~$0$. There are no $P,Q\in L[x,y]$ and no $\lambda\ne0$ with \texttt{NewtonNFc}$(P,Q)$ and $[P,Q]=\lambda x^2$.
\end{theorem}

\begin{leanlisting}{\texttt{branch\_c\_eliminated}: unconditional branch-(c) elimination}{bc:lst:unconditional}
theorem branch_c_eliminated
    {L : Type*} [Field L] [CharZero L] :
    ¬ ∃ (P Q : MvPolynomial (Fin 2) L) (lam : L),
      lam ≠ 0 ∧ NewtonNFc P Q ∧
      jac P Q = C lam * X 0 ^ 2
\end{leanlisting}

\emph{Proof sketch; kernel record.} Apply \texttt{main\_theorem\_c\_of\_chartEmpty\_T1ne0} (the conditional form, Theorem~\ref{bc:thm:main}) to \texttt{challenge\_target}. The latter is the Fpoly $=$ Fpoly2 bridge plus instantiation of \texttt{chartEmptyC\_T1ne0\_of\_certificate}: the chart \texttt{Fpoly} and the comparator's \texttt{Fpoly2} are the same six \texttt{polyOfTerms} generators in the same order, built from the single \texttt{tl\_*} definitions in \texttt{Chart/Gens.lean}; the certificate is instantiated at $\rho = \mathrm{Fin}\,3199$, $p = 32003$, $w_0 = 11147$ with A1 (\texttt{Cinv} $\cdot$ \texttt{MS} $= 1$), A2, and the six data hypotheses (hw$_0$, hscale, prime32003, hmon, hr$_1$, hsupp2 in Fpoly2 form, transported along the bridge). \texttt{\#print axioms challenge\_target} $=$ \texttt{[propext, Classical.choice, Quot.sound]}, and \texttt{\#print axioms branch\_c\_eliminated} is the same permitted triple (logs \texttt{bridge\_challenge\_target\_2026-11-04.log}, \texttt{branch\_c\_eliminated\_2026-11-04.log}, November 4, 2026, lean-build-2). \emph{Scope:} this is one case --- case~(1) --- of one proposition; it does not by itself exclude $(72,108)$ unconditionally and does not prove the planar Jacobian conjecture, both of which additionally require GGHV Proposition~4.3 to be correct and exhaustive and the reduction chain, ledgered here, not re-derived. \emph{Provenance:} proved as scratch files (\texttt{bridge\_attempt\_2026\_11\_04.lean}, \texttt{composition\_2026-11-04.lean}); incorporation as tracked repo modules is open item~(5) in \S\ref{bc:sec:epistemic}. \qed

\subsection{Logical status}

\begin{remark}\label{bc:rem:conditional}
Theorem~\ref{bc:thm:main} does not depend on GGHV: its Lean proof goes from \texttt{NewtonNFc} and $[P,Q]=\lambda x^2$ through the descent to the chart statement, all inside Lean. In Lean, \texttt{main\_theorem\_c\_of\_chartEmpty\_T1ne0} (listing~\ref{bc:lst:main}) derives the conclusion from the explicit hypothesis \texttt{ChartEmptyC\_T1ne0} (listing~\ref{bc:lst:chartempty}), with \texttt{\#print axioms} reporting only \texttt{propext}, \texttt{Classical.choice} and \texttt{Quot.sound}, and no \texttt{sorry}, \texttt{axiom} or \texttt{native\_decide} anywhere in the $790$ modules. The chain is:
\texttt{NewtonNFc} $+\ [P,Q]=\lambda x^2 \xrightarrow{\text{B0--B3}}$ \texttt{DescentClaimC} hypotheses $\xrightarrow{\text{B4--B5}}$ \texttt{ChartEmptyC} $\xrightarrow{\text{B6--B7}}$ \texttt{ChartEmptyC\_T1ne0} $\Rightarrow$ contradiction. Theorem~\ref{bc:thm:unconditional} (\texttt{challenge\_target}) now discharges \texttt{ChartEmptyC\_T1ne0} with no hypotheses, closing the chain unconditionally in Lean.
Only Corollary~\ref{bc:cor:gghv} additionally assumes GGHV Proposition~4.3.
\end{remark}

\begin{remark}[A retraction]\label{bc:rem:retraction}
A preliminary version of this work claimed a ``reduction lemma'': that a unit ideal modulo a prime lifts to a unit ideal over $\mathbb{Q}$. That statement is \textbf{false}. Counterexample: $F=1+109X$ is $1$ modulo~$109$ (unit ideal) but vanishes at $X=-1/109$ over $\mathbb{Q}$ (proper ideal). The Lean file deriving \texttt{False} from the stated axiom is archived for the record, and every theorem depending on it is vacuous. The correct replacement is the rank lemma (Lemma~\ref{bc:lem:rank}): \emph{full row rank} of a fixed-degree Macaulay matrix lifts through a nonzero minor, because nonvanishing of a determinant is an open condition. Mere consistency of a modular system does not lift. \S\ref{bc:sec:epistemic} gives the precise distinction.
\end{remark}

\subsection{Outline}

Section~\ref{bc:sec:source} source-gates all base objects against GGHV v1. Section~\ref{bc:sec:system} sets up the polynomial system and the grading. Section~\ref{bc:sec:layers} gives the exact layer reduction and the $132$ raw equations. Section~\ref{bc:sec:normalization} fixes normalizations and reaches the chart. Section~\ref{bc:sec:top} treats the top layer, inherited from branches~(a,b). Section~\ref{bc:sec:descent} carries out the descent to the twelve conditions, proved by kernel reflection. Section~\ref{bc:sec:conditions} studies the chart conditions. Section~\ref{bc:sec:t1zero} eliminates the stratum $b_{11,20}=0$ inside Lean. Section~\ref{bc:sec:wholechart} eliminates the whole chart by the rank lemma (now corroborating the Lean certificate). Section~\ref{bc:sec:proof} assembles the proof of Theorem~\ref{bc:thm:main} and records the unconditional Lean certificate (Theorem~\ref{bc:thm:unconditional}). Section~\ref{bc:sec:epistemic} gives the epistemic accounting: grades, the retraction, and what is not claimed. Section~\ref{bc:sec:conclusion} concludes. Section~\ref{bc:sec:attribution} covers prior work, attribution, and independence. Figure~\ref{bc:fig:bridges} shows the bridge architecture.

\begin{figure}[htbp]
\centering
\includegraphics[width=0.97\textwidth]{figures_c/fig2_bridges.pdf}
\caption{Bridge architecture. Top band: the eight Lean bridges B0--B7 (grade~A), each labelled with its certified Lean name. Middle band: the external rank lemma B8 (grade~B), formerly the only non-Lean component and now corroborating the Lean certificate of Theorem~\ref{bc:thm:unconditional}. Bottom band: the Lean propositions they target, composing left to right to the final contradiction. Grade colors are consistent with Table~\ref{bc:tab:bridges}. Rendered by \texttt{figures/fig2\_bridges.py}.}
\label{bc:fig:bridges}
\end{figure}

\section{Source data}\label{bc:sec:source}

Every base object below is source-gated directly against GGHV v1 before any proof is drafted, following the discipline of~\cite{BranchABv19}. The branch-(c) Newton polygons~\eqref{bc:eq:polygons} are \cite[Proposition~4.3, case~(1)]{GGHV}: the vertex lists are asserted at GGHV v1 line~494 (case~(1) list) and lines~617--618 (post-$\varphi$ case~c list), and the checker \filename{source\_gate\_verify\_c.py} re-asserts them byte-for-byte against the downloaded v1 text. The post-$\varphi$ bracket target $[P,Q]=\lambda x^2$ is \cite[\S5]{GGHV}.

\begin{leanlisting}{The branch-(c) Newton normal form (\texttt{NewtonNFc})}{bc:lst:newtonnfc}
-- Jacobian/BranchC/DescentClaim.lean (excerpt, exact statement)
def NewtonNFc (P Q : MvPolynomial (Fin 2) L) : Prop :=
  (P.support \subseteq NPc) \land (Q.support \subseteq NQc) \land
  TenVertexCoeffsNonzero P Q
-- NPc = conv {(0,0),(1,0),(8,14),(8,16),(0,8)}
-- NQc = conv {(0,0),(2,1),(12,21),(12,24),(0,12)}
\end{leanlisting}

The Lean formalization works over an arbitrary field $L$ of characteristic~$0$. The quintic field
\begin{equation}\label{bc:eq:K5}
K_5=\mathbb{Q}[w]/(R),\qquad R(w)=w^5-w^4+3w^3+3w^2+26,
\end{equation}
enters through the top layer (\S\ref{bc:sec:top}): $R$ is irreducible over $\mathbb{Q}$, and it remains irreducible modulo~$109$, so $\mathbb{F}_{109^5}=\mathbb{F}_{109}[w]/(R)$ is a field. The good prime/root pairs used in this audit are
\begin{equation}\label{bc:eq:primes}
(101,9),\qquad (1000003,806739),\qquad (32003,11147),
\end{equation}
each satisfying $R(w_0)\equiv0\pmod p$, plus the inert prime $109$.

\section{The polynomial system}\label{bc:sec:system}

\subsection{Lattice points}

Write $P=\sum a_{i,j}x^iy^j$, $Q=\sum b_{i,j}x^iy^j$. The branch-(c) supports are the lattice points of the polygons~\eqref{bc:eq:polygons}; the supports strictly contain those of branches~(a,b) because of the vertices $(0,8)$ and $(0,12)$. The Lean files record the supports as \texttt{rangeP}, \texttt{rangeQ} (\filename{Jacobian/BranchC/DescentClaim.lean}).

\subsection{The bracket equations}

Grade by $w(x)=2$, $w(y)=-1$. Write the layers
\[
P=\sum_k y^{-k}A_k,\qquad Q=\sum_l y^{-l}B_l,\qquad A_k,B_l\in L[u]
\]
with $u$ the residual variable. The bracket of weighted-homogeneous pieces is weighted-homogeneous, so $[P,Q]=\lambda x^2$ splits into one identity per weight:
\begin{equation}\label{bc:eq:eident}
\texttt{EIdent}\ A\ B\ \lambda\ n:\quad \sum_{k+l=n}\mathrm{LT}_{k,l}(A_k,B_l)=[n=5]\cdot\lambda u^2,\qquad \mathrm{LT}_{k,l}(A,B)=kAB'-lA'B.
\end{equation}
Layer $k$ of $P$ corresponds to $W=j-2i$ via $n=1-W$; the raw equation key $(i,j)$ corresponds to $(n=2i+1-j,\ m=i)$, and the coefficient of the product term is $c=k\cdot(m+1-i_1)-l\cdot i_1$ with $l=n-k$ (\S5.3 of the v3.1 guide). The top layer $n=5$ carries the $\lambda u^2$ term; all other layers are homogeneous.

\section{Exact layer reduction}\label{bc:sec:layers}

\textbf{B0 (Lean, v1).} From \texttt{jac}$(P,Q)=C\,\lambda\cdot x^2$ and \texttt{NewtonNFc}, the layer identities $\texttt{EIdent}\ A\ B\ \lambda\ n$ hold for every layer $n\ge-20$ (\texttt{layers\_of\_jac\_c}, \texttt{layer\_bracket\_gen} in \filename{LayersGen.lean}). This is an exact symbolic identity, verified by the Lean kernel.

\textbf{B4 (Lean, v3).} Coefficient extraction turns the polynomial identities $\texttt{EIdent}\ A\ B\ 0\ n$, $n=4,\dots,-2$, into scalar equations. For each monomial of $\mathrm{LT}_{k,l}(A_k,B_l)$, \texttt{coeff\_layerTermZ} computes its coefficient with integer indices; coefficients outside the branch-(c) supports vanish by \texttt{rangeP}, \texttt{rangeQ} ($179$ zero facts; $a_{0,0},b_{0,0}$ are in the support but always carry coefficient~$0$). Each of the $132$ raw equations then follows by \texttt{linear\_combination}. This is the recipe of the repository's \texttt{BranchAbFinal.lean}, replayed by \texttt{certgen\_c/gen\_bridge\_c.py} (\texttt{Descent2R/Bridge.lean}).

\section{Normalization}\label{bc:sec:normalization}

\textbf{B1 (Lean, v1).} The top layer $n=5$ is the branch-(a,b) problem, classified exactly over $K_5$ in~\cite{BranchABv19} and formalized in the repository (\texttt{ChartProof/Final.lean}, \texttt{BranchAb*}). Branch~(c) inherits its exact $K_5$ solution; this is where $w$ and $R$ enter (\texttt{eIdent\_five}).

\textbf{B2 (Lean, v1/v2).} Degree-$19$ rigidity: $b_{12,24}\ne0$ forces $t_2=b_{12,22}\ne0$. The theorems are \texttt{vertex\_12\_24\_forces\_t2} and \texttt{edge19\_rigidity\_of\_jac} in \filename{Edge19.lean}, and \texttt{t0\_vertex\_rigidity\_of\_jac} in \filename{LayerE2.lean}, \filename{Degree19.lean}).

\textbf{B3 (Lean, v1).} Torus and weighted scaling remove the continuous symmetries, normalizing to the rescaled $K_5$ top layer and the chart $b_{12,22}=1$ (\filename{eIdent\_transport}, \filename{main\_theorem\_c\_of\_claim} in \filename{DescentClaim.lean}). These are \texttt{DescentClaimC}'s hypotheses: the $19$ top-layer values, the $123$ unknowns, and the $132$ raw equations.

\section{The top layer}\label{bc:sec:top}

At $n=5$ the branch-(c) layer equations coincide with the branch-(a,b) top layer, whose complete $K_5$ classification is proved in~\cite{BranchABv19} (five torus orbits; the degree-$35$ eliminant irreducible over $\mathbb{Q}$). The classification is used here as a black box: the $19$ top-layer values enter \texttt{Descent2R/Main.lean} as hypotheses \texttt{h\_v} at minimal integer scale, and the $K_5$ arithmetic they satisfy is exactly the arithmetic of~\eqref{bc:eq:K5}. No new top-layer analysis is needed for branch~(c); the new mathematics is the descent of \S\ref{bc:sec:descent} and the chart analysis of \S\S\ref{bc:sec:conditions}--\ref{bc:sec:wholechart}.

\paragraph{Formalization seam.} Julian's 2026-09-26 audit identified one genuine seam at this step: the degree-$19$ rigidity lemma is \emph{not} formalized inside Lean. The parity-constrained Davenport--Stothers classification (the step ruling out $\deg f > 19$) is supplied by the external certified point in \texttt{e5\_cert.json} (exact equality over $K_5$), which Lean imports and checks by \texttt{decide}; the mathematical classification argument itself is grade~B, not grade~A. This is the only non-kernel step below the chart (the other gap flagged by the audit, the missing GGHV citation line, is bibliographic and is repaired in~\S\ref{bc:sec:attribution}).

\section{Descent}\label{bc:sec:descent}

\textbf{B5 (Lean, v3; kernel reflection).} Each layer $n=4,\dots,-2$ is a linear system in its new unknowns, with coefficients polynomial in the earlier unknowns and the chart parameters. The v2 operator-rank theorems (\filename{Rank/**}, 96 modules) give each system's exact rank over $K_5$; the left null vectors give the elimination conditions:
{\small
\begin{center}
\begin{tabular}{lll}
\textbf{Layer} & \textbf{Modules} & \textbf{Conditions produced} \\ \midrule
$E_4$ ($n=4$) & \texttt{Descent2R/L4} (54) & --- \\
$E_3$ ($n=3$) & \texttt{L3} (57) & --- \\
$E_2$ ($n=2$) & \texttt{L2c} (57) & none (left null vector annihilates RHS) \\
$E_1$ ($n=1$) & \texttt{L1} (60) & $\Omega$ \\
$E_0$ ($n=0$) & \texttt{L0} (60) & $\Psi$ and pure row \texttt{E0pure\_18\_37} \\
$E_{-1}$ ($n=-1$) & \texttt{Lm1} (99) & $\Phi_1,\Phi_2$ and pure rows \texttt{Em1pure\_17\_36}, \\
& & \texttt{Em1pure\_18\_38} \\
$E_{-2}$ ($n=-2$) & \texttt{Lm2} (110) & $\Theta_1,\Theta_2,\Theta_3$ and pure rows \\
& & \texttt{Em2pure\_16\_35}, \texttt{Em2pure\_17\_37} \\
\end{tabular}
\end{center}
}
\texttt{lower\_c.py} computed these eliminations in Python; \texttt{Descent2R} replays them as $600$ polynomial identities $\mathrm{target}=\sum m_i\cdot\mathrm{fact}_i$, each checked by the Lean kernel. Reflection uses the repository's \filename{Jacobian/ChartProof/Reflect.lean} (built on core's \texttt{Lean.Grind.CommRing}): \texttt{Expr}/\texttt{Poly}, normalized by \texttt{toPolyK} and compared by \texttt{decide +kernel}. Each step theorem has the form $(\texttt{Expr.mul}\ (\texttt{.num}\ K)\ X)\texttt{.denote}\ \mathrm{ctx}=0$ proved by \texttt{lc\_zero} with a kernel \texttt{decide}; \texttt{Main} removes $K$ with \texttt{cancel\_int} when $K\ne1$ ($231$ of $752$ \texttt{have}s). The context is $143$ atoms in a balanced \texttt{RArray} (\texttt{gctx2}). A tactic proof of the same identities was tried first and abandoned: $214$~MB of source, chunks over $10$~min each. The kernel evaluates the normalizer directly with GMP arithmetic; \texttt{decide +kernel} is not \texttt{native\_decide}---no compiler is trusted.

The main theorem of the descent (\texttt{Descent2R/Main.lean}):
\begin{leanlisting}{\texttt{chart\_descent\_refl}: the twelve conditions from the $132$ raw equations}{bc:lst:descent}
-- Descent2R/Main.lean (statement; 752 haves elided)
theorem chart_descent_refl
  (w : L) (hw : R w = 0)
  (v : Fin 19 \to L) (h_v : v = K5_top_layer_values)
  (u : Fin 123 \to L)
  (h4_ : raw_equations_n4) ... (hm2_ : raw_equations_nm2) :
  cond_Omega w ... = 0 \land cond_Psi w ... = 0 \land
  cond_Phi1 w ... = 0 \land cond_Phi2 w ... = 0 \land
  cond_Theta1 w ... = 0 \land cond_Theta2 w ... = 0 \land cond_Theta3 w ... = 0 \land
  cond_E0pure_18_37 w ... = 0 \land cond_Em1pure_17_36 w ... = 0 \land
  cond_Em1pure_18_38 w ... = 0 \land cond_Em2pure_16_35 w ... = 0 \land
  cond_Em2pure_17_37 w ... = 0
\end{leanlisting}
The twelve conjuncts are \texttt{lower\_c}'s conditions, defined once in \texttt{CondsC.lean} (generated by \texttt{certgen\_c/}\linebreak\texttt{gen\_conds\_c.py}): each \texttt{cond\_*} is the explicit polynomial with the same balanced sum tree as the reflected fact (subtrees of at most $32$ terms are separate definitions: $110$ definitions for $12$ conditions), so \texttt{ek.denote ctx = 0} and \texttt{cond\_* \dots\ = 0} are definitionally equal. \texttt{Bridge} (\texttt{Descent2R/Bridge.lean}) then gives \texttt{descentClaimC\_of\_chartEmpty}: \texttt{DescentClaimC} $\Leftarrow$ \texttt{ChartEmptyC}.

\section{The chart conditions}\label{bc:sec:conditions}

\subsection{Chart variables}

On the normalized chart $b_{12,22}=1$ the conditions involve seven parameters:

\begin{center}
\begin{tabular}{lccccccc}
\toprule
Symbol & $T_1$ & $T_2$ & $S_1$ & $S_2$ & $R_1$ & $R_2$ & $Q$ \\
Coefficient & $b_{11,20}$ & $b_{12,22}$ & $b_{11,21}$ & $b_{12,23}$ & $a_{6,13}$ & $a_{7,15}$ & $b_{10,21}$ \\
Weight & $1$ & $1$ & $2$ & $2$ & $3$ & $3$ & $4$ \\
\bottomrule
\end{tabular}
\end{center}

\subsection{The twelve conditions}\label{bc:sec:twelve}

\texttt{certgen\_c/conds\_c.json} (md5 \texttt{168299d255af361578423c0ed4852c11}, written by \texttt{certgen\_c/lower\_c.py}):

\begin{center}
\begin{tabular}{lcc}
\toprule
Condition & $K_5$-terms & Integer terms (minimal scale, $w$ expanded) \\
\midrule
$\Omega$ & $3$ & $15$ \\
$\Psi$ & $32$ & $160$ \\
$\Phi_1,\ \Phi_2$ & $54,\ 54$ & $270,\ 270$ \\
$\Theta_1,\ \Theta_2,\ \Theta_3$ & $84,\ 84,\ 84$ & $420,\ 420,\ 420$ \\
\texttt{E0pure\_18\_37} & $9$ & $41$ \\
\texttt{Em1pure\_17\_36} & $39$ & $195$ \\
\texttt{Em1pure\_18\_38} & $12$ & $60$ \\
\texttt{Em2pure\_16\_35} & $77$ & $381$ \\
\texttt{Em2pure\_17\_37} & $55$ & $275$ \\
\bottomrule
\end{tabular}
\end{center}

\subsection{\texorpdfstring{The $\Omega$ parabola}{The Omega parabola}}\label{bc:sec:omega}

Exactly over $K_5$,
\begin{equation}\label{bc:eq:omega}
\Omega=o_1S_2^2+o_2T_2^2S_2+o_3T_2^4,\qquad o_2^2=4o_1o_3,\quad o_1\ne0,
\end{equation}
so that with $\kappa=-o_2/(2o_1)\in K_5$,
\begin{equation}\label{bc:eq:omegafactor}
\Omega=o_1(S_2-\kappa T_2^2)^2.
\end{equation}
The relation $o_2^2=4o_1o_3$ is checked in \texttt{step1/lowerc\_chart.py} and independently in \texttt{step3c\_chart\_identity.py}; $o_1(w)\ne0$ is proved in the Lean kernel by a Bezout identity $a\cdot N_{o_1}+b\cdot R=c$ with $c$ a $151$-digit constant (\texttt{T1Zero.s\_X2}). On the chart $T_2=1$, $\Omega=0$ forces $S_2=\kappa$: the chart collapses to the parabola's vertex line (Figure~\ref{bc:fig:omega}).

\begin{figure}[htbp]
\centering
\includegraphics[width=0.95\textwidth]{figures_c/fig3_omega.pdf}
\caption{Chart geometry near $\Omega=0$. Main panel: the $(T_2,S_2)$-plane drawn schematically over $\mathbb{R}$ for illustration (the mathematics is over $K_5$), showing the parabola $S_2=\kappa T_2^2$ from~\eqref{bc:eq:omegafactor}, the chart line $T_2=1$ (normalization B3), and their intersection $(1,\kappa)$; the stratum $T_1=b_{11,20}=0$ is indicated as a transverse slice, with the $T_1=0$ stratum (Lean, grade~A) and the $T_1\ne0$ stratum (rank lemma, grade~B; now also Lean certificate, grade~A) labelled. Inset: the six chart generators' $K_5$-term counts ($22,35,35,52,52,52$), the input to the Macaulay matrix of \S\ref{bc:sec:wholechart}. Rendered by \texttt{figures/fig3\_omega.py}.}
\label{bc:fig:omega}
\end{figure}

\section{\texorpdfstring{The stratum $b_{11,20}=0$}{The stratum b11,20=0}}\label{bc:sec:t1zero}

\textbf{B7 (Lean kernel).} On the stratum $T_1=b_{11,20}=0$ the conditions admit an exact $K_5$ certificate small enough for the kernel: degree~$2$, $75$ $K_5$-terms, coefficient heights at most $9{,}681$ digits (\texttt{chart\_certificates/step1/certB\_lowerc.json}, produced by \texttt{certB\_lowerc.py} via pivots from a mod-$p$ RREF and an exact FLINT solve; independently rechecked by \texttt{check\_certB\_lowerc.py} with \texttt{fmpq\_mpoly} expansion, $w$ as a variable, reduced mod~$R$ at the end). \texttt{T1Zero} checks the certificate chain end to end in the kernel (\texttt{gen\_t1zero\_lean.py}):
{\footnotesize
\begin{center}
\begin{tabular}{lll}
Step & Identity (integer-polynomial, kernel-checked) & Build \\ \midrule
\texttt{s\_sq} & $D_O D_{o_1} D_k^2\Omega\equiv D_O N_{o_1}X^2\bmod(T_2-1,R)$, $X=D_kS_2-N_\kappa(w)$, $D_k=851565312$ & $1.5$~s \\
\texttt{s\_X2} & Bezout $a\cdot e_{\mathrm{sq}}+b\cdot D_OX^2R=c\cdot D_OX^2$ ($c$: $151$ digits) $\Rightarrow X=0$ & --- \\
\texttt{s\_f\_n} (6) & $D_k^d\cdot ck_n=D_k^dT_1A_1+D_k^d(T_2-1)A_2+XA_3+f_n+RA_4$ & $2.4$--$5.7$~s \\
\texttt{s\_p\_n} (6) & $p_n=\nu_n f_n-q_nR$ ($250$--$375$ terms, $\approx5{,}440$ digits) & $17$--$28$~s \\
\texttt{s\_final} & $\sum_n p_n=DD$, a $5387$-digit integer, exact & $55.6$~s \\
\end{tabular}
\end{center}
}
\texttt{chartEmpty\_t1\_zero} chains these: $(DD:L)=0$ contradicts characteristic~$0$ ($26$~s under \texttt{lake}). The hypotheses are exactly the \texttt{cond\_*} of \texttt{CondsC}, definitionally the conditions \texttt{Descent2R} derives. Five perturbation controls (Bezout constant $+1$, a dropped equation, a perturbed cofactor coefficient, $DD+1$, a dropped $T_2$ equation) are all rejected by the kernel (\texttt{logs/t1z\_controls.log}: five of five ``Tactic \texttt{decide} failed'', as required).

\textbf{B6 (Lean).} \texttt{chartEmptyC\_of\_T1ne0} (\texttt{T1Zero/Combine.lean}) splits on $b_{11,20}$: \texttt{ChartEmptyC} $\Leftarrow$ \texttt{ChartEmptyC\_T1ne0}, isolating exactly what Lean does not check as a precise \texttt{Prop}:
\begin{leanlisting}{\texttt{ChartEmptyC\_T1ne0}: what Lean labels; formerly proved in \S\ref{bc:sec:wholechart} (grade~B), now proved in Lean by \texttt{challenge\_target} (grade~A, Theorem~\ref{bc:thm:unconditional})}{bc:lst:chartempty}
-- T1Zero/Combine.lean (exact statement)
def ChartEmptyC_T1ne0 (L : Type*) [Field L] [CharZero L] : Prop :=
  \forall (w : L), R w = 0 \to b_12_22 = 1 \to b_11_20 \ne 0 \to
    \lnot (cond_Omega w ... = 0 \land ... \land cond_Theta3 w ... = 0)
\end{leanlisting}

\section{The whole chart}\label{bc:sec:wholechart}

A canonical $K_5$ certificate for the whole chart is too large to build on this machine: at weight $W=24$ the Hadamard bound is $4{,}392{,}160$ digits, and calibrating against the $t_1=0$ slices ($W=2,3,4$, where the true-to-bound height ratios are $0.061,0.061,0.046$) gives heights of $2.0$--$2.7\times10^5$ digits, about $2\times10^9$ digits in all (grade~C estimate; the Hadamard bound itself is rigorous). Exact solution runs out of memory at $7$~GB; multimodular reconstruction would need $\approx3\times10^4$ primes at $15$--$20$~s each, about six days. Instead we prove the \emph{existence} of the certificate.

\subsection{The rank lemma}

\begin{lemma}[Fixed-degree Macaulay rank lifts]\label{bc:lem:rank}
Let $A=\mathbb{Z}_{(p)}[w]/(R)$, free of rank~$5$ over $\mathbb{Z}_{(p)}$, embedding in $K_5$. Let $\varphi:A\to\mathbb{F}_p$, $w\mapsto w_0$, with $R(w_0)\equiv0\pmod p$. Let $M$ be the weighted Macaulay matrix of the chart generators at fixed weight $W$ (columns $(n,m)$ with $\mathrm{wdeg}\,m\le W-\mathrm{wt}\,F_n$; rows the monomials of the products $m\cdot F_n$, plus~$1$; entries the coefficients). If $\varphi(M)$ has full row rank $N$, then $M$ has full row rank over $K_5$.
\end{lemma}
\begin{proof}
Some $N\times N$ minor satisfies $\det\varphi(M_S)=\varphi(\det M_S)\ne0$. Hence $\det M_S\ne0$ in $A\subset K_5$, and $M_S$ is invertible over the field $K_5$. \qed
\end{proof}

Consequences. $Mx=e_1$ is solvable over $K_5$, so $1=\sum c_nF_n$ in $K_5[t,s,r,u,q]$ with $\mathrm{wdeg}\,c_n\le24-\mathrm{wt}\,F_n$. Mapping $K_5\to L$ by $w\mapsto w$, no $L$ admits a common zero. With $S_2=\kappa$ from~\eqref{bc:eq:omegafactor} ($o_1\ne0$), no chart point satisfies $\Omega=\Psi=\dots=\Theta_3=0$: that is \texttt{ChartEmptyC}, hence \texttt{ChartEmptyC\_T1ne0}, even without the pure rows.

\begin{remark}[Why this lemma succeeds where the retracted one failed]
A unit ideal mod~$p$ says only that the Macaulay system is \emph{consistent} mod~$p$; consistency does not lift ($1+109X$ over $\mathbb{Z}_{(109)}$). Full row rank at fixed degree is the nonvanishing of a minor, an \emph{open} condition, and open conditions lift. The planted-zero control below shows the sensitivity: one common zero costs exactly one rank.
\end{remark}

\subsection{The instance}

\texttt{step3b\_rank\_lift.py} ($6.5$~min, $0.5$~GB; \texttt{step3b\_rank\_lift.log}):

{\small
\begin{center}
\begin{tabular}{lcc}
\toprule
 & $p=1000003$ & $p=32003$ \\
 & $w_0=806739$ & $w_0=11147$ \\
\midrule
All coefficients $p$-integral, $R(w_0)\equiv0$ & asserted & asserted \\
$W=22$: rank / augmented / rows & $2277/2278/2281$: inconsistent & same \\
$W=23$: rank / augmented / rows & $2707/2708/2708$: inconsistent & same \\
$\mathbf{W=24}$: $\mathbf{3199\times6054}$ & $\mathbf{3199/3199/3199}$: \textbf{full row rank} & \textbf{same} \\
NumPy pivot-block elimination (independent) & $\det\ne0$ & $\det\ne0$ \\
Planted common zero, $W=24$ (control) & rank $3198$, inconsistent & same \\
\bottomrule
\end{tabular}
\end{center}
}

The protocol was frozen before the run: the lemma, the second prime, the independent NumPy elimination and the controls were all fixed in \texttt{step3b\_rank\_lift.py} before it executed. The $W=24$ threshold is sharp: $W=22$ and $W=23$ are inconsistent (Figure~\ref{bc:fig:ranks}).

\begin{figure}[htbp]
\centering
\includegraphics[width=0.95\textwidth]{figures_c/fig4_ranks.pdf}
\caption{Macaulay rank profile at weights $W=22,23,24$, identical at both primes ($p=1000003$ and $p=32003$). For each weight: bars for the number of rows ($2281$, $2708$, $3199$), the computed rank ($2277$, $2707$, $3199$), and the augmented rank, with the $W=24$ group highlighted (full row rank $3199/3199$, the obstruction). The fourth group shows the planted-common-zero control at $W=24$ (rank $3198$, exactly one short of full). Data source: \texttt{step3b\_rank\_lift.log}. Rendered by \texttt{figures/fig4\_ranks.py}; the $y$-axis is truncated with a break mark to make the single-rank deficit visible.}
\label{bc:fig:ranks}
\end{figure}

The chart generators are independently verified as the right polynomials (exact evaluation at $12$ random points; an independent $\kappa$; the exact $\Omega$ square; controls including $\kappa+1$ detected) by the script \texttt{step3c\_chart\_identity.py}. This is \textbf{B8}, proved outside Lean at grade~B. B8 remains valid corroboration; it is no longer the sole route to \texttt{ChartEmptyC\_T1ne0} (see Theorem~\ref{bc:thm:unconditional}).

\section{Proof of the main theorem}\label{bc:sec:proof}

\begin{leanlisting}{\texttt{main\_theorem\_c\_of\_chartEmpty\_T1ne0}}{bc:lst:main}
-- Jacobian/BranchC/Combine (exact statement)
theorem main_theorem_c_of_chartEmpty_T1ne0
    {L : Type*} [Field L] [CharZero L]
    (h : ChartEmptyC_T1ne0 L) :
    \lnot \exists (P Q : MvPolynomial (Fin 2) L) (lam : L),
      lam \ne 0 \land NewtonNFc P Q \land jac P Q = C lam * X 0 ^ 2
\end{leanlisting}

\emph{Proof.} Let $P,Q,\lambda$ satisfy \texttt{NewtonNFc} and $[P,Q]=\lambda x^2$, $\lambda\ne0$. B0--B3 give \texttt{DescentClaimC}'s hypotheses (the $19$ top-layer values, $123$ unknowns, $132$ raw equations). B4 extracts the raw equations; B5 (\texttt{chart\_descent\_refl}) yields the twelve conditions, i.e.\ \texttt{ChartEmptyC}. B6 splits on $b_{11,20}$: on $b_{11,20}=0$, B7 (\texttt{chartEmpty\_t1\_zero}) gives \texttt{False} in the kernel; on $b_{11,20}\ne0$, the hypothesis \texttt{ChartEmptyC\_T1ne0} gives \texttt{False}. Hence no such $P,Q,\lambda$ exist. \qed

Corollary~\ref{bc:cor:unconditional} follows by supplying \texttt{ChartEmptyC\_T1ne0} from \texttt{challenge\_target} (Theorem~\ref{bc:thm:unconditional}); the \S\ref{bc:sec:wholechart} rank route supplies it independently at grade~B. Corollary~\ref{bc:cor:gghv} by applying \cite[Proposition~4.3]{GGHV}.

The bridge composition is summarized in Table~\ref{bc:tab:bridges}.

{\small
\begin{center}
\begin{tabular}{lllll}
\toprule
\textbf{Bridge} & \textbf{From} & \textbf{To} & \textbf{Lean name} & \textbf{Grade} \\
\midrule
B0 & $[P,Q]=\lambda x^2$, \texttt{NewtonNFc} & layer identities $E_n$ & \texttt{layers\_of\_jac\_c} & A \\
B1 & $E_5$, vertices & $K_5$ top layer & \texttt{eIdent\_five} & A \\
B2 & B0, B1 & $t_2=b_{12,22}\ne0$ & \texttt{vertex\_12\_24\_forces\_t2} & A \\
B3 & B0--B2 & chart $b_{12,22}=1$ & \texttt{main\_theorem\_c\_of\_claim} & A \\
B4 & $E_n$, $n=4\dots-2$ & $132$ raw equations & \texttt{eIdent\_c4}$\dots$\texttt{eIdent\_cm2} & A \\
B5 & $132$ equations & twelve conditions & \texttt{chart\_descent\_refl} & A \\
B6 & \texttt{ChartEmptyC\_T1ne0} & \texttt{ChartEmptyC} & \texttt{chartEmptyC\_of\_T1ne0} & A \\
B7 & conditions, $b_{11,20}=0$ & \texttt{False} & \texttt{chartEmpty\_t1\_zero} & A \\
B8 & chart generators & \texttt{False} & rank lemma + \texttt{step3b\_rank} & B \\
& & & \texttt{\_lift.py} & \\
\bottomrule
\end{tabular}
\end{center}
}
\label{bc:tab:bridges}

B8 is now corroborating rather than load-bearing: the Lean certificate \texttt{challenge\_target} discharges \texttt{ChartEmptyC\_T1ne0} at grade~A (Theorem~\ref{bc:thm:unconditional}); the table grades for B0--B7 and B8 as originally proved are unchanged.

\section{Epistemic status and what is not claimed}\label{bc:sec:epistemic}

HLRE bookkeeping: grades A (machine-checked or exact proof from stated axioms), B (derived under explicit premises), C (heuristic or fitted), D (tautological or refuted).

{\small
\begin{center}
\begin{tabular}{llp{5.2cm}c}
\toprule
\textbf{ID} & \textbf{Claim} & \textbf{Evidence} & \textbf{Grade} \\
\midrule
C1 & \texttt{main\_theorem\_c\_of\_claim} & Lean kernel (v1) & A \\
C2 & \texttt{descentClaimC\_of\_chartEmpty} & Lean kernel, reflective; $600$ identities pre-checked in Python & A \\
C3 & \texttt{chartEmpty\_t1\_zero} & Lean kernel; FLINT certificate rechecked in Python; $5$ controls & A \\
C4 & \texttt{chartEmptyC\_of\_T1ne0} & Lean & A \\
C5 & \texttt{ChartEmptyC} holds (all char.-$0$ $L$) & rank lemma + two-prime, two-implementation ranks + input checks & \textbf{B} \\
C6 & canonical certificate needs $\approx2\times10^9$ digits & Hadamard bound (rigorous) $\times$ calibrated ratio (fitted) & \textbf{C} \\
C7 & the stated reduction lemma is inconsistent & Lean counterexample ($1+109X$) & A \\
C8 & \texttt{lower\_c}'s ideal differs from the pipeline's & exact comparison & A \\
C9 & \texttt{ChartEmptyC\_T1ne0} holds (all char.-$0$ $L$, no hypotheses) & Lean certificate \texttt{challenge\_target} (bridge + A1, A2, hsupp2) & A \\
C10 & branch~(c) eliminated: no \texttt{NewtonNFc} pair with $[P,Q]=\lambda x^2$, $\lambda\ne0$ & \texttt{branch\_c\_eliminated} from C9 + \texttt{main\_theorem\_c\_of\_chartEmpty\_T1ne0} & A \\
\bottomrule
\end{tabular}
\end{center}
}
\label{bc:tab:grades}

Verbs follow HLRE control: ``proves'' only for Lean-checked or exact deductions; ``proves under stated computational premises'' for C5; ``estimates'' for C6.

What is \emph{not} claimed: not the Jacobian conjecture; not any other case of GGHV Proposition~4.3; not the reduction from the conjecture to Proposition~4.3; not $(72,108)$ excluded unconditionally (that needs GGHV correctness and exhaustiveness, ledgered not re-derived); not an explicit $K_5$ certificate for the whole chart (its existence at grade~B is C5; the load-bearing whole-chart step is now the Lean certificate C9, not an exhibited $K_5$ certificate); not a review of the deferred single-prime $p$-adic argument or of the retracted \texttt{char0\_cert} bundle. A Lean proof of the stratum $b_{11,20}\ne0$ now exists (C9) and is therefore removed from this list. The mod-$p$ Gr\"obner/lift certificates ($\mathbb{F}_{101}$, $\mathbb{F}_{1000003}$, $\mathbb{F}_{109^5}$) are certificates, not proofs: they exclude only solutions integral at that prime. The rank argument is what lifts.

Open items, in order: (1)~I3 replication of C5---an independent rank computation of the $W=24$ Macaulay matrix in Magma, Sage or Singular from \texttt{conds\_c.json} and $\kappa$; (2)~optional formalization of the rank lemma's instance in Lean (\texttt{RingHom.map\_det}; the instance needs the $3199\times6054$ matrix as Lean data); (3)~deferred reviews; (4)~the v2 corrected-commit proposal for \texttt{branch\_c/}, on the author's explicit word only; (5)~incorporation of \texttt{challenge\_target} and \texttt{branch\_c\_eliminated} (currently 2026-11-04 scratch files, not tracked repo modules) as tracked modules.

\section{Conclusion}\label{bc:sec:conclusion}

Branch~(c), the last open configuration of GGHV Proposition~4.3 at degree pair $(72,108)$, is eliminated, unconditionally in Lean. The descent from $[P,Q]=\lambda x^2$ to twelve explicit chart conditions is proved twice: symbolically, as $600$ kernel-checked reflective identities in Lean~4 over $505$ generated modules, and numerically, by exact finite-field linear algebra. The chart is empty: in the Lean kernel on the stratum $b_{11,20}=0$, and on the whole chart since November 4, 2026 by the Lean certificate \texttt{challenge\_target} (Theorem~\ref{bc:thm:unconditional}); the rank lemma---full row rank of the $W=24$ Macaulay matrix at two primes, verified by two independent implementations---remains as independent corroboration. An earlier, false reduction lemma is retracted and replaced by the rank lemma; the paper's logical dependencies are those of Table~\ref{bc:tab:grades} as updated (C9, C10 grade~A). With branches~(a,b) eliminated in~\cite{BranchABv19}, no $(72,108)$ counterexample arises from GGHV Proposition~4.3, conditional on that proposition's correctness and exhaustiveness.

\section{Prior work, attribution, and independence}\label{bc:sec:attribution}

The GGHV classification~\cite{GGHV} is the sole external mathematical input, used only in Corollary~\ref{bc:cor:gghv}. The top-layer $K_5$ classification and the reflection library are inherited from the branch-(a,b) formalization~\cite{BranchABv19} without modification; \texttt{verify\_branch\_c\_lean.sh} step~0 checks these prerequisites by content fingerprint. The \texttt{branch\_c/} pipeline scripts were read-only throughout: the pipeline parametrizes the chart differently, and the v3 proof uses only \texttt{lower\_c}'s conditions, derived in Lean. The $W=24$ rank at $p=1000003$ was known from the step-3 measurement; the lemma, the second prime ($32003$), the independent NumPy elimination and the controls were frozen in \texttt{step3b\_rank\_lift.py} before it ran---the second prime agreed, as expected, since full row rank over $K_5$ implies full row rank at all but finitely many primes.

\part{\texorpdfstring{Dicriticals at Infinity of Keller Maps: Theorem R and Local Obstructions}{Dicriticals at Infinity of Keller Maps: Theorem R and Local Obstructions}}

\section{Introduction: setting, definitions, and the corrected statement}\label{td:sec:statement}

\subsection{Context}

The planar Jacobian conjecture asserts that a polynomial map
$(P,Q) \colon \mathbb{C}^2 \to \mathbb{C}^2$ with
$[P,Q] := P_x Q_y - P_y Q_x \in \mathbb{C}^\times$ is invertible. A \emph{Keller map} is
such a pair with constant nonzero Jacobian. This paper studies the dicritical divisors of the
pencil $\lambda P + \mu Q$ of a Keller map, and in particular how many of them can lie at
infinity. The October 1 draft claimed this line of work was closed. It was not, and this
revision says exactly what is proved, what is only sketched or computed at fixed degree, and
what is open.

Throughout, $K$ is an algebraically closed field of characteristic $0$; statements over any
field of characteristic $0$ reduce to its algebraic closure. We write
$F = (P,Q) \colon \mathbb{A}^2 \to \mathbb{A}^2$ for a Keller map, so
$[P,Q] = P_x Q_y - P_y Q_x \in K^\times$.

\subsection{Standard facts used}

We use the following standard facts, labelled (F1)--(F3) as in the restatement note:
\begin{itemize}[nosep]
\item[(F1)] $F$ is \'{e}tale. Its generic degree $e = [K(x,y) : K(P,Q)]$ is at least $1$.
\item[(F2)] The \emph{non-properness set} $S_F$ is the set of points of the target over which
$F$ is not proper. It is either empty or a curve whose components are rational
(Jelonek~\cite{Jelonek1993}, for dominant maps $\mathbb{A}^n \to \mathbb{A}^n$).
\item[(F3)] $F$ is an automorphism if and only if $S_F = \emptyset$. Indeed, if
$S_F = \emptyset$ then $F$ is a finite \'{e}tale cover of the simply connected
$\mathbb{A}^2$, hence an isomorphism; if $F$ is not an automorphism then $e \geq 2$, since
with $e = 1$ the map $F$ is birational and \'{e}tale, hence an open immersion by Zariski's
main theorem, so it is injective, and then bijective
(Bia\l{}ynicki-Birula--Rosenlicht~\cite{BBR1962}).
\end{itemize}

\subsection{Compactification, the pencil, and the standard definition}

Source and target are $\mathbb{A}^2 \subset \mathbb{P}^2$, with lines at infinity $L_\infty$
and $L'_\infty$. Let $\sigma \colon X \to \mathbb{P}^2$ be a sequence of point blow-ups over
$L_\infty$ such that $G = \bar{F} \circ \sigma \colon X \to \mathbb{P}^2$ is a morphism. The
\emph{curves at infinity} of $X$ are the components of $X \setminus \mathbb{A}^2$: the strict
transform of $L_\infty$ and the exceptional curves.

For $c \in \mathbb{A}^2$ (target), $\Lambda_c$ is the pencil whose members are $F^{-1}(\ell)$
for the lines $\ell \ni c$. Its members are the fibers of
$\varphi_c = \pi_c \circ F$, where $\pi_c$ is the projection from $c$. The paper's pencil
$\lambda P + \mu Q$ is $\Lambda_0$.

\begin{definition}[Dicritical divisor, standard definition]\label{td:def:dicritical}
A divisor $E$ over the source plane is \emph{dicritical} for $\Lambda_c$ if $\varphi_c$ is
not constant on $E$. Its \emph{degree} $m(E)$ is the degree of
$\varphi_c|_E \colon E \to \mathbb{P}^1$. The divisor $E$ is \emph{at infinity} if its center
on $\mathbb{P}^2$ lies on $L_\infty$; $L_\infty$ itself counts. Otherwise $E$ is
\emph{affine}. The \emph{depth} of $L_\infty$ is $0$; the exceptional curve of the blow-up at
a point $p$ has depth one more than the largest depth of the curves through $p$. Thus $E_1$
over a point of $L_\infty$ has depth $1$. A point of an exceptional curve is a
\emph{satellite} point if it lies on another component of the total transform of $L_\infty$,
and \emph{free} otherwise.
\end{definition}

\begin{remark}
The October 1 draft defined dicriticals as ``exceptional divisors on which the pulled-back
pencil has no base points.'' That is not the standard notion, and it is replaced by
Definition~\ref{td:def:dicritical} throughout this revision.
\end{remark}

\subsection{Counting dicriticals}

\begin{proposition}[The count]\label{td:prop:count}
\textit{Status: grade B, proof by hand from \textnormal{(F1)--(F3)}.}
Let $F$ be a Keller map and $c \in \mathbb{A}^2 \setminus S_F$.
\begin{enumerate}[label=(\alph*),nosep]
\item The affine dicriticals of $\Lambda_c$ are exactly the $e$ exceptional curves $E_q$ of
the blow-ups at the points $q \in F^{-1}(c)$. Each has degree $1$.
\item The dicriticals at infinity are the curves at infinity $D$ of $X$ for which $G(D)$ is a
curve. Then $G(D)$ is $L'_\infty$ or a component of $\bar{S}_F$, and
\begin{equation}\label{td:eq:count-degree}
m(D) = \deg\big(D \to G(D)\big) \cdot \deg G(D).
\end{equation}
\item $L'_\infty$ and every component of $\bar{S}_F$ are $G(D)$ for at least one such $D$.
So the number $N_\infty(c)$ of dicriticals at infinity satisfies
\begin{equation}\label{td:eq:count-lower}
N_\infty(c) \geq 1 + \#\{\text{components of } S_F\}.
\end{equation}
\end{enumerate}
\end{proposition}

\begin{proof}
We prove (a), (b), (c) in order, spelling out each step.

\emph{Part (a).} Since $c \notin S_F$, the map $F$ is proper over a neighbourhood of $c$.
By (F1), $F$ is \'{e}tale, so the fiber $F^{-1}(c)$ has exactly $e$ points, and $F$ is a local
isomorphism at each of them. Fix $q \in F^{-1}(c)$. Near $q$, the map $\varphi_c$ is the
projection from the point $c$ composed with a local isomorphism. One blow-up at $q$ resolves
it: in the blow-up chart, the exceptional curve $E_q$ parametrizes the directions through $q$,
and $\varphi_c$ restricts to the projection itself, hence maps $E_q$ isomorphically onto
$\mathbb{P}^1$. So $E_q$ is dicritical of degree $1$. At every other affine point,
$\varphi_c$ is regular, so no other affine divisor is dicritical. This gives exactly the $e$
exceptional curves $E_q$, each of degree $1$.

\emph{Part (b).} Let $D$ be a curve at infinity. First, $G(D)$ is not the point $c$: the map $G$
sends no curve at infinity into $\mathbb{A}^2 \setminus S_F$, and $c \notin S_F$ by hypothesis.
If $G(D)$ is a point different from $c$, then $\varphi_c = \pi_c \circ G$ is constant on $D$,
since $\pi_c$ sends every point other than $c$ to a single point of $\mathbb{P}^1$ only after
fixing the line through $c$ --- more precisely, $\pi_c$ is constant on no curve through a point
other than $c$ except by collapsing it, and here $G(D)$ is a single point $\neq c$, so
$\varphi_c|_D$ is that constant value. Hence $D$ is not dicritical in this case.

If instead $G(D)$ is a curve, then either $G(D) = L'_\infty$, or $G(D)$ meets
$\mathbb{A}^2$. In the second case, every point of $G(D) \cap \mathbb{A}^2$ is a limit of $F$
along a path to infinity, so it lies in $S_F$; since $G(D)$ is irreducible, $G(D)$ is a
component of the closure $\bar{S}_F$. Neither $L'_\infty$ nor a component of $S_F$ is a line
through $c$. Therefore the projection $\pi_c$ restricted to the curve $G(D)$ has degree
$\deg G(D)$, and composing with $D \to G(D)$ gives
\[
m(D) = \deg(\varphi_c|_D) = \deg\big(D \to G(D)\big) \cdot \deg G(D),
\]
which is \eqref{td:eq:count-degree}. So $D$ is dicritical exactly when $G(D)$ is a curve, with
the stated degree.

\emph{Part (c).} The map $G$ is proper and surjective, so $L'_\infty$ is dominated by some
curve at infinity: there is $D$ with $G(D) = L'_\infty$. By the description of $S_F$ as
$G(X \setminus \mathbb{A}^2) \cap \mathbb{A}^2$, each component of $S_F$ is likewise
dominated by a curve at infinity. Counting the distinct images $L'_\infty$ and the components
of $S_F$ gives the lower bound \eqref{td:eq:count-lower}.
\end{proof}

\begin{corollary}[Automorphisms]\label{td:cor:automorphisms}
\textit{Status: grade B.}
If $F$ is an automorphism, then $\Lambda_c$ has exactly two dicriticals for every $c$: one
affine, over $F^{-1}(c)$, and one at infinity, the unique curve $D$ with
$G(D) = L'_\infty$ (here $G$ is birational). Both have degree $1$.
\end{corollary}

\begin{proof}
If $F$ is an automorphism, then $e = 1$ and $S_F = \emptyset$ by (F3). Part (a) of
Proposition~\ref{td:prop:count} gives exactly one affine dicritical, over the single point
$F^{-1}(c)$, of degree $1$. Part (c) with $S_F = \emptyset$ gives at least one dicritical at
infinity, with image $L'_\infty$; since $G$ is birational, there is exactly one such curve,
and \eqref{td:eq:count-degree} gives its degree as $1 \cdot 1 = 1$.
\end{proof}

\subsection{The October 1 main theorem is false}

\begin{corollary}[The theorem as worded fails]\label{td:cor:count-fails}
\textit{Status: grade B in general; exact in the examples below.}
Every Keller map of degree at least $2$ has at least two exceptional dicriticals in
$\Lambda_0$, after a translation of the target that puts $0$ in
$F(\mathbb{A}^2) \setminus S_F$. An automorphism needs no translation. Hence the October 1
draft's theorem (``no two distinct dicritical divisors,'' with no at-infinity restriction) is
false for every Keller map of degree at least $2$. If only exceptional curves count as
divisors of the resolution, it holds exactly for affine maps, whose dicritical at infinity is
$L_\infty$ itself; if $L_\infty$ counts, it fails for them too.
\end{corollary}

\begin{proof}
Let $F$ have degree at least $2$. Translate the target so that $0 \notin S_F$ and
$0 \in F(\mathbb{A}^2)$; then $F^{-1}(0)$ is nonempty. By Proposition~\ref{td:prop:count}(a),
each point of $F^{-1}(0)$ contributes an affine dicritical of degree $1$, and by part (c)
there is at least one dicritical at infinity. That is already two dicriticals in total. The
example below exhibits both of them explicitly for an automorphism, where no translation is
needed.
\end{proof}

\begin{example}[Two dicriticals of an automorphism]\label{td:ex:automorphism}
Take
\begin{equation}\label{td:eq:auto-map}
P = x - y^2, \qquad Q = y.
\end{equation}
We compute the Jacobian step by step. First the partial derivatives:
\[
P_x = 1, \qquad P_y = -2y, \qquad Q_x = 0, \qquad Q_y = 1.
\]
Hence
\[
[P,Q] = P_x Q_y - P_y Q_x = 1 \cdot 1 - (-2y) \cdot 0 = 1 \in K^\times,
\]
so $F = (P,Q)$ is a Keller map. It is an automorphism: solving $P = x - y^2$ and $Q = y$
gives $y = Q$ and $x = P + Q^2$, so the inverse is $(P,Q) \mapsto (P + Q^2, Q)$, and indeed
$P + Q^2 = (x - y^2) + y^2 = x$.

\emph{The affine dicritical.} The fiber over $0$ is the single point $(0,0)$, since
$Q = y = 0$ and $P = x - y^2 = 0$ force $x = 0$. Blow up $\mathbb{A}^2$ at the origin and
use the chart
\[
x = s, \qquad y = s t,
\]
so the exceptional divisor $E$ is $\{s = 0\}$ with coordinate $t$. Substituting:
\begin{align}
P &= s - (st)^2 = s - s^2 t^2 = s\,(1 - s t^2), \\
Q &= s t.
\end{align}
The map $\varphi_0$ is the projection from $0$, so on this chart it is the ratio
\[
\varphi_0 = [P : Q] = [\,s(1 - s t^2) : s t\,] = [\,1 - s t^2 : t\,].
\]
Restricting to the exceptional divisor, where $s = 0$, gives
\[
\varphi_0|_E = [1 : t],
\]
which is non-constant in $t$. Hence $E$ is dicritical, of degree $1$.

\emph{The dicritical at infinity.} By Corollary~\ref{td:cor:automorphisms}, $\Lambda_0$ has
exactly one dicritical at infinity, the unique curve $D$ with $G(D) = L'_\infty$, also of
degree $1$. Together with the affine dicritical above, this is two distinct dicriticals.
The October 1 theorem, which forbade two distinct dicritical divisors without any
at-infinity restriction, is therefore false, already for automorphisms.
\end{example}

\subsection{The corrected statement: Theorem R}

\begin{theorem}[Theorem R --- the corrected main theorem; OPEN]\label{td:thm:R}
\textit{Status: open, not proved in this paper or elsewhere.}
For every Keller map $(P,Q)$, the pencil $\lambda P + \mu Q$ has at most one dicritical
divisor at infinity.
\end{theorem}

\begin{proposition}\label{td:prop:onedic}
\textit{Status: grade B.}
Let $F$ be a Keller map and $c \notin S_F$. Then $\Lambda_c$ has at most one dicritical at
infinity if and only if $F$ is an automorphism. In that case it has exactly one, of degree $1$.
\end{proposition}

\begin{proof}
If $F$ is an automorphism, Corollary~\ref{td:cor:automorphisms} gives exactly one dicritical at
infinity, of degree $1$. Conversely, suppose $\Lambda_c$ has at most one dicritical at
infinity. By Proposition~\ref{td:prop:count}(c), the number of dicriticals at infinity is at
least $1 + \#\{\text{components of } S_F\}$. At most one therefore forces
$\#\{\text{components of } S_F\} = 0$, that is, $S_F = \emptyset$. By (F3), $F$ is an
automorphism.
\end{proof}

\begin{proposition}[Theorem R is equivalent to the planar Jacobian conjecture]\label{td:prop:equiv}
\textit{Status: grade B.}
Theorem~R holds for all Keller maps if and only if the planar Jacobian conjecture holds.
\end{proposition}

\begin{proof}
If the conjecture holds, then every Keller map is an automorphism, and
Corollary~\ref{td:cor:automorphisms} gives exactly one dicritical at infinity, in particular at
most one. So Theorem~R holds.

Conversely, let $F$ be a Keller map and pick $c \notin S_F$. Such a $c$ exists because $S_F$
is either empty or a curve, never the whole target. Then $F - c$ is a Keller map with
$0 \notin S_{F-c}$, and its pencil $\lambda (P - c_1) + \mu (Q - c_2)$ is $\Lambda_c(F)$.
Applying Theorem~R to $F - c$ and Proposition~\ref{td:prop:onedic} to $\Lambda_c(F)$ makes $F$
an automorphism. Since $F$ was arbitrary, the conjecture holds.
\end{proof}

\begin{remark}
Proposition~\ref{td:prop:equiv} sets the bar for this paper: a complete proof of the corrected
theorem would be a proof of the conjecture. Everything that follows is therefore measured
against that bar, and no step below is presented as closing it.
\end{remark}

\subsection{Where a counterexample's dicriticals would sit}\label{td:sec:where}

Let $F$ be a Keller map that is not an automorphism, and let $c \notin S_F$. Then
$\Lambda_c$ has:
\begin{itemize}[nosep]
\item $e \geq 2$ affine dicriticals, by (F3) and Proposition~\ref{td:prop:count}(a);
\item at least one $E_\infty$ with $G(E_\infty) = L'_\infty$, of degree
$\deg(E_\infty \to L'_\infty)$;
\item at least one $E_S$ over each component $C$ of $S_F$, of degree
$\deg(E_S \to C) \cdot \deg C$, by \eqref{td:eq:count-degree}.
\end{itemize}
$E_\infty$ always exists, so Theorem~R amounts to showing that $E_S$ cannot exist. The two
halves of the October 1 draft, read against this picture, were: the first half claimed there
is no dicritical of degree $m \geq 2$ for any Keller map, which applied to $E_S$ would
already exclude every counterexample whose $S_F$ has a component of degree at least $2$; the
second half claimed there are no two dicriticals of degree $1$, which would handle the rest.
Either half, to succeed, must apply to $E_\infty$ and $E_S$ wherever they sit in the
resolution tree: over any point of $L_\infty$, at any depth, and after free or satellite
blow-ups. Sections~\ref{td:sec:depth01}--\ref{td:sec:bridge} record how far each half actually goes.

\section{Dicriticals at infinity of small depth}\label{td:sec:depth01}

This section proves the two depth $0$ and depth $1$ steps toward Theorem~R, each with every
intermediate computation shown, and then records the constant-leading-form argument in its
honest state: a sketch for free centers only. The conclusion of this section is that every
dicritical at infinity of a non-affine Keller map has depth at least $2$. Nothing here touches
depth $2$ or beyond, except for the sketch in \S\ref{td:sec:constant-form}, whose limits are
stated there.

\subsection{The line at infinity}\label{td:sec:linf}

\begin{lemma}[The line at infinity]\label{td:lem:linf}
\textit{Status: grade B, proof by hand; the binary-form step is executed in
\filename{base\_point\_at\_infinity\_check.py}, check $1$.}
For a non-affine Keller map, $L_\infty$ is not dicritical for any $\Lambda_c$.
\end{lemma}

\begin{proof}
Let $d = \max(\deg P, \deg Q) \geq 2$, the non-affine case. Write $P_d$ and $Q_d$ for the
homogeneous components of degree $d$ of $P$ and $Q$ (zero if the corresponding degree is
smaller than $d$). The restriction $\varphi_c|_{L_\infty}$ is the rational function
$P_d / Q_d$ on $L_\infty \cong \mathbb{P}^1$, where the point $c$ contributes only a constant
shift that does not affect constancy. We show this ratio is constant, in two cases.

\emph{Case 1: unequal degrees.} Suppose $\deg P \neq \deg Q$, say $\deg P > \deg Q$ (the other
case is symmetric). Then on $L_\infty$, the term of highest degree in the pencil is $P_d$
alone, so $P_d / Q_d$ is $\infty$ wherever $Q_d$ vanishes to lower order, and the restriction
is the constant $\infty$. Explicitly: homogenizing, the leading part of a member
$\lambda P + \mu Q$ at infinity is $\lambda P_d$ (since $\deg Q < d$), so the pencil has the
single point at infinity determined by $P_d$ as its value, independent of the member. The
restriction is constant, so $L_\infty$ is not dicritical. The case $\deg Q > \deg P$ gives
the constant $0$ by the same computation.

\emph{Case 2: equal degrees.} Suppose $\deg P = \deg Q = d$. The Jacobian $[P,Q]$ is a nonzero
constant, so in particular its homogeneous component of degree $2d - 2$ vanishes. That
component is exactly $[P_d, Q_d]$, because differentiating lowers degree by $1$ in each factor
and the bracket of the top components is the top component of the bracket:
\begin{equation}\label{td:eq:top-bracket}
[P,Q]_{2d-2} = [P_d, Q_d] = 0.
\end{equation}
So it suffices to prove the following binary-form lemma.

\begin{lemma}[Proportionality of binary forms with vanishing bracket]\label{td:lem:binary}
Let $f, g$ be binary forms (homogeneous polynomials in two variables) of the same degree
$\geq 1$. If $[f,g] = 0$, then $f = r g$ for some constant $r \in K$.
\end{lemma}

\begin{proof}[Proof of Lemma~\ref{td:lem:binary}, step by step]
Work in the rational function field and set $h = f / g$, a rational function of degree $0$,
meaning $h(tx, ty) = h(x,y)$ for all $t$. Differentiating this homogeneity relation with
respect to $t$ at $t = 1$ gives Euler's relation for $h$:
\begin{equation}\label{td:eq:euler-h}
x h_x + y h_y = 0.
\end{equation}
Next, the bracket is bilinear and satisfies the quotient rule in the first argument:
\[
[h, g] = \frac{[f,g]}{g} - \frac{f [g,g]}{g^2} = \frac{[f,g]}{g} = 0,
\]
using $[g,g] = 0$ and the hypothesis $[f,g] = 0$. Writing the bracket out,
\begin{equation}\label{td:eq:bracket-hg}
h_x g_y - h_y g_x = 0.
\end{equation}
Now suppose $(h_x, h_y) \neq (0,0)$. Equation~\eqref{td:eq:bracket-hg} says the vector
$(h_x, h_y)$ is proportional to $(g_x, g_y)$: both are orthogonal to the same direction, so
there is a scalar $\lambda$ with $(g_x, g_y) = \lambda (h_x, h_y)$ wherever the vectors are
defined. Equation~\eqref{td:eq:euler-h} says $(h_x, h_y)$ is proportional to $(-y, x)$: indeed,
$x h_x + y h_y = 0$ means $(h_x, h_y)$ is orthogonal to $(x,y)$, and in two variables the
orthogonal direction is spanned by $(-y, x)$. Combining the two proportionalities,
$(g_x, g_y)$ is proportional to $(-y, x)$, so
\[
x g_x + y g_y = 0.
\]
But Euler for the homogeneous form $g$ of degree $\deg g \geq 1$ says
\[
x g_x + y g_y = (\deg g) \cdot g \neq 0.
\]
This is a contradiction. Hence the supposition is false, and $(h_x, h_y) = (0,0)$. In
characteristic $0$, a rational function with vanishing partial derivatives is constant, so
$h = r$ for some $r \in K$, that is, $f = r g$.
\end{proof}

Returning to Case 2: Lemma~\ref{td:lem:binary} applied to $f = P_d$, $g = Q_d$ gives
$P_d = r Q_d$, so $\varphi_c|_{L_\infty} = P_d / Q_d = r$ is constant. Hence $L_\infty$ is not
dicritical. This completes both cases.
\end{proof}

\subsection{First blow-ups at infinity}\label{td:sec:e1lemma}

\begin{lemma}[First blow-ups at infinity]\label{td:lem:e1}
\textit{Status: algebra grade A, in \filename{lean/TwoDicritical/E1Keller.lean}
(\texttt{e1\_lead\_proportional}, \texttt{jac2\_eq\_zero\_of\_natDegree\_zero}); chart step
grade B, executed in \filename{lean/check\_e1\_index.py}, check $5$, not Lean-formalized.}
For a Keller map, the exceptional curve of the blow-up at any point of $L_\infty$ is not
dicritical. This holds for every degree $m$, not only for $m = 2$.
\end{lemma}

\begin{proof}
A linear change of coordinates keeps the Keller condition (the Jacobian transforms by the
constant determinant of the change), so we may move the point of $L_\infty$ to
$[1 : 0 : 0]$. Let $E_1$ be the exceptional curve of the blow-up at this point, with the
chart coordinate $t_1$ on $E_1$ as in \filename{lean/check\_e1\_index.py}. The proof has three
steps: the restriction formula, the Wronskian computation, and the degenerate cases.

\emph{Step 1: the restriction formula.} Suppose $P$ and $Q$ have the same $x$-degree
$n \geq 1$, with $x^n$-coefficients $A, D \in K[y]$, so that
\begin{equation}\label{td:eq:lead-expansion}
P = A(y)\, x^n + (\text{terms of $x$-degree} < n), \qquad
Q = D(y)\, x^n + (\text{terms of $x$-degree} < n).
\end{equation}
In the $E_1$ chart, the ratio $P/Q$ restricts to $E_1$ as
\begin{equation}\label{td:eq:e1-restriction}
\frac{P}{Q}\bigg|_{E_1} = \frac{A(1/t_1)}{D(1/t_1)}.
\end{equation}
The reason is that on $E_1$ only the top $x$-degree terms survive: lower $x$-degree terms
vanish to higher order along $E_1$, and the chart coordinate $t_1$ inverts the $y$-variable,
which is why $A$ and $D$ are evaluated at $1/t_1$. If instead the $x$-degrees of $P$ and $Q$
differ, say $\deg_x P > \deg_x Q$, then the numerator dominates and the restriction is the
constant $\infty$; if $\deg_x P < \deg_x Q$, it is the constant $0$. In both differing cases
$E_1$ is already not dicritical, so only the equal-degree case remains.

\emph{Step 2: the Wronskian computation.} Differentiate the expansions
\eqref{td:eq:lead-expansion} term by term. The top $x$-degree contributions are:
\begin{align}
P_x &= n A(y)\, x^{n-1} + (\text{$x$-degree} < n-1), \\
P_y &= A'(y)\, x^{n} + (\text{$x$-degree} < n), \\
Q_x &= n D(y)\, x^{n-1} + (\text{$x$-degree} < n-1), \\
Q_y &= D'(y)\, x^{n} + (\text{$x$-degree} < n).
\end{align}
Now form the Jacobian and collect the top $x$-degree. The product $P_x Q_y$ has top
$x$-degree $(n-1) + n = 2n - 1$, with coefficient
\[
(n A) \cdot D' = n A D'.
\]
The product $P_y Q_x$ likewise has top $x$-degree $n + (n-1) = 2n - 1$, with coefficient
\[
A' \cdot (n D) = n A' D.
\]
Every other product of a term from $P_x, P_y$ with a term from $Q_x, Q_y$ has $x$-degree
strictly less than $2n - 1$: lowering the $x$-degree of either factor by at least $1$ lowers
the total by at least $1$. Therefore the $x^{2n-1}$-coefficient of $[P,Q]$ is exactly
\begin{equation}\label{td:eq:wronskian-coeff}
[P,Q]_{x^{2n-1}} = n\,\big(A D' - A' D\big).
\end{equation}
Since $(P,Q)$ is Keller, $[P,Q]$ is a nonzero constant, so every coefficient of a positive
power of $x$ in $[P,Q]$ vanishes. In particular the coefficient \eqref{td:eq:wronskian-coeff}
vanishes, and since $n \geq 1$,
\[
A D' - A' D = 0.
\]
This is the vanishing of the Wronskian of $A$ and $D$. In characteristic $0$ it forces $A$
and $D$ to be proportional: indeed, where $D \neq 0$,
\[
\left(\frac{A}{D}\right)' = \frac{A' D - A D'}{D^2} = 0,
\]
so $A/D$ is constant, say $A = r D$ for $r \in K$. (This is the Lean lemma
\texttt{e1\_lead\_proportional}.) Substituting into the restriction formula
\eqref{td:eq:e1-restriction} gives
\[
\frac{P}{Q}\bigg|_{E_1} = \frac{r D(1/t_1)}{D(1/t_1)} = r,
\]
a constant. Hence $E_1$ is not dicritical in the equal-degree case either.

\emph{Step 3: the degenerate case.} If both $x$-degrees are $0$, then $P, Q \in K[y]$ are
functions of $y$ alone, so $P_x = Q_x = 0$ and $[P,Q] = 0$. This contradicts the Keller
condition $[P,Q] \in K^\times$ (Lean lemma \texttt{jac2\_eq\_zero\_of\_natDegree\_zero}). So
this case cannot occur for a Keller map.

Since a linear change of coordinates moves any point of $L_\infty$ to $[1:0:0]$ while
preserving the Keller condition, the same conclusion holds over every point at infinity.
\end{proof}

\begin{corollary}[Depth at least $2$]\label{td:cor:depth2}
\textit{Status: follows from Lemmas~\ref{td:lem:linf} and~\ref{td:lem:e1}.}
Every dicritical at infinity of a non-affine Keller map has depth at least $2$.
\end{corollary}

\begin{proof}
Depth $0$ is $L_\infty$ itself, excluded by Lemma~\ref{td:lem:linf} for non-affine maps.
Depth $1$ curves are the exceptional curves of blow-ups at points of $L_\infty$, excluded by
Lemma~\ref{td:lem:e1}. The smallest remaining depth is $2$.
\end{proof}

\subsection{The constant-leading-form argument: a sketch for free centers only}\label{td:sec:constant-form}

\begin{proposition}[Constant leading form --- proof sketch]\label{td:prop:constant-form}
\textit{Status: grade C, proof sketch; free centers only. This is not a proof for satellite
chains or for depth $\geq 2$ in general (obligation O3).}
Suppose $E_k$ ($k \geq 2$) is dicritical of order $m \geq 1$ for the pencil. Then after the
standard normalization, the leading coefficient satisfies $\deg_{u'} U(u',0) = 0$.
\end{proposition}

\begin{proof}[Proof sketch]
Let $E_k$ ($k \geq 2$) be an exceptional divisor in the resolution, with local coordinates
$(u',v')$ in which $E_k = \{v' = 0\}$. For a polynomial $S \in K[x,y]$ of degree $d$, write
the pullback as
\begin{equation}\label{td:eq:pullback}
S = {v'}^{-\nu}\,U(u',v')\cdot(\text{monomial in }u'),
\end{equation}
where $\nu = \nu_{E_k}(S)$ is the order of vanishing and $U$ is a unit (nonzero at the
center). The dicritical condition of order $m$ requires the $u'$-degree of the pencil
restriction to equal $m$.

Track the unit $U$ through the blow-up sequence, one blow-up at a time. Consider a single
blow-up at a center with $u_0 \neq 0$ (a \emph{free} center: the center does not lie on another
component of the total transform, so the local coordinate $u'$ is nonzero there). In the
blow-up chart, substituting the center shifts the leading coefficient by the constant factor
$u_0^{-2}$: schematically, if before the blow-up the leading term is $c \cdot {u'}^j$, then
after the blow-up it is $c\, u_0^{-2} \cdot {u''}^{j}$ in the new coordinate $u''$, with the
same exponent $j$ and a nonzero constant multiple. The point is that multiplication by the
nonzero constant $u_0^{-2}$ changes no $u'$-degree: the exponent $j$ is carried through
unchanged, and no new $u'$-dependence is created, because the center coordinate $u_0$ is a
nonzero constant rather than a function of $u'$.

Repeating this step along the chain: each blow-up at a center with $u_0 \neq 0$ multiplies
the leading coefficient by a nonzero constant and preserves its $u'$-degree. A direct
computation in the normalized chart then shows the $u'$-degree of $U(u',0)$ is $0$. The
verification of the degree count is in \filename{general\_constant\_form.py}
(\S\ref{td:sec:correspondence}).

\emph{Where the sketch stops.} The step above uses $u_0 \neq 0$ at every blow-up. If a center
lies on another component of the total transform (a satellite center), then the corresponding
coordinate vanishes at the center, the factor is no longer a nonzero constant, and the
degree-tracking argument as written does not go through. The sketch therefore covers free
chains only. It does not prove that no $E_k$ ($k \geq 2$) is dicritical of order
$m \geq 2$ in general, and it is not presented as doing so: the general $m \geq 2$ statement
at depth $\geq 2$ is obligation O3 in \S\ref{td:sec:open}. The middle-divisor ($E_2$) and leaf
cases ($E_3$, $E_4$) are checked separately as computations in \filename{middle\_e2.py},
\filename{leaf\_e3.py}, and \filename{leaf\_e4.py}; those are fixed instances, not the
general argument.
\end{proof}

\begin{remark}
In the October 1 draft, Proposition~\ref{td:prop:constant-form} was followed by a corollary
claiming no $E_k$ ($k \geq 2$) is dicritical of order $m \geq 2$ for any Keller map, and that
corollary was then used to close the $m \geq 2$ half of the main theorem. That inference is
withdrawn: a sketch for free centers does not yield the general corollary, and the general
corollary is not claimed anywhere in this revision.
\end{remark}

\section{The first exceptional divisor: a false obstruction and its correction}\label{td:sec:e1}

This section does two things the October 1 draft got backwards. First, it exhibits the draft's
$E_1$ degree obstruction as false, shows the exact step where its proof goes wrong, and
computes the counterexample in full. Second, it records the corrected statement, which was
already proved as Lemma~\ref{td:lem:e1} in \S\ref{td:sec:e1lemma}: first blow-ups at infinity are
never dicritical for a Keller map, for every degree $m$. The correction is strictly stronger
than the draft's claim in scope (every $m$, not only $m = 2$) and strictly weaker in what it
asserts about $E_1$ itself (it does not claim anything about a hypothetical $m = 2$ open
subset, which, as the counterexample shows, contains Keller pairs).

\subsection{What the October 1 draft claimed}

The draft's $E_1$ degree obstruction theorem read, in substance: let
$P, Q$ lie in the $E_1$ $m = 2$ open subset, meaning $\nu_{E_1} = -2$ and both pencil
restrictions have exact degree $2$; then $[P,Q]$ is never a nonzero constant. Its proof argued
as follows: every monomial $x^J y^K$ allowed by the order condition $\nu_{E_1} \geq -2$
satisfies $J + K \geq 2$; hence $P_x, P_y, Q_x, Q_y$ have no constant term; hence $[P,Q]$ has
minimum total degree $2$ and cannot be a nonzero constant. Every step after the first
sentence is correct arithmetic. The first sentence is false, and this subsection shows why,
step by step.

\subsection{The order condition, computed correctly}

\begin{lemma}[The $E_1$ order condition]\label{td:lem:e1-order}
In the $E_1$ chart, the valuation satisfies
\begin{equation}\label{td:eq:e1-valuation}
\nu_{E_1}(x^i y^j) = -i.
\end{equation}
Consequently $\nu_{E_1} \geq -2$ holds for a monomial $x^i y^j$ if and only if $i \leq 2$,
that is, if and only if the $x$-degree is at most $2$. It says nothing about the total degree
$i + j$.
\end{lemma}

\begin{proof}
Equation~\eqref{td:eq:e1-valuation} is the normalization of the $E_1$ valuation used in the
$E_1$ chart (see \filename{lean/check\_e1\_index.py} and
\filename{lean/TwoDicritical/E1Keller.lean}): the variable $x$ has a pole of order $1$ along
$E_1$, contributing $-1$ per power of $x$, while $y$ contributes $0$. Hence for the monomial
$x^i y^j$,
\[
\nu_{E_1}(x^i y^j) = i \cdot \nu_{E_1}(x) + j \cdot \nu_{E_1}(y) = i \cdot (-1) + j \cdot 0 = -i.
\]
The inequality $\nu_{E_1} \geq -2$ is therefore
\[
-i \geq -2 \quad\Longleftrightarrow\quad i \leq 2.
\]
In words: the condition bounds the $x$-degree by $2$. It does not bound $j$, and it does not
bound the total degree $i + j$ from below: the monomial $y$ itself, with $i = 0$ and $j = 1$,
satisfies $\nu_{E_1}(y) = 0 \geq -2$ while having total degree $1$. The draft's claim that
every allowed monomial has $J + K \geq 2$ is thus false on its face, once the valuation is
written out.
\end{proof}

\subsection{How the mix-up happened}

The draft's error is an index mix-up, identified in the Lean analysis
(\filename{lean/README.md}, item 2). The computation \filename{e1\_solitary.py} indexes its
unknowns $p[J,K]$ by the homogeneous monomial $X^{d-J-K} Y^J Z^K$, that is, by the affine
monomial
\[
x^{d-J-K} y^J.
\]
The draft's proof, and the scripts \filename{e1\_keller\_search.py} and the Lean file
\filename{lean/TwoDicritical/E1Obstruction.lean}, read $p[J,K]$ as the monomial $x^J y^K$
instead. Under that misreading, the condition on the indices becomes a condition on the total
degree $J + K$, and the draft concluded that every allowed monomial has total degree at least
$2$. At $d = 4$, both index sets have $12$ monomials, so the count agreed and the mix-up
was hidden: the $12$-dimensional span at $d = 4$ is real, but it is a span of monomials of
$x$-degree at most $2$, not of total degree at least $2$. The file
\filename{lean/TwoDicritical/E1Obstruction.lean} proves a true lemma about the misread set of
monomials (if every monomial of $P$ and $Q$ has total degree at least $2$, then $[P,Q]$ has
no constant term), but that set is not the $E_1$ condition, so the lemma does not prove
anything about $E_1$.

\subsection{The counterexample, computed in full}

\begin{example}[A Keller pair in the draft's $E_1$ open subset]\label{td:ex:e1-counter}
\textit{Status: grade A in Lean (\filename{lean/TwoDicritical/E1Keller.lean},
\texttt{thm\_e1\_false}); the computation below is exact.}
Set
\begin{equation}\label{td:eq:e1-counter}
P = (x + y^2)^2 + y, \qquad Q = P + x + y^2.
\end{equation}
We expand $P$ first, term by term:
\[
P = x^2 + 2 x y^2 + y^4 + y.
\]
Its partial derivatives are:
\begin{align}
P_x &= 2x + 2y^2 = 2\,(x + y^2), \\
P_y &= 4x y + 4y^3 + 1 = 4y\,(x + y^2) + 1.
\end{align}
Now compute $[P,Q]$. Since $Q = P + x + y^2$ and $[P,P] = 0$, bilinearity of the bracket
gives
\[
[P,Q] = [P, x] + [P, y^2].
\]
Each bracket is computed separately. For $[P,x]$: $x$ has derivatives $1$ and $0$, so
\[
[P,x] = P_x \cdot 0 - P_y \cdot 1 = -P_y = -\big(4y(x+y^2) + 1\big).
\]
For $[P, y^2]$: $y^2$ has derivatives $0$ and $2y$, so
\[
[P, y^2] = P_x \cdot 2y - P_y \cdot 0 = 2y P_x = 2y \cdot 2(x+y^2) = 4y(x+y^2).
\]
Adding the two contributions, the $4y(x+y^2)$ terms cancel exactly:
\begin{equation}\label{td:eq:e1-bracket}
[P,Q] = -\big(4y(x+y^2) + 1\big) + 4y(x+y^2) = -1 \in K^\times.
\end{equation}
So $(P,Q)$ is a Keller pair.

It lies in the draft's $E_1$ $m = 2$ open subset at $d = 4$: its $x$-degree is $2$ (the term
$x^2$), so $\nu_{E_1} = -2$ by Lemma~\ref{td:lem:e1-order}, and each pencil restriction is the
function $C(t_1) = t_1^2$ of exact degree $2$ (this is the membership statement formalized as
\texttt{thm\_e1\_false} in \filename{lean/TwoDicritical/E1Keller.lean}, with the open subset
denoted \texttt{E1Open4} there). Yet $[P,Q] = -1$ is a nonzero constant, contradicting the
draft's theorem directly.

Two further observations pin down the failure. First, $P$ contains the monomial $y$, of total
degree $1$: the draft's claimed lower bound $J + K \geq 2$ fails for a monomial that is
actually present. Second, $P/Q$ restricts to $1$ on $E_1$ for this pair, so $E_1$ is not
dicritical for them at all: the draft's ``open subset'' does not even capture an $m = 2$
dicritical. The draft's proof argued about the wrong monomials, and its conclusion is false.
\end{example}

\begin{remark}[Withdrawn theorem]
The draft's $E_1$ degree obstruction theorem and its corollary (``no dicritical of order
$m \geq 2$ exists for any Keller map'') are both withdrawn. The corollary depended on the false theorem for the $E_1$ case and on the
free-only sketch of Proposition~\ref{td:prop:constant-form} for the $E_k$, $k \geq 2$, case;
neither leg supports it as stated.
\end{remark}

\subsection{The corrected \texorpdfstring{$E_1$}{E1} statement}

The intended conclusion of the draft's $E_1$ argument --- that the first blow-up cannot be
dicritical for a Keller map --- is true, and it is Lemma~\ref{td:lem:e1} of
\S\ref{td:sec:e1lemma}, proved there for every degree $m$. To recap the logic in one place: the
restriction of $P/Q$ to $E_1$ is $A(1/t_1)/D(1/t_1)$ when the $x$-degrees agree (and $0$ or
$\infty$ when they differ), and the Keller condition forces the Wronskian $AD' - A'D$ to
vanish via the $x^{2n-1}$-coefficient \eqref{td:eq:wronskian-coeff}, hence $A = rD$ and the
restriction is the constant $r$. The algebra is Lean-grade-A; the chart step is executed, not
Lean-formalized (obligation O5). No $m = 2$ open subset, and no total-degree argument, is
used anywhere in the corrected proof.

\section{The \texorpdfstring{$m=1$}{m=1} bridge, under its premises}\label{td:sec:bridge}

This section keeps the October 1 draft's Vandermonde and block-rank argument, which is correct
as linear algebra, but restates its scope exactly. The argument proves a collapse of a system
of equations. Whether that system is the system of two $m = 1$ dicriticals depends on two
premises, (P1) and (P2), stated below. Premise (P2) is not derived anywhere in this paper; it
is an open obligation (O2). The theorem of this section is therefore a theorem about equations
under premises, not an unconditional geometric statement, and its corollaries are fenced
accordingly.

\subsection{The chart and the two premises}

Read through its chart, the bridge concerns the following geometry. The paper's pullback is
\[
S^h\big(1, b, b^2(a + \tilde{a} b)\big),
\]
the chart at the point $w_2 = a$ of $E_2$ over $[1 : 0 : 0]$, where $w_2 = x / y^2$. The two
centers $a_3 \neq a_4$ are points of this $E_2$. The bridge argument rests on:
\begin{itemize}[nosep]
\item[(P1)] the two dicriticals lie over two distinct \emph{free} points of this $E_2$;
\item[(P2)] every member of the pencil has pole order at most $1$ on them.
\end{itemize}
Premise (P1) is part of the geometric reading of the equations: the chart above is the chart
at a free point of $E_2$, and the argument is set in that chart. Premise (P2) is a hypothesis
on the members, not a consequence of the definitions used here: the paper defines $m$ as the
degree of $\varphi_c|_E$, and degree $1$ does not force pole order $1$
(\filename{lean/check\_satellite.py}, check $7$). Deriving (P2) from the definition of $m$,
or supplying a separate argument for higher pole order, is obligation O2 in
\S\ref{td:sec:open}. Nothing in this section assumes (P2) silently: every statement that needs
it says so.

\subsection{Setup of the equations}

Consider two distinct centers $a_3 \neq a_4$ (both nonzero), each requiring order
$\operatorname{ord}_b(S) \geq 2d - 1$ (the $m = 1$ condition, under (P2)). For
$S = \sum p_{JK} x^J y^K$, the pullback at center $a$ is
\begin{equation}\label{td:eq:bridge-pullback}
S_a(b) = \sum_{J,K} p_{JK}\,(a + \tilde{a} b)^{d-J-K}\, b^{2d-2J-K},
\end{equation}
and we require the coefficients of $b^0, \dots, b^{2d-2}$ to vanish at $a = a_3$ and at
$a = a_4$. The exponent $d - J - K$ is the remaining degree after the monomial $x^J y^K$ is
factored out, and the exponent $2d - 2J - K$ is the $b$-order contributed by that monomial in
the chart; the condition $\operatorname{ord}_b(S) \geq 2d - 1$ is exactly the vanishing of
all coefficients below $b^{2d-1}$.

\begin{theorem}[Bridge collapse, under (P1) and (P2)]\label{td:thm:bridge}
\textit{Status: equations grade A, in \filename{lean/TwoDicritical/BridgeCollapse.lean};
geometric reading grade B, executed in \filename{lean/check\_satellite.py}.}
For arbitrary distinct nonzero $a_3, a_4$ and arbitrary $d \geq 2$, the simultaneous
conditions $\operatorname{ord}_b(S) \geq 2d - 1$ at both centers force $p_{JK} = 0$ whenever
$J \geq 1$. The nullspace is exactly $\operatorname{span}\{1, y\}$.
\end{theorem}

In particular, every $P, Q$ in the span is a function of $y$ alone, so $[P,Q] \equiv 0$: no
Keller pair exists in the system of equations. Whether this rules out two geometric $m = 1$
dicriticals depends on (P1) and (P2); see \S\ref{td:sec:bridge-scope}.

\subsection{Descending induction on weighted degree}

Order monomials by the weighted degree $W = 2J + K$, descending. For fixed $W$, write the
variables as $(J_k, K_k)$, $k = 1, \dots, n$, ordered by descending $J_k$. Define
\begin{equation}\label{td:eq:ek-def}
e_k = d - J_k - K_k = d - W + J_k.
\end{equation}
Since $W$ is fixed and the $J_k$ are strictly decreasing by the ordering, the $e_k$ are
strictly decreasing, hence distinct. This distinctness is the whole input to the rank lemma
below, so we record it explicitly: for $k < l$,
\[
e_k - e_l = J_k - J_l > 0.
\]

The equations for weight $W$ come from the pairs $(t, i)$ with $2d - t + i = W$, where $t$ is
the $b$-power and $i$ the $\tilde{a}$-power, at both centers. To see why, expand the factor
$(a + \tilde{a} b)^{e_k}$ in \eqref{td:eq:bridge-pullback} by the binomial theorem, term by
term:
\[
(a + \tilde{a} b)^{e_k} = \sum_{i \geq 0} \binom{e_k}{i} a^{e_k - i} (\tilde{a} b)^i.
\]
Multiplying by $b^{2d - 2J_k - K_k}$ and using $W = 2J_k + K_k$, the monomial $p_{J_k K_k}$
contributes to the coefficient of $b^t$ exactly the terms with $2d - W + i = t$, that is,
$i = t - (2d - W)$. Collecting over $k$ gives the $(t, i, a)$ equation:
\begin{equation}\label{td:eq:block}
\sum_{k=1}^n p_{J_k K_k}\binom{e_k}{i} a^{e_k-i} = 0, \qquad a \in \{a_3, a_4\}.
\end{equation}
Crucially, the equations \eqref{td:eq:block} for weight $W$ involve only variables of weight
exactly $W$ at the top: variables of weight $> W$ contribute to higher $b$-powers or have
already been eliminated, as the induction below makes precise.

\begin{lemma}[Block full rank]\label{td:lem:block-rank}
For distinct $e_1, \dots, e_n$ and $m \geq n$, the $2m \times n$ matrix with entries
$\binom{e_k}{i} a^{e_k-i}$ ($i = 0, \dots, m-1$, $a \in \{a_3, a_4\}$) has rank $n$.
\end{lemma}

\begin{proof}
We reduce the matrix in three explicit steps, each preserving rank.

\emph{Step 1: factor the columns and rows.} Factor column $k$ by $a^{e_k} \neq 0$ (note $a$
is $a_3$ or $a_4$, both nonzero by hypothesis) and row $(i, a)$ by $a^{-i} \neq 0$. Rank is
preserved under nonzero column and row scalings. After this step, the entry in row
$(i, a)$, column $k$ is $\binom{e_k}{i}$, independent of $a$. The two centers now give two
copies of the same rows, so it suffices to show that the $m \times n$ matrix
$[\binom{e_k}{i}]$ ($i = 0, \dots, m-1$) has rank $n$; taking the first $n$ rows
($i = 0, \dots, n-1$) suffices, since $m \geq n$.

\emph{Step 2: triangularize by row operations.} For fixed $i$, the binomial coefficient
$\binom{e}{i}$, viewed as a function of $e$, is a polynomial in $e$ of exact degree $i$:
\[
\binom{e}{i} = \frac{e (e-1) \cdots (e - i + 1)}{i!}
= \frac{1}{i!} e^i + (\text{terms of degree} < i \text{ in } e).
\]
Hence row $i$ of the matrix $[\binom{e_k}{i}]$ equals $\frac{1}{i!}$ times the row
$[e_k^i]$ plus a linear combination of the rows $[e_k^j]$ for $j < i$. Proceeding from
$i = 0$ upward, subtract suitable multiples of earlier rows from row $i$ to kill all
lower-degree terms: row $0$ is already $[1]$; row $1$ is $[e_k] + (\text{constant row})$,
so subtracting the constant multiple of row $0$ leaves $[e_k]$; row $2$ is
$\frac{1}{2}[e_k^2] + (\text{combination of rows } 0, 1)$, so subtracting leaves
$\frac{1}{2}[e_k^2]$; and so on. These are elementary row operations, so they preserve rank
and multiply the determinant by $1$. After all $n$ steps, the matrix is the scaled
Vandermonde matrix $[e_k^i / i!]$ ($i = 0, \dots, n-1$).

\emph{Step 3: the Vandermonde determinant.} The determinant of $[e_k^i]$ is the classical
Vandermonde product $\prod_{j < k} (e_k - e_j)$, nonzero because the $e_k$ are distinct
(shown above from the ordering by $J_k$). Factoring $1/i!$ out of row $i$ gives
\begin{equation}\label{td:eq:vandermonde-det}
\det[\binom{e_k}{i}] = \frac{\prod_{j<k}(e_k - e_j)}{\prod_{i=0}^{n-1} i!} \neq 0.
\end{equation}
A nonzero $n \times n$ minor means the matrix has rank $n$. The $n = 3$ case is verified
symbolically in \filename{m1\_block\_rank.py}, giving
$-(e_1-e_2)(e_1-e_3)(e_2-e_3)/2$, which agrees with \eqref{td:eq:vandermonde-det} since
$\prod_{i=0}^{2} i! = 0! \cdot 1! \cdot 2! = 2$ and the sign is the ordering of the factors.
This completes the proof.
\end{proof}

\begin{proof}[Proof of Theorem~\ref{td:thm:bridge}]
Descend on $W$ from $2d$ to $2$, maintaining the inductive hypothesis that all variables
with weight $> W$ are zero. The base case is $W = 2d$: there are no variables of weight
$> 2d$ (since $W = 2J + K \leq 2d$ for monomials of degree at most $d$), so the hypothesis
holds vacuously.

For the inductive step at weight $W$, the equations \eqref{td:eq:block} for weight $W$ involve
only variables of weight $\geq W$: a variable of weight $W' > W$ contributes to $b$-powers
above those in the weight-$W$ system, and by the inductive hypothesis all such variables are
already $0$. Hence the weight-$W$ system is exactly the block system of
Lemma~\ref{td:lem:block-rank}, which has full column rank $n$. Therefore
$p_{J_k K_k} = 0$ for all $k$, that is, all variables of weight $W$ vanish. This restores the
hypothesis for the next lower weight.

Descending to $W = 2$ eliminates every variable with $J \geq 1$: any monomial with $J \geq 1$
has $W = 2J + K \geq 2$, and the descent covers all weights $W \geq 2$. The only monomials
with $W < 2$ are $(J,K) = (0,0)$ with $W = 0$ and $(J,K) = (0,1)$ with $W = 1$. These are never
constrained by the descent, so the nullspace is exactly $\operatorname{span}\{1, y\}$.
\end{proof}

\subsection{Computational verification}

The induction is verified computationally, at the fixed degrees and instances stated (these
are checks of the linear algebra, not replacements for the proof above):
\begin{itemize}[nosep]
\item Symbolic nullspace at general $a_3 \neq a_4$: $d = 2$ gives nullity $2$; $d = 3$ gives
nullity $2$ (\filename{m1\_bridge\_general.py}).
\item Numeric nullspace at $a_3 = 2$, $a_4 = 5$: nullity $2$ for $d = 2, \dots, 8$
(\filename{m1\_induction.py}).
\item Block-triangular structure exhibited at $d = 6$ (\filename{m1\_triangular.py}).
\item Vandermonde determinant verified symbolically (\filename{m1\_block\_rank.py}).
\end{itemize}

\subsection{Scope of the bridge result}\label{td:sec:bridge-scope}

Under premises (P1) and (P2), the bridge has the following consequences, and no others are
claimed:
\begin{itemize}[nosep]
\item With a single center imposed, the same block argument gives the larger nullspace
$\operatorname{span}\{1, y, x - a y^2\}$ (\filename{lean/TwoDicritical/OneCenter.lean},
\texttt{bridge\_one\_center}). So one such dicritical already forces degree at most $2$:
no pencil of degree at least $3$ has even one dicritical of this chart type.
\item In degree $2$, the automorphism $(x - a y^2, y)$ has exactly one such dicritical,
consistent with Theorem~R (it has one dicritical at infinity, as
Corollary~\ref{td:cor:automorphisms} requires, and no second).
\item The bridge as proved here says nothing about pencils of degree at least $3$ beyond the
one-center remark above, because (P2) is not derived: a dicritical of degree $1$ with pole
order greater than $1$ is not governed by the system \eqref{td:eq:block}. Nor does it cover
dicriticals over satellite points, over different points of $L_\infty$, or at different
depths; those are obligation O1.
\end{itemize}

\subsection{The adjacent case: a computation at \texorpdfstring{$d = 3$}{d=3} only}

The adjacent $m = 1$ case --- two dicriticals sharing a curve in the resolution tree rather
than lying over two free points of one $E_2$ --- is treated in this paper only by computation
at $d = 3$ (\filename{e3\_adjacent.py}, \filename{routeB\_cancellation.py}), covering the
$E_2$--$E_3$ adjacent pair at that degree. \textit{Status: grade C, fixed-degree
computation.} There is no argument here for other degrees; the general adjacent case is
obligation O4 in \S\ref{td:sec:open}. In the October 1 draft, this computation was cited as
closing the adjacent case for every degree. That citation is withdrawn.

\section{What remains open}\label{td:sec:open}

This section replaces the October 1 draft's ``Proof of the main theorem,'' which assembled
the false $E_1$ theorem, the free-only constant-form sketch, the premises-dependent bridge,
and the $d = 3$ adjacent computation into a contradiction, via a case split (``an adjacent
pair or a bridge'') that was asserted, not proved. No such assembly is attempted here. The
obligations below are what a proof of Theorem~R still needs; the tests after them are checks
any proposed proof must pass.

\subsection{Obligations}

\begin{itemize}[nosep]
\item[\textbf{O1.}] \textbf{The case split.} The positions of two dicriticals in a resolution
tree are more varied than ``an adjacent pair or a bridge'': dicriticals can lie over
\emph{different points} of $L_\infty$ (a pencil can have several base points at infinity in
general, as the non-Keller map $(x, xy)$ shows, and whether a Keller counterexample's pencil
can must be settled); over \emph{satellite points} (the point $a = \infty$ of $E_2$ in
$(x + y^3, y)$, and deeper ones); and at \emph{different depths}. A proof must cover every
configuration of $E_\infty$ and $E_S$ from \S\ref{td:sec:where}.
\item[\textbf{O2.}] \textbf{Premise (P2).} Derive pole order at most $1$ from the definition of
$m$, or restate the bridge with (P2) as a hypothesis and supply a separate argument for
higher pole order. Until then, the bridge governs only the system
\eqref{td:eq:bridge-pullback}--\eqref{td:eq:block} under (P1) and (P2).
\item[\textbf{O3.}] \textbf{$m \geq 2$ at depth $\geq 2$.} Prove the constant-leading-form
statement of Proposition~\ref{td:prop:constant-form} for every chain of centers, satellite ones
included. Today there is a sketch for free chains only.
\item[\textbf{O4.}] \textbf{The adjacent case} for every degree, not only $d = 3$.
\item[\textbf{O5.}] \textbf{Two small formalizations:} the chart step of Lemma~\ref{td:lem:e1}
and the binary-form step of Lemma~\ref{td:lem:linf} in Lean. The algebra of Lemma~\ref{td:lem:e1}
is formalized; the chart step is executed, not formalized.
\item[\textbf{O6.}] \textbf{The Keller condition must be used globally.} Theorem~R is false
without it. The map $(x, xy)$, with Jacobian $x$, has two dicriticals at infinity in
$\Lambda_{(1,1)}$, over $[1 : 0 : 0]$ and $[0 : 1 : 0]$; the second maps onto the line $u = 0$
of non-proper values (\filename{restatement\_checks.py}, check $2$). Any argument that never
uses $[P,Q] \in K^\times$ cannot prove Theorem~R. This is a standing check on every step
above: Lemmas~\ref{td:lem:linf} and~\ref{td:lem:e1} use the Keller condition (vanishing of
$[P_d, Q_d]$, vanishing of the Wronskian coefficient), while the bridge's linear algebra and
the constant-form sketch must be read with their premises attached.
\end{itemize}

\subsection{Tests any proof must pass}

\begin{itemize}[nosep]
\item[\textbf{T1.}] Automorphisms have exactly one dicritical at infinity, of degree $1$
(Corollary~\ref{td:cor:automorphisms}), but its position varies: over the free point $a$ of
$E_2$ for $(x - a y^2, y)$ (where $P/Q = \tilde{a}/a$ there), and over the satellite point
$a = \infty$ of $E_2$, at depth $5$, for $(x + y^3, y)$. A lemma that rules out a dicritical
in either position is false.
\item[\textbf{T2.}] The identity. Its dicritical at infinity is $L_\infty$ itself;
Lemma~\ref{td:lem:linf} excludes only non-affine maps.
\item[\textbf{T3.}] Non-Keller maps can have two dicriticals at infinity (O6).
\item[\textbf{T4.}] Every Keller map of degree at least $2$ has two exceptional dicriticals in
total (Corollary~\ref{td:cor:count-fails}), so only the count \emph{at infinity} can be
bounded. Any argument that bounds the total count is refuted by Example~\ref{td:ex:automorphism}.
\end{itemize}

\section{Script-to-arithmetic correspondence}\label{td:sec:correspondence}

Each script below was actually executed; outputs are archived in \filename{outputs/}. The
table maps every script to the exact paper statement it certifies, \emph{under the corrected
statements of this revision}. The grade column follows the restatement note, \S 8: A means
machine-checked, B means a proof by hand from stated facts or an exact computation outside
Lean, C means a sketch or a computation at fixed degrees. A script that certified a withdrawn
statement is listed with what it actually certifies now, not with the draft's reading of it.

\begin{center}
\small
\begin{tabular}{llll}
\textbf{Script} & \textbf{Paper locus} & \textbf{Certifies} & \textbf{Grade} \\ \midrule
\filename{general\_constant\_form.py} & Prop.~\ref{td:prop:constant-form} & Leading-unit degree count, free chains only & C \\
\filename{middle\_e2.py} & Prop.~\ref{td:prop:constant-form} & $E_2$ instance of the sketch & C \\
\filename{leaf\_e3.py} & Prop.~\ref{td:prop:constant-form} & $E_3$ leaf instance of the sketch & C \\
\filename{leaf\_e4.py} & Prop.~\ref{td:prop:constant-form} & $E_4$ leaf instance of the sketch & C \\
\filename{e1\_solitary.py} & \S\ref{td:sec:e1} & $12$-monomial count at $d = 4$ (indexing as computed) & C \\
\filename{e1\_global.py} & \S\ref{td:sec:e1} & Explicit pair, $\nu = -2$, $\deg = 2$ (membership data) & C \\
\filename{e1\_keller\_search.py} & \S\ref{td:sec:e1} & Searched the misread monomial set; superseded & --- \\
\filename{m1\_bridge\_general.py} & Thm.~\ref{td:thm:bridge} & Symbolic nullspace, $d = 2, 3$ & B \\
\filename{m1\_induction.py} & \S\ref{td:sec:bridge} & Nullity $2$ for $d = 2, \dots, 8$ & C \\
\filename{m1\_triangular.py} & \S\ref{td:sec:bridge} & Block-triangular structure at $d = 6$ & C \\
\filename{m1\_block\_rank.py} & Lem.~\ref{td:lem:block-rank} & Vandermonde determinant $\neq 0$ & B \\
\filename{at\_degree\_bound.py} & Prop.~\ref{td:prop:constant-form} & Finite bridge table, $m \leq 4$ & C \\
\filename{at\_degree\_bound\_m5.py} & Prop.~\ref{td:prop:constant-form} & Finite bridge table, $m \leq 5$ & C \\
\filename{mixed\_poles.py} & Prop.~\ref{td:prop:constant-form} & Mixed $(2,3)$ bridge instance & C \\
\filename{depth5\_global.py} & Prop.~\ref{td:prop:constant-form} & Depth-$5$ enumeration & C \\
\filename{depth6\_enum.py} & Prop.~\ref{td:prop:constant-form} & Depth-$6$ enumeration & C \\
\filename{routeB\_cancellation.py} & \S\ref{td:sec:bridge} & Adjacent $m = 1$ at $d = 3$ only & C \\
\filename{e3\_adjacent.py} & \S\ref{td:sec:bridge} & $E_2$--$E_3$ adjacent pair at $d = 3$ & C \\
\end{tabular}
\end{center}

The Lean evidence, which is not a script in the table above, is:
\begin{itemize}[nosep]
\item \filename{lean/TwoDicritical/E1Keller.lean}: \texttt{thm\_e1\_false} (the draft's $E_1$
theorem is false, Example~\ref{td:ex:e1-counter}); \texttt{e1\_lead\_proportional} and
\texttt{jac2\_eq\_zero\_of\_natDegree\_zero} (the algebra of Lemma~\ref{td:lem:e1}); grade A for
the algebra.
\item \filename{lean/TwoDicritical/BridgeCollapse.lean} and
\filename{lean/TwoDicritical/OneCenter.lean}: the bridge equations and the one-center span
(Theorem~\ref{td:thm:bridge}, \S\ref{td:sec:bridge-scope}); grade A for the equations.
\item \filename{lean/check\_e1\_index.py} (check $5$) and
\filename{lean/check\_satellite.py} (check $7$): the $E_1$ chart step and the
degree-versus-pole-order distinction; executed checks, grade B.
\item \filename{base\_point\_at\_infinity\_check.py} (check $1$) and
\filename{restatement\_checks.py} (check $2$): the binary-form step of
Lemma~\ref{td:lem:linf} and the non-Keller example of O6; executed checks, grade B.
\end{itemize}

\section{Computational archive}\label{td:sec:archive}

All scripts used in this work are archived in \filename{scripts/} with their outputs in
\filename{outputs/}, as archived with the October 1 draft:
\begin{itemize}[nosep]
\item 18 Python scripts (SymPy-based), each with a header stating its epistemic status;
\item 18 output files recording exact stdout and exit codes;
\item All scripts exited with code $0$ on re-execution (October 1, 2026).
\end{itemize}
This revision changes no script and claims no new re-execution: the corrections are to the
statements the scripts are read as certifying (\S\ref{td:sec:correspondence}), not to the
scripts themselves.

\subsection*{Reproducibility}
All artifacts are archived in \href{https://github.com/brandonmccraryresearch-cloud/2D-Jacobian-Conjecture-Workspace}{the public repository}, directory \filename{branch\_two\_dicritical/}. Each script runs with \filename{\textasciitilde/miniconda3/envs/physics/bin/python} (Python 3.12, SymPy). The verification script \filename{verify\_all.sh} re-runs every script and checks exit codes.

\section{Computational and AI-assisted research disclosure}\label{td:sec:disclosure}

Muse (Meta), the author's personal AI assistant, assisted with computational scripting,
manuscript drafting, and proof verification; all script outputs were verified by actual
execution. The author (B.~D.~McCrary) directed the investigation, verified the results, and
takes responsibility for the manuscript. This revision was drafted with the same assistance:
the October 1 draft was corrected against the October 13, 2026 restatement note, the
derivations were expanded step by step, and no new theorem, computation, or script output
was introduced. The mathematical argument stands on the artifacts above, not on these
attributions.

\begin{thebibliography}{99}

\bibitem{GGHV} J.~A.~Guccione, J.~J.~Guccione, R.~Horruitiner, C.~Valqui, \textit{Increasing the degree of a possible counterexample to the Jacobian conjecture from 100 to 108}, arXiv:2204.14178v1 (2022).

\bibitem{Stothers} W.~W.~Stothers, \textit{Polynomial identities and Hauptmoduln}, Quart.\ J.\ Math.\ Oxford Ser.\ (2) \textbf{32} (1981), 349--370.

\bibitem{Davenport} H.~Davenport, \textit{On $f^3(t)-g^2(t)$}, Det Kongelige Norske Videnskabers Selskabs Forhandlinger \textbf{38} (1965), 86--87.

\bibitem{Moh} T.~T.~Moh, \textit{On the Jacobian conjecture and the configurations of roots}, J.\ Reine Angew.\ Math.\ \textbf{340} (1983), 140--212.

\bibitem{Santibanez} F.~Santiba\~nez-Leal, \textit{CAOS records EXP-061 and EXP-065}, GitHub repository \href{https://github.com/fsantibanezleal/CAOS_RESEARCH}{\filename{github.com/fsantibanezleal/CAOS\_RESEARCH}} (2026).

\bibitem{SantibanezPaper} F.~Santiba\~nez-Leal, \textit{CAOS v0.34}, Zenodo \href{https://doi.org/10.5281/zenodo.22834215}{doi:10.5281/zenodo.22834215} (2026).

\bibitem{JC125note} B.~D.~McCrary, Muse, \textit{The $(72,108)$ candidate for the planar Jacobian conjecture: a containment note}, Zenodo \href{https://doi.org/10.5281/zenodo.22941057}{doi:10.5281/zenodo.22941057} (2026).

\bibitem{ClaudeBlind} Claude (Anthropic), \textit{Independent blind replication: GGHV Proposition~4.3, normal form~(2)}, \href{https://github.com/brandonmccraryresearch-cloud/2D-Jacobian-Conjecture-Workspace/tree/main/tier1_blind_gghv43_nf2}{\textcolor{magenta}{\filename{tier1\_blind\_gghv43\_nf2/}}} in B.~D.~McCrary, \textit{2D-Jacobian-Conjecture-Workspace}, GitHub (2026), commit \texttt{37c73b67} (2026-09-25).

\bibitem{BranchABv19} B.~D.~McCrary, \textit{Elimination of the GGHV branch-(a,b) normal form at degree pair $(72,108)$: a certified computational audit}, Zenodo \href{https://doi.org/10.5281/zenodo.23023490}{doi:10.5281/zenodo.23023490} (2026).

\bibitem{Helali} H.~Helali, \textit{Exact characteristic-zero certificate for the GGHV $(72,108)$ configuration}, Zenodo \href{https://doi.org/10.5281/zenodo.21479814}{doi:10.5281/zenodo.21479814} (2026).

\bibitem{wstrinz} wstrinz, \textit{Structural analysis of the plane Jacobian $(72,108)$ candidate}, GitHub repository \href{https://github.com/wstrinz/plane-jacobian-72-108}{\filename{github.com/wstrinz/plane-jacobian-72-108}} (2026).

\bibitem{Jelonek1993} Jelonek, Z. \textit{The set of points at which a polynomial map is not proper.} Ann.\ Polon.\ Math.\ 58 (1993), no.~3, 259--266. doi:10.4064/ap-58-3-259-266.

\bibitem{Chau2008} Nguyen Van Chau. \textit{Plane Jacobian problem for rational polynomials.} arXiv:0801.4138.

\bibitem{Chau2009} Nguyen Van Chau. \textit{Pencil of irreducible rational curves and plane Jacobian conjecture.} arXiv:0905.3939.

\bibitem{AH2015} Abhyankar, S.~S., Heinzer, W.~J. \textit{Existence of dicritical divisors revisited.} arXiv:1508.06015.

\bibitem{AbhyankarSurvey} \textit{Abhyankar's work on dicritical divisors.} arXiv:1707.06733. A survey; it states no equivalence between dicriticals and the conjecture.

\bibitem{BBR1962} Bia\l{}ynicki-Birula, A., Rosenlicht, M. \textit{Injective morphisms of real algebraic varieties.} Proc.\ Amer.\ Math.\ Soc.\ 13 (1962), 200--203.

\bibitem{Restatement} McCrary, B.~D. \textit{The two-dicritical theorem, restated.} Unpublished working note, \filename{branch\_two\_dicritical/RESTATEMENT.md}, October 13, 2026. Not published or deposited.

\end{thebibliography}

\end{document}
